% Options for packages loaded elsewhere
\PassOptionsToPackage{unicode}{hyperref}
\PassOptionsToPackage{hyphens}{url}
%
\documentclass[
  9pt,
  a4paper,
]{article}
\usepackage{amsmath,amssymb}
\usepackage{iftex}
\ifPDFTeX
  \usepackage[T1]{fontenc}
  \usepackage[utf8]{inputenc}
  \usepackage{textcomp} % provide euro and other symbols
\else % if luatex or xetex
  \usepackage{unicode-math} % this also loads fontspec
  \defaultfontfeatures{Scale=MatchLowercase}
  \defaultfontfeatures[\rmfamily]{Ligatures=TeX,Scale=1}
\fi
\usepackage{lmodern}
\ifPDFTeX\else
  % xetex/luatex font selection
  \setmainfont[]{DejaVu Sans}
  \setmonofont[]{DejaVu Sans Mono}
\fi
% Use upquote if available, for straight quotes in verbatim environments
\IfFileExists{upquote.sty}{\usepackage{upquote}}{}
\IfFileExists{microtype.sty}{% use microtype if available
  \usepackage[]{microtype}
  \UseMicrotypeSet[protrusion]{basicmath} % disable protrusion for tt fonts
}{}
\makeatletter
\@ifundefined{KOMAClassName}{% if non-KOMA class
  \IfFileExists{parskip.sty}{%
    \usepackage{parskip}
  }{% else
    \setlength{\parindent}{0pt}
    \setlength{\parskip}{6pt plus 2pt minus 1pt}}
}{% if KOMA class
  \KOMAoptions{parskip=half}}
\makeatother
\usepackage{xcolor}
\usepackage[margin=19mm]{geometry}
\setlength{\emergencystretch}{3em} % prevent overfull lines
\providecommand{\tightlist}{%
  \setlength{\itemsep}{0pt}\setlength{\parskip}{0pt}}
\setcounter{secnumdepth}{-\maxdimen} % remove section numbering
\ifLuaTeX
\usepackage[bidi=basic]{babel}
\else
\usepackage[bidi=default]{babel}
\fi
\babelprovide[main,import]{french}
\ifPDFTeX
\else
\babelfont{rm}[]{DejaVu Sans}
\fi
% get rid of language-specific shorthands (see #6817):
\let\LanguageShortHands\languageshorthands
\def\languageshorthands#1{}
\usepackage{etoolbox}
\usepackage{microtype}
\usepackage{float}
\floatplacement{figure}{H}
\sloppy
\setlength{\emergencystretch}{5em}
\AtBeginEnvironment{longtable}{\small}
\ifLuaTeX
  \usepackage{selnolig}  % disable illegal ligatures
\fi
\IfFileExists{bookmark.sty}{\usepackage{bookmark}}{\usepackage{hyperref}}
\IfFileExists{xurl.sty}{\usepackage{xurl}}{} % add URL line breaks if available
\urlstyle{same}
\hypersetup{
  pdftitle={Dictionnaire progressif des domaines},
  pdfauthor={Vincent P. P. NOËL - dossier de recherche, rédaction assistée},
  pdflang={fr-FR},
  hidelinks,
  pdfcreator={LaTeX via pandoc}}

\title{Dictionnaire progressif des domaines}
\usepackage{etoolbox}
\makeatletter
\providecommand{\subtitle}[1]{% add subtitle to \maketitle
  \apptocmd{\@title}{\par {\large #1 \par}}{}{}
}
\makeatother
\subtitle{La langue des passages - V6.44 - Document 69}
\author{Vincent P. P. NOËL - dossier de recherche, rédaction assistée}
\date{12 octobre 2026}

\begin{document}
\maketitle

{
\setcounter{tocdepth}{1}
\tableofcontents
}
\hypertarget{comment-apprendre-cette-langue}{%
\section{Comment apprendre cette
langue}\label{comment-apprendre-cette-langue}}

Une fiche commence par une image simple, puis donne un symbole, son
domaine et la condition qui rend le geste autorisé. Le nom imagé aide à
comprendre ; la définition et la garde portent le calcul. Les mots
usuels restent canoniques.

La langue des passages relie le \textbf{moule} P, la \textbf{valise} U,
la \textbf{portée} Fourier ou Mellin, le \textbf{puzzle} de changement
de base et la \textbf{balance} de Schur. Son arithmétique distingue
addition, produit ordinaire et convolution. Sa logique écrit ce qui est
conservé, ce qui se perd et le budget du passage. Les preuves L1--L8 et
les domaines exacts sont dans le document 68.

Une phrase possède quatre étiquettes : support S, ordre de jets r aux
±1/2, ordre q en 0, et unité d'amplitude A\^{}u. Le geste P gagne r ; la
dérivée gagne q ; la convolution additionne les ordres. Le produit
ordinaire peut perdre les moments. Les jets de test aux ±1/2, les jets
de test en 0 et les jets des zéros critiques sont trois entrées
différentes de la langue.

Le noyau ci-dessous explique les objets de l'expérience actuelle et les
principaux objets du corpus. Tous les 95 noms du registre canonique sont
repris avec leur ancrage et leur garde. L'index technique séparé donne
toutes les déclarations Python analysées : une fonction de programme
n'est pas automatiquement un opérateur mathématique. Cette édition ne
prétend pas avoir désambiguïsé chaque symbole de toutes les archives
historiques.

Une fiche marquée définition n'est pas une preuve nouvelle ; un statut
du registre est hérité, sans validation additionnelle. Les calculs V644
sont exploratoires, sauf l'algèbre rationnelle et le contrat analytique
entier de queue explicités.

Cette édition contient \textbf{219 fiches du noyau}, \textbf{95
métaphores nommées} et un index de \textbf{986 fonctions}, \textbf{15
classes}, dans \textbf{155 fichiers avec déclarations}, avec 0 erreur de
parseur.

\hypertarget{niveau-1---les-petits-mots-et-les-unituxe9s}{%
\section{Niveau 1 - Les petits mots et les
unités}\label{niveau-1---les-petits-mots-et-les-unituxe9s}}

\hypertarget{nombre}{%
\subsection{Nombre}\label{nombre}}

\textbf{Comme un premier jeu.} Une place sur une règle.

\textbf{Symbole.} x∈ℝ

\textbf{Sens et domaine.} Un réel sert de coordonnée ; les nombres
complexes ajoutent une seconde coordonnée.

\textbf{Garde.} Une lettre peut changer de rôle ; son domaine doit être
annoncé.

\textbf{Ancrage.} Documents 67-68 et noyau du corpus ; définition locale
explicitée dans cette fiche ; définition pédagogique ; voir portée
précise dans ancrage et garde.

\hypertarget{nombre-complexe}{%
\subsection{Nombre complexe}\label{nombre-complexe}}

\textbf{Comme un premier jeu.} Une flèche qui va aussi vers le côté.

\textbf{Symbole.} z=a+ib, i²=-1

\textbf{Sens et domaine.} ℂ est un plan avec addition et multiplication
; Re z=a et Im z=b.

\textbf{Garde.} Complexe ne signifie ni positif ni négatif ; seul un
réel peut recevoir cet ordre.

\textbf{Ancrage.} Documents 67-68 et noyau du corpus ; définition locale
explicitée dans cette fiche ; définition pédagogique ; voir portée
précise dans ancrage et garde.

\hypertarget{conjugaison}{%
\subsection{Conjugaison}\label{conjugaison}}

\textbf{Comme un premier jeu.} La flèche se regarde dans un miroir.

\textbf{Symbole.} z̄=a-ib

\textbf{Sens et domaine.} Involution de ℂ ; dans un produit scalaire
elle rend la norme réelle.

\textbf{Garde.} Oublier la conjugaison peut transformer une forme
hermitienne en autre chose.

\textbf{Ancrage.} Documents 67-68 et noyau du corpus ; définition locale
explicitée dans cette fiche ; définition pédagogique ; voir portée
précise dans ancrage et garde.

\hypertarget{module}{%
\subsection{Module}\label{module}}

\textbf{Comme un premier jeu.} La longueur de la flèche.

\textbf{Symbole.} \textbar z\textbar²=z̄z

\textbf{Sens et domaine.} Quantité réelle non négative ; homogène au
carré de l'amplitude.

\textbf{Garde.} La longueur oublie la phase, donc deux flèches
différentes peuvent avoir le même module.

\textbf{Ancrage.} Documents 67-68 et noyau du corpus ; définition locale
explicitée dans cette fiche ; définition pédagogique ; voir portée
précise dans ancrage et garde.

\hypertarget{fonction}{%
\subsection{Fonction}\label{fonction}}

\textbf{Comme un premier jeu.} Une machine prend un jouet et rend un
résultat.

\textbf{Symbole.} f:X→Y

\textbf{Sens et domaine.} À chaque entrée du domaine X correspond une
sortie de Y.

\textbf{Garde.} Composer g et f exige que la sortie de f entre dans le
domaine de g.

\textbf{Ancrage.} Documents 67-68 et noyau du corpus ; définition locale
explicitée dans cette fiche ; définition pédagogique ; voir portée
précise dans ancrage et garde.

\hypertarget{somme}{%
\subsection{Somme}\label{somme}}

\textbf{Comme un premier jeu.} On rassemble des petites piles.

\textbf{Symbole.} Σ\_k a\_k

\textbf{Sens et domaine.} Somme finie ou limite de sommes partielles
pour une série.

\textbf{Garde.} Une série infinie exige une règle de convergence ;
réordonner peut changer une série conditionnelle.

\textbf{Ancrage.} Documents 67-68 et noyau du corpus ; définition locale
explicitée dans cette fiche ; définition pédagogique ; voir portée
précise dans ancrage et garde.

\hypertarget{intuxe9grale}{%
\subsection{Intégrale}\label{intuxe9grale}}

\textbf{Comme un premier jeu.} On rassemble toutes les fines bandes d'un
dessin.

\textbf{Symbole.} ∫ f(x) dx

\textbf{Sens et domaine.} Intégrale de Lebesgue lorsque f est intégrable
; remplace une somme discrète.

\textbf{Garde.} Une intégrale impropre ne peut pas être manipulée sans
vérifier sa convergence.

\textbf{Ancrage.} Documents 67-68 et noyau du corpus ; définition locale
explicitée dans cette fiche ; définition pédagogique ; voir portée
précise dans ancrage et garde.

\hypertarget{duxe9rivuxe9e}{%
\subsection{Dérivée}\label{duxe9rivuxe9e}}

\textbf{Comme un premier jeu.} On regarde comment la pente change.

\textbf{Symbole.} Df=f′

\textbf{Sens et domaine.} D est une dérivation, classique ou faible
selon le domaine.

\textbf{Garde.} Df et f n'ont pas automatiquement les mêmes unités ; dx
doit être typé.

\textbf{Ancrage.} Documents 67-68 et noyau du corpus ; définition locale
explicitée dans cette fiche ; définition pédagogique ; voir portée
précise dans ancrage et garde.

\hypertarget{courbure}{%
\subsection{Courbure}\label{courbure}}

\textbf{Comme un premier jeu.} On regarde si le chemin se courbe.

\textbf{Symbole.} D²f=f″

\textbf{Sens et domaine.} Deux dérivations ; définie classiquement pour
C² ou faiblement dans H².

\textbf{Garde.} Un coin peut ajouter une distribution quand on prolonge
une fonction par zéro.

\textbf{Ancrage.} Documents 67-68 et noyau du corpus ; définition locale
explicitée dans cette fiche ; définition pédagogique ; voir portée
précise dans ancrage et garde.

\hypertarget{duxe9rivuxe9e-partielle}{%
\subsection{Dérivée partielle}\label{duxe9rivuxe9e-partielle}}

\textbf{Comme un premier jeu.} On tourne un bouton en gardant les autres
immobiles.

\textbf{Symbole.} ∂\_t F

\textbf{Sens et domaine.} Dérivation selon une variable déclarée.

\textbf{Garde.} Une dérivée en temps physique est une puissance si la
grandeur est une énergie ; un paramètre arithmétique n'est pas un temps
physique.

\textbf{Ancrage.} Documents 67-68 et noyau du corpus ; définition locale
explicitée dans cette fiche ; définition pédagogique ; voir portée
précise dans ancrage et garde.

\hypertarget{support}{%
\subsection{Support}\label{support}}

\textbf{Comme un premier jeu.} L'endroit où le dessin n'est pas effacé.

\textbf{Symbole.} supp f⊂{[}-S,S{]}

\textbf{Sens et domaine.} Fermeture des points où la fonction est non
nulle, au sens essentiel pour L².

\textbf{Garde.} Déplacer le dessin change son adresse, pas son diamètre.

\textbf{Ancrage.} Documents 67-68 et noyau du corpus ; définition locale
explicitée dans cette fiche ; définition pédagogique ; voir portée
précise dans ancrage et garde.

\hypertarget{limite}{%
\subsection{Limite}\label{limite}}

\textbf{Comme un premier jeu.} Les dessins se rapprochent de plus en
plus.

\textbf{Symbole.} f\_n→f

\textbf{Sens et domaine.} Le sens dépend de la norme ou de la topologie
annoncée.

\textbf{Garde.} Convergence ponctuelle, L² et H¹ ne sont pas
interchangeables.

\textbf{Ancrage.} Documents 67-68 et noyau du corpus ; définition locale
explicitée dans cette fiche ; définition pédagogique ; voir portée
précise dans ancrage et garde.

\hypertarget{norme}{%
\subsection{Norme}\label{norme}}

\textbf{Comme un premier jeu.} Une règle qui mesure un dessin.

\textbf{Symbole.} \textbar\textbar f\textbar\textbar{}

\textbf{Sens et domaine.} Positive, homogène, triangulaire ; la métrique
choisie fixe ce que signifie proche.

\textbf{Garde.} Changer la règle conserve parfois le signe, pas les
tailles ni les marges.

\textbf{Ancrage.} Documents 67-68 et noyau du corpus ; définition locale
explicitée dans cette fiche ; définition pédagogique ; voir portée
précise dans ancrage et garde.

\hypertarget{produit-scalaire}{%
\subsection{Produit scalaire}\label{produit-scalaire}}

\textbf{Comme un premier jeu.} On mesure combien deux dessins se
ressemblent.

\textbf{Symbole.} ⟨f,g⟩=∫f̄g

\textbf{Sens et domaine.} Forme sesquilinéaire hermitienne ; ici
convention antilinéaire en première variable.

\textbf{Garde.} Une métrique externe doit être définie sans présupposer
le signe de Q\_W.

\textbf{Ancrage.} Documents 67-68 et noyau du corpus ; définition locale
explicitée dans cette fiche ; définition pédagogique ; voir portée
précise dans ancrage et garde.

\hypertarget{composition}{%
\subsection{Composition}\label{composition}}

\textbf{Comme un premier jeu.} Deux machines attachées l'une à l'autre.

\textbf{Symbole.} g∘f

\textbf{Sens et domaine.} Appliquer f puis g ; domaines et codomaines
compatibles.

\textbf{Garde.} Le nom d'une machine ne suffit pas pour autoriser un
passage.

\textbf{Ancrage.} Documents 67-68 et noyau du corpus ; définition locale
explicitée dans cette fiche ; définition pédagogique ; voir portée
précise dans ancrage et garde.

\hypertarget{unituxe9s-fictives}{%
\subsection{Unités fictives}\label{unituxe9s-fictives}}

\textbf{Comme un premier jeu.} Chaque quantité porte une étiquette sur
sa boîte.

\textbf{Symbole.} {[}f{]}=A, {[}Q{]}=E=A²

\textbf{Sens et domaine.} x=log u est sans dimension ; A et E sont des
unités formelles de tenue des calculs.

\textbf{Garde.} Une unité cohérente ne prouve pas un signe et ne donne
pas une énergie physique.

\textbf{Ancrage.} Documents 67-68 et noyau du corpus ; définition locale
explicitée dans cette fiche ; définition pédagogique ; voir portée
précise dans ancrage et garde.

\hypertarget{puissance-par-duxe9rivation}{%
\subsection{Puissance par
dérivation}\label{puissance-par-duxe9rivation}}

\textbf{Comme un premier jeu.} On mesure combien la pile change à chaque
tic.

\textbf{Symbole.} dE/dt

\textbf{Sens et domaine.} Si t est un temps, unité E/T ; si t est un
indice arithmétique, unité formelle différente.

\textbf{Garde.} E≥0 n'impose pas dE/dt≥0 ; par exemple
E(t)=e\^{}(-t)\textgreater0 et E′\textless0.

\textbf{Ancrage.} Documents 67-68 et noyau du corpus ; définition locale
explicitée dans cette fiche ; définition pédagogique ; voir portée
précise dans ancrage et garde.

\hypertarget{uxe9nergie-arithmuxe9tique-formelle}{%
\subsection{Énergie arithmétique
formelle}\label{uxe9nergie-arithmuxe9tique-formelle}}

\textbf{Comme un premier jeu.} La balance porte l'étiquette énergie.

\textbf{Symbole.} E\_ar=Q\_W(f)

\textbf{Sens et domaine.} Unité fictive A² pour une forme quadratique ;
changement d'amplitude c multiplie par \textbar c\textbar².

\textbf{Garde.} Le nom énergie ne donne pas Q\_W≥0 : c'est précisément
le signe recherché.

\textbf{Ancrage.} Documents 67-68 et noyau du corpus ; définition locale
explicitée dans cette fiche ; définition pédagogique ; voir portée
précise dans ancrage et garde.

\hypertarget{positivituxe9}{%
\subsection{Positivité}\label{positivituxe9}}

\textbf{Comme un premier jeu.} La balance ne descend pas sous sa marque.

\textbf{Symbole.} x≥0 ; Q(f)≥0

\textbf{Sens et domaine.} Ordre sur les réels ; positivité d'une forme
pour toutes les entrées du domaine.

\textbf{Garde.} Produit de nombres positifs positif ; une somme signée
ou une dérivée n'hérite pas de cette règle.

\textbf{Ancrage.} Documents 67-68 et noyau du corpus ; définition locale
explicitée dans cette fiche ; définition pédagogique ; voir portée
précise dans ancrage et garde.

\hypertarget{ruxe8gle-du-produit}{%
\subsection{Règle du produit}\label{ruxe8gle-du-produit}}

\textbf{Comme un premier jeu.} Deux dessins bougent, il faut compter
deux effets.

\textbf{Symbole.} D(fg)=f′g+fg′

\textbf{Sens et domaine.} Identité de Leibniz sous régularité
suffisante.

\textbf{Garde.} Les cutoffs et les phases apportent donc des termes
supplémentaires.

\textbf{Ancrage.} Documents 67-68 et noyau du corpus ; définition locale
explicitée dans cette fiche ; définition pédagogique ; voir portée
précise dans ancrage et garde.

\hypertarget{cauchyschwarz}{%
\subsection{Cauchy--Schwarz}\label{cauchyschwarz}}

\textbf{Comme un premier jeu.} L'attache de deux dessins ne dépasse pas
leurs tailles.

\textbf{Symbole.}
\textbar⟨f,g⟩\textbar≤\textbar\textbar f\textbar\textbar\textbar\textbar g\textbar\textbar{}

\textbf{Sens et domaine.} Inégalité Hilbert, utilisée pour majorer
corrélations et primes.

\textbf{Garde.} La majoration de taille ne décide pas du signe de
l'attache.

\textbf{Ancrage.} Documents 67-68 et noyau du corpus ; définition locale
explicitée dans cette fiche ; définition pédagogique ; voir portée
précise dans ancrage et garde.

\hypertarget{identituxe9}{%
\subsection{Identité}\label{identituxe9}}

\textbf{Comme un premier jeu.} Le jouet revient pareil.

\textbf{Symbole.} If=f

\textbf{Sens et domaine.} Opérateur neutre dans le domaine choisi.

\textbf{Garde.} Dans I-K, la présence de I ne prouve pas que K≤I.

\textbf{Ancrage.} Documents 67-68 et noyau du corpus ; définition locale
explicitée dans cette fiche ; définition pédagogique ; voir portée
précise dans ancrage et garde.

\hypertarget{partie-entiuxe8re}{%
\subsection{Partie entière}\label{partie-entiuxe8re}}

\textbf{Comme un premier jeu.} On garde seulement les piles entières.

\textbf{Symbole.} ⌊x⌋

\textbf{Sens et domaine.} Plus grand entier ≤x ; détermine le plafond
discret des puissances premières.

\textbf{Garde.} Près d'un entier, il faut encadrer x avant de décider.

\textbf{Ancrage.} Documents 67-68 et noyau du corpus ; définition locale
explicitée dans cette fiche ; définition pédagogique ; voir portée
précise dans ancrage et garde.

\hypertarget{partie-ruxe9elle}{%
\subsection{Partie réelle}\label{partie-ruxe9elle}}

\textbf{Comme un premier jeu.} On garde le pas droit de la flèche.

\textbf{Symbole.} Re z=(z+z̄)/2

\textbf{Sens et domaine.} Projection ℂ→ℝ ; intervient dans Q\_pr et
polarisation.

\textbf{Garde.} Oublie la composante imaginaire ; ne permet pas de
reconstruire la phase seule.

\textbf{Ancrage.} Documents 67-68 et noyau du corpus ; définition locale
explicitée dans cette fiche ; définition pédagogique ; voir portée
précise dans ancrage et garde.

\hypertarget{linuxe9arituxe9}{%
\subsection{Linéarité}\label{linuxe9arituxe9}}

\textbf{Comme un premier jeu.} La machine respecte l'assemblage des
dessins.

\textbf{Symbole.} A(af+bg)=aAf+bAg

\textbf{Sens et domaine.} Propriété des opérateurs P, translations,
transforms et corrections fixes.

\textbf{Garde.} Q est quadratique, pas linéaire ; il contient des
attaches croisées.

\textbf{Ancrage.} Documents 67-68 et noyau du corpus ; définition locale
explicitée dans cette fiche ; définition pédagogique ; voir portée
précise dans ancrage et garde.

\hypertarget{forme-quadratique}{%
\subsection{Forme quadratique}\label{forme-quadratique}}

\textbf{Comme un premier jeu.} La pile grandit au carré de la taille.

\textbf{Symbole.} Q(cf)=\textbar c\textbar²Q(f)

\textbf{Sens et domaine.} Forme issue d'un couplage hermitien ; domaine
annoncé.

\textbf{Garde.} Homogénéité quadratique seule n'implique pas Q≥0.

\textbf{Ancrage.} Documents 67-68 et noyau du corpus ; définition locale
explicitée dans cette fiche ; définition pédagogique ; voir portée
précise dans ancrage et garde.

\hypertarget{niveau-2---les-fonctions-et-leurs-miroirs}{%
\section{Niveau 2 - Les fonctions et leurs
miroirs}\label{niveau-2---les-fonctions-et-leurs-miroirs}}

\hypertarget{exponentielle}{%
\subsection{Exponentielle}\label{exponentielle}}

\textbf{Comme un premier jeu.} Une taille qui se multiplie quand on
ajoute un pas.

\textbf{Symbole.} e\^{}x

\textbf{Sens et domaine.} Fonction réelle positive ; e\^{}(x+y)=e\^{}x
e\^{}y ; extension entière sur ℂ.

\textbf{Garde.} e\^{}z peut être complexe ; « exponentielle positive »
concerne x réel.

\textbf{Ancrage.} Documents 67-68 et noyau du corpus ; définition locale
explicitée dans cette fiche ; définition pédagogique ; voir portée
précise dans ancrage et garde.

\hypertarget{logarithme}{%
\subsection{Logarithme}\label{logarithme}}

\textbf{Comme un premier jeu.} On compte les pas multiplicatifs avec une
règle additive.

\textbf{Symbole.} x=log u

\textbf{Sens et domaine.} Inverse réel de exp sur u\textgreater0 ;
log(uv)=log u+log v.

\textbf{Garde.} Le logarithme complexe demande un choix de branche.

\textbf{Ancrage.} Documents 67-68 et noyau du corpus ; définition locale
explicitée dans cette fiche ; définition pédagogique ; voir portée
précise dans ancrage et garde.

\hypertarget{sinus-et-cosinus}{%
\subsection{Sinus et cosinus}\label{sinus-et-cosinus}}

\textbf{Comme un premier jeu.} Une roue qui tourne.

\textbf{Symbole.} e\^{}(it)=cos t+i sin t

\textbf{Sens et domaine.} Fonctions périodiques réelles ; décrivent les
phases Fourier.

\textbf{Garde.} Une phase modifie les couplages si elle dépend de x.

\textbf{Ancrage.} Documents 67-68 et noyau du corpus ; définition locale
explicitée dans cette fiche ; définition pédagogique ; voir portée
précise dans ancrage et garde.

\hypertarget{sinus-hyperbolique}{%
\subsection{Sinus hyperbolique}\label{sinus-hyperbolique}}

\textbf{Comme un premier jeu.} Deux tailles poussent en sens opposé.

\textbf{Symbole.} 2 sinh(x/2)=e\^{}(x/2)-e\^{}(-x/2)

\textbf{Sens et domaine.} Solution élémentaire utilisée pour l'inverse
causal de D²-1/4.

\textbf{Garde.} Ce noyau inverse exige les deux moments pour refermer le
support à droite.

\textbf{Ancrage.} Documents 67-68 et noyau du corpus ; définition locale
explicitée dans cette fiche ; définition pédagogique ; voir portée
précise dans ancrage et garde.

\hypertarget{fonction-gamma}{%
\subsection{Fonction Gamma}\label{fonction-gamma}}

\textbf{Comme un premier jeu.} Une machine prolonge la pile des
factorielles.

\textbf{Symbole.} Γ(s)=∫₀∞ t\^{}(s-1)e\^{}(-t)dt

\textbf{Sens et domaine.} Définie par l'intégrale pour Re s\textgreater0
puis prolongement méromorphe ; Γ(n+1)=n!.

\textbf{Garde.} Ses pôles et branches des facteurs associés doivent être
pris en compte.

\textbf{Ancrage.} Documents 67-68 et noyau du corpus ; définition locale
explicitée dans cette fiche ; définition pédagogique ; voir portée
précise dans ancrage et garde.

\hypertarget{fonction-zuxeata}{%
\subsection{Fonction zêta}\label{fonction-zuxeata}}

\textbf{Comme un premier jeu.} Une chorale où chaque entier chante.

\textbf{Symbole.} ζ(s)=Σₙ≥₁ n\^{}(-s)

\textbf{Sens et domaine.} Série absolument convergente pour Re
s\textgreater1 ; prolongement méromorphe, pôle simple en 1.

\textbf{Garde.} La série brute n'est pas la définition convergente dans
la bande critique.

\textbf{Ancrage.} Documents 67-68 et noyau du corpus ; définition locale
explicitée dans cette fiche ; définition pédagogique ; voir portée
précise dans ancrage et garde.

\hypertarget{zuxeata-compluxe9tuxe9e}{%
\subsection{Zêta complétée}\label{zuxeata-compluxe9tuxe9e}}

\textbf{Comme un premier jeu.} On habille la chorale pour voir son
miroir.

\textbf{Symbole.} ξ(s)=½s(s-1)π\^{}(-s/2)Γ(s/2)ζ(s)

\textbf{Sens et domaine.} Fonction entière ; ξ(s)=ξ(1-s), zéros
correspondant aux zéros non triviaux de ζ.

\textbf{Garde.} Symétrie ne suffit pas à imposer que tous les zéros
soient sur l'axe.

\textbf{Ancrage.} Documents 67-68 et noyau du corpus ; définition locale
explicitée dans cette fiche ; définition pédagogique ; voir portée
précise dans ancrage et garde.

\hypertarget{hypothuxe8se-de-riemann}{%
\subsection{Hypothèse de Riemann}\label{hypothuxe8se-de-riemann}}

\textbf{Comme un premier jeu.} Toutes les notes recherchées devraient
rester sur une ligne.

\textbf{Symbole.} ζ(ρ)=0 ⇒ Re ρ=1/2, 0\textless Re ρ\textless1

\textbf{Sens et domaine.} Énoncé global sur tous les zéros non triviaux.

\textbf{Garde.} Un nombre fini de notes, matrices ou certificats ne
couvre pas l'infini.

\textbf{Ancrage.} Documents 67-68 et noyau du corpus ; définition locale
explicitée dans cette fiche ; définition pédagogique ; voir portée
précise dans ancrage et garde.

\hypertarget{zuxe9ro-et-multiplicituxe9}{%
\subsection{Zéro et multiplicité}\label{zuxe9ro-et-multiplicituxe9}}

\textbf{Comme un premier jeu.} Le dessin touche zéro ; parfois il s'y
accroche deux fois.

\textbf{Symbole.} F(ρ)=0, F′(ρ)=0

\textbf{Sens et domaine.} Un zéro de multiplicité m annule les m
premiers jets et pas le suivant.

\textbf{Garde.} Zéro double d'une fonction n'est pas automatiquement
collision de branches suivies.

\textbf{Ancrage.} Documents 67-68 et noyau du corpus ; définition locale
explicitée dans cette fiche ; définition pédagogique ; voir portée
précise dans ancrage et garde.

\hypertarget{suxe9rie-de-dirichlet}{%
\subsection{Série de Dirichlet}\label{suxe9rie-de-dirichlet}}

\textbf{Comme un premier jeu.} Chaque entier reçoit un volume de voix.

\textbf{Symbole.} F(s)=Σ a\_n n\^{}(-s)

\textbf{Sens et domaine.} Fonction sur un demi-plan de convergence ;
coefficients a\_n déclarés.

\textbf{Garde.} Une équation fonctionnelle seule ne donne pas un produit
d'Euler.

\textbf{Ancrage.} Documents 67-68 et noyau du corpus ; définition locale
explicitée dans cette fiche ; définition pédagogique ; voir portée
précise dans ancrage et garde.

\hypertarget{produit-deuler}{%
\subsection{Produit d'Euler}\label{produit-deuler}}

\textbf{Comme un premier jeu.} Les briques premières fabriquent les
entiers.

\textbf{Symbole.} ζ(s)=∏ₚ(1-p\^{}(-s))\^{}(-1)

\textbf{Sens et domaine.} Identité absolument convergente pour Re
s\textgreater1.

\textbf{Garde.} La multiplicativité locale ne prouve pas RH.

\textbf{Ancrage.} Documents 67-68 et noyau du corpus ; définition locale
explicitée dans cette fiche ; définition pédagogique ; voir portée
précise dans ancrage et garde.

\hypertarget{poids-de-von-mangoldt}{%
\subsection{Poids de von Mangoldt}\label{poids-de-von-mangoldt}}

\textbf{Comme un premier jeu.} Seules les tours de briques premières
portent une étiquette.

\textbf{Symbole.} Λ(p\^{}m)=log p, Λ(n)=0 sinon

\textbf{Sens et domaine.} Fonction arithmétique ; -ζ′/ζ=ΣΛ(n)n\^{}(-s)
pour Re s\textgreater1.

\textbf{Garde.} Λ(n) ici n'est pas la constante de de Bruijn-Newman.

\textbf{Ancrage.} Documents 67-68 et noyau du corpus ; définition locale
explicitée dans cette fiche ; définition pédagogique ; voir portée
précise dans ancrage et garde.

\hypertarget{fonction-de-muxf6bius}{%
\subsection{Fonction de Möbius}\label{fonction-de-muxf6bius}}

\textbf{Comme un premier jeu.} On donne un signe aux sacs de briques
sans doublon.

\textbf{Symbole.} μ(n)=0 ou (-1)\^{}r

\textbf{Sens et domaine.} Nulle si un carré premier divise n ; sinon
signe selon le nombre de facteurs premiers.

\textbf{Garde.} Une covariance de μ ne révèle pas chaque zéro.

\textbf{Ancrage.} Documents 67-68 et noyau du corpus ; définition locale
explicitée dans cette fiche ; définition pédagogique ; voir portée
précise dans ancrage et garde.

\hypertarget{caractuxe8re-de-dirichlet}{%
\subsection{Caractère de Dirichlet}\label{caractuxe8re-de-dirichlet}}

\textbf{Comme un premier jeu.} Une couleur suit la règle de
multiplication.

\textbf{Symbole.} χ(mn)=χ(m)χ(n)

\textbf{Sens et domaine.} Caractère modulo q, périodique et nul sur les
non-unités ; L(s,χ)=Σχ(n)n\^{}(-s).

\textbf{Garde.} Ne pas identifier un mélange de fonctions L à un produit
d'Euler.

\textbf{Ancrage.} Documents 67-68 et noyau du corpus ; définition locale
explicitée dans cette fiche ; définition pédagogique ; voir portée
précise dans ancrage et garde.

\hypertarget{coefficients-de-li}{%
\subsection{Coefficients de Li}\label{coefficients-de-li}}

\textbf{Comme un premier jeu.} On rassemble les notes dans une suite de
balances.

\textbf{Symbole.} λ\_n=Σρ{[}1-(1-1/ρ)\^{}n{]}

\textbf{Sens et domaine.} Somme sur les zéros selon la convention
symétrique ; critère global de Li pour ξ.

\textbf{Garde.} Sommes tronquées et intervalle fini de n ne prouvent pas
toutes les inégalités.

\textbf{Ancrage.} Documents 67-68 et noyau du corpus ; définition locale
explicitée dans cette fiche ; définition pédagogique ; voir portée
précise dans ancrage et garde.

\hypertarget{flot-de-chaleur}{%
\subsection{Flot de chaleur}\label{flot-de-chaleur}}

\textbf{Comme un premier jeu.} On floute ou on accentue un dessin.

\textbf{Symbole.} ∂\_t H=-∂\_x²H selon convention

\textbf{Sens et domaine.} Flot de de Bruijn-Newman avec paramètre et
signe déclarés, lié à un seuil de réalité des zéros.

\textbf{Garde.} Le paramètre n'est pas une température physique et la
convention doit être fixée.

\textbf{Ancrage.} Documents 67-68 et noyau du corpus ; définition locale
explicitée dans cette fiche ; définition pédagogique ; voir portée
précise dans ancrage et garde.

\hypertarget{entropie-et-conventions}{%
\subsection{Entropie et conventions}\label{entropie-et-conventions}}

\textbf{Comme un premier jeu.} On compte combien un choix nous surprend.

\textbf{Symbole.} H(p)=-Σp\_i log p\_i

\textbf{Sens et domaine.} Entropie de Shannon non négative pour une loi
discrète ; entropie différentielle peut être négative.

\textbf{Garde.} Il n'existe pas une positivité universelle de toute
quantité appelée entropie.

\textbf{Ancrage.} Documents 67-68 et noyau du corpus ; définition locale
explicitée dans cette fiche ; définition pédagogique ; voir portée
précise dans ancrage et garde.

\hypertarget{constante-deuler}{%
\subsection{Constante d'Euler}\label{constante-deuler}}

\textbf{Comme un premier jeu.} Une petite correction reste après deux
grands comptes.

\textbf{Symbole.} γ\_E=lim(Σₖ≤n1/k-log n)

\textbf{Sens et domaine.} Constante réelle présente dans C₀.

\textbf{Garde.} Ne pas l'identifier au rabais γ de Bellman.

\textbf{Ancrage.} Documents 67-68 et noyau du corpus ; définition locale
explicitée dans cette fiche ; définition pédagogique ; voir portée
précise dans ancrage et garde.

\hypertarget{uxe9quation-fonctionnelle}{%
\subsection{Équation fonctionnelle}\label{uxe9quation-fonctionnelle}}

\textbf{Comme un premier jeu.} La chorale se répond dans un miroir.

\textbf{Symbole.} ξ(s)=ξ(1-s)

\textbf{Sens et domaine.} Symétrie analytique exacte de la fonction
complétée.

\textbf{Garde.} Le miroir permet des paires hors de la ligne ; il ne
force pas RH seul.

\textbf{Ancrage.} Documents 67-68 et noyau du corpus ; définition locale
explicitée dans cette fiche ; définition pédagogique ; voir portée
précise dans ancrage et garde.

\hypertarget{factorielle}{%
\subsection{Factorielle}\label{factorielle}}

\textbf{Comme un premier jeu.} On compte tous les rangements de n
jouets.

\textbf{Symbole.} n!=1·2·\ldots·n

\textbf{Sens et domaine.} Fonction sur les entiers non négatifs, 0!=1 ;
intervient dans Taylor.

\textbf{Garde.} Une série Taylor tronquée exige une borne de reste.

\textbf{Ancrage.} Documents 67-68 et noyau du corpus ; définition locale
explicitée dans cette fiche ; définition pédagogique ; voir portée
précise dans ancrage et garde.

\hypertarget{prolongement-analytique}{%
\subsection{Prolongement analytique}\label{prolongement-analytique}}

\textbf{Comme un premier jeu.} Le dessin continue quand ses règles le
forcent.

\textbf{Symbole.} F prolongée depuis un domaine

\textbf{Sens et domaine.} Unicité sur domaines connexes par identité
analytique ; méromorphe si pôles.

\textbf{Garde.} Cela ne donne pas automatiquement l'emplacement des
zéros.

\textbf{Ancrage.} Documents 67-68 et noyau du corpus ; définition locale
explicitée dans cette fiche ; définition pédagogique ; voir portée
précise dans ancrage et garde.

\hypertarget{puissance-premiuxe8re}{%
\subsection{Puissance première}\label{puissance-premiuxe8re}}

\textbf{Comme un premier jeu.} Une tour faite d'une seule brique
première.

\textbf{Symbole.} n=p\^{}m

\textbf{Sens et domaine.} Indices portant Λ(n) ; liste exacte finie à
support fixé.

\textbf{Garde.} Inclure toutes les puissances, pas seulement les
premiers simples.

\textbf{Ancrage.} Documents 67-68 et noyau du corpus ; définition locale
explicitée dans cette fiche ; définition pédagogique ; voir portée
précise dans ancrage et garde.

\hypertarget{arc-tangente-et-artanh}{%
\subsection{Arc tangente et artanh}\label{arc-tangente-et-artanh}}

\textbf{Comme un premier jeu.} Deux petites règles inversent des
courbes.

\textbf{Symbole.} atan z, artanh z

\textbf{Sens et domaine.} Pour les queues archimédiennes réelles ici
z=e\^{}(-S)∈(0,1), formules usuelles.

\textbf{Garde.} Les branches complexes ne sont pas utilisées dans ces
quadratures réelles.

\textbf{Ancrage.} Documents 67-68 et noyau du corpus ; définition locale
explicitée dans cette fiche ; définition pédagogique ; voir portée
précise dans ancrage et garde.

\hypertarget{branche-du-logarithme}{%
\subsection{Branche du logarithme}\label{branche-du-logarithme}}

\textbf{Comme un premier jeu.} On choisit où la rotation commence.

\textbf{Symbole.} Log F sur domaine sans zéro

\textbf{Sens et domaine.} Un logarithme holomorphe exige domaine
approprié et absence de zéro ; sinon branches locales.

\textbf{Garde.} La phase globale et les passages de branches ne sont pas
résolus par une formule locale.

\textbf{Ancrage.} Documents 67-68 et noyau du corpus ; définition locale
explicitée dans cette fiche ; définition pédagogique ; voir portée
précise dans ancrage et garde.

\hypertarget{puxe9riodicituxe9}{%
\subsection{Périodicité}\label{puxe9riodicituxe9}}

\textbf{Comme un premier jeu.} Le dessin recommence après le même pas.

\textbf{Symbole.} f(x+T)=f(x)

\textbf{Sens et domaine.} Relation exacte à période déclarée ; utilisée
dans des mécanismes de rigidité.

\textbf{Garde.} Le support compact, une fonction entière et un caractère
modulo q ont des domaines différents.

\textbf{Ancrage.} Documents 67-68 et noyau du corpus ; définition locale
explicitée dans cette fiche ; définition pédagogique ; voir portée
précise dans ancrage et garde.

\hypertarget{ligne-critique}{%
\subsection{Ligne critique}\label{ligne-critique}}

\textbf{Comme un premier jeu.} La porte finale ressemble à une ligne
dans un plan.

\textbf{Symbole.} Re s=1/2

\textbf{Sens et domaine.} Ensemble vertical où RH place tous les zéros
non triviaux.

\textbf{Garde.} La symétrie autour de cette ligne n'impose pas que les
zéros y restent.

\textbf{Ancrage.} Documents 67-68 et noyau du corpus ; définition locale
explicitée dans cette fiche ; définition pédagogique ; voir portée
précise dans ancrage et garde.

\hypertarget{rh-et-simplicituxe9}{%
\subsection{RH et simplicité}\label{rh-et-simplicituxe9}}

\textbf{Comme un premier jeu.} La ligne et les doubles notes sont deux
règles.

\textbf{Symbole.} RH ≠ tous les zéros simples

\textbf{Sens et domaine.} RH porte sur la partie réelle ; conjecture de
simplicité est plus forte et distincte.

\textbf{Garde.} Résoudre une collision de zéros n'est pas nécessairement
nécessaire à prouver RH par un critère de forme.

\textbf{Ancrage.} Documents 67-68 et noyau du corpus ; définition locale
explicitée dans cette fiche ; définition pédagogique ; voir portée
précise dans ancrage et garde.

\hypertarget{niveau-3---les-piuxe8ces-ouxf9-vivent-les-fonctions}{%
\section{Niveau 3 - Les pièces où vivent les
fonctions}\label{niveau-3---les-piuxe8ces-ouxf9-vivent-les-fonctions}}

\hypertarget{espace-luxb2}{%
\subsection{Espace L²}\label{espace-luxb2}}

\textbf{Comme un premier jeu.} Les dessins ont une taille totale finie.

\textbf{Symbole.} ∫\textbar f\textbar²\textless∞

\textbf{Sens et domaine.} Fonctions modulo égalité presque partout ;
Hilbert avec norme quadratique.

\textbf{Garde.} La dérivée et la forme de Weil ne sont pas contrôlées
par la seule norme L².

\textbf{Ancrage.} Documents 67-68 et noyau du corpus ; définition locale
explicitée dans cette fiche ; définition pédagogique ; voir portée
précise dans ancrage et garde.

\hypertarget{espace-huxb9}{%
\subsection{Espace H¹}\label{espace-huxb9}}

\textbf{Comme un premier jeu.} La règle mesure aussi les pentes.

\textbf{Symbole.}
\textbar\textbar f\textbar\textbar H¹²=\textbar\textbar f\textbar\textbar₂²+\textbar\textbar f′\textbar\textbar₂²

\textbf{Sens et domaine.} Fonctions L² avec dérivée faible L².

\textbf{Garde.} Sa positivité est externe ; un opérateur compact peut
avoir une forme positive sans marge uniforme.

\textbf{Ancrage.} Documents 67-68 et noyau du corpus ; définition locale
explicitée dans cette fiche ; définition pédagogique ; voir portée
précise dans ancrage et garde.

\hypertarget{espace-hux2080uxb3}{%
\subsection{Espace H₀³}\label{espace-hux2080uxb3}}

\textbf{Comme un premier jeu.} Le dessin et ses deux premières pentes
s'attachent au bord.

\textbf{Symbole.} ψ,ψ′,ψ″=0 aux bords

\textbf{Sens et domaine.} Fermeture de Cc∞(-S,S) dans H³ ; traces nulles
jusqu'à l'ordre 2.

\textbf{Garde.} La définition dépend du domaine ouvert et du
prolongement par zéro.

\textbf{Ancrage.} Documents 67-68 et noyau du corpus ; définition locale
explicitée dans cette fiche ; définition pédagogique ; voir portée
précise dans ancrage et garde.

\hypertarget{cux153ur-lisse}{%
\subsection{Cœur lisse}\label{cux153ur-lisse}}

\textbf{Comme un premier jeu.} Des dessins sans coins, rangés loin du
bord.

\textbf{Symbole.} Cc∞(-S,S)

\textbf{Sens et domaine.} Fonctions infiniment dérivables à support
compact intérieur.

\textbf{Garde.} Une fonction polynomiale coupée au bord n'est
généralement pas dans ce cœur.

\textbf{Ancrage.} Documents 67-68 et noyau du corpus ; définition locale
explicitée dans cette fiche ; définition pédagogique ; voir portée
précise dans ancrage et garde.

\hypertarget{translation}{%
\subsection{Translation}\label{translation}}

\textbf{Comme un premier jeu.} La valise change d'adresse sans changer
son contenu.

\textbf{Symbole.} U\_a f(x)=f(x-a)

\textbf{Sens et domaine.} Unitaire dans L² et H¹ sur ℝ ; mêmes
autocorrélations.

\textbf{Garde.} E\_S est transporté dans une fenêtre plus grande, pas
forcément dans la même E\_S.

\textbf{Ancrage.} Documents 67-68 et noyau du corpus ; définition locale
explicitée dans cette fiche ; définition pédagogique ; voir portée
précise dans ancrage et garde.

\hypertarget{dilatation-beurling}{%
\subsection{Dilatation Beurling}\label{dilatation-beurling}}

\textbf{Comme un premier jeu.} Le dessin tient dans une boîte plus
petite.

\textbf{Symbole.} D\_a f(t)=a\^{}(-1/2)f(t/a)1\_(0,a)(t)

\textbf{Sens et domaine.} Isométrie L²(0,1) pour 0\textless a≤1 ; sous J
elle devient translation.

\textbf{Garde.} Le facteur a\^{}(-1/2) est nécessaire et les domaines ne
sont pas celui de U sans J.

\textbf{Ancrage.} Documents 67-68 et noyau du corpus ; définition locale
explicitée dans cette fiche ; définition pédagogique ; voir portée
précise dans ancrage et garde.

\hypertarget{fraction-beurling}{%
\subsection{Fraction Beurling}\label{fraction-beurling}}

\textbf{Comme un premier jeu.} On garde le petit reste d'un partage.

\textbf{Symbole.} ρ\_θ(t)=\{θ/t\}, Σc\_jθ\_j=0

\textbf{Sens et domaine.} Combinaisons de parties fractionnaires dans
L²(0,1), condition admissible de Nyman-Beurling.

\textbf{Garde.} Leur densité est un critère de RH, donc ne peut pas être
supposée pour prouver RH.

\textbf{Ancrage.} Documents 67-68 et noyau du corpus ; définition locale
explicitée dans cette fiche ; définition pédagogique ; voir portée
précise dans ancrage et garde.

\hypertarget{passage-logarithmique-j}{%
\subsection{Passage logarithmique J}\label{passage-logarithmique-j}}

\textbf{Comme un premier jeu.} Une règle courbe devient une règle
droite.

\textbf{Symbole.} Jf(x)=e\^{}(-x/2)f(e\^{}(-x))

\textbf{Sens et domaine.} Unitaire L²(0,1)→L²(0,∞) ; change la mesure
avec son facteur.

\textbf{Garde.} L² ne devient pas H¹ par simple changement de
coordonnées.

\textbf{Ancrage.} Documents 67-68 et noyau du corpus ; définition locale
explicitée dans cette fiche ; définition pédagogique ; voir portée
précise dans ancrage et garde.

\hypertarget{passage-multiplicatif}{%
\subsection{Passage multiplicatif}\label{passage-multiplicatif}}

\textbf{Comme un premier jeu.} Des pas deviennent des multiplications.

\textbf{Symbole.} ℒf(u)=u\^{}(-1/2)f(log u)

\textbf{Sens et domaine.} Unitaire L²(ℝ,dx)→L²(ℝ₊,du).

\textbf{Garde.} Ne pas confondre ℒ avec J, qui utilise x=-log t et une
demi-droite.

\textbf{Ancrage.} Documents 67-68 et noyau du corpus ; définition locale
explicitée dans cette fiche ; définition pédagogique ; voir portée
précise dans ancrage et garde.

\hypertarget{transformuxe9e-de-fourier}{%
\subsection{Transformée de Fourier}\label{transformuxe9e-de-fourier}}

\textbf{Comme un premier jeu.} On entend les rythmes cachés dans le
dessin.

\textbf{Symbole.} f̂(τ)=(2π)\^{}(-1/2)∫f(x)e\^{}(-iτx)dx

\textbf{Sens et domaine.} Unitaire L² par Plancherel ; translations
deviennent phases, D devient iτ.

\textbf{Garde.} Une convention de normalisation différente change les
constantes.

\textbf{Ancrage.} Documents 67-68 et noyau du corpus ; définition locale
explicitée dans cette fiche ; définition pédagogique ; voir portée
précise dans ancrage et garde.

\hypertarget{transformuxe9e-de-mellin}{%
\subsection{Transformée de Mellin}\label{transformuxe9e-de-mellin}}

\textbf{Comme un premier jeu.} On entend les rythmes de changements de
taille.

\textbf{Symbole.} Mf(s)=∫₀∞ f(u)u\^{}(s-1)du

\textbf{Sens et domaine.} Pour ℒf, Mℒf(s)=∫f(x)e\^{}((s-1/2)x)dx ;
entier si support compact.

\textbf{Garde.} Les moments admissibles sont les valeurs en 0 et 1, pas
toute la transformée.

\textbf{Ancrage.} Documents 67-68 et noyau du corpus ; définition locale
explicitée dans cette fiche ; définition pédagogique ; voir portée
précise dans ancrage et garde.

\hypertarget{convolution}{%
\subsection{Convolution}\label{convolution}}

\textbf{Comme un premier jeu.} Un petit pinceau passe sur le dessin.

\textbf{Symbole.} (η*f)(x)=∫η(x-y)f(y)dy

\textbf{Sens et domaine.} Lissage et produit Fourier selon convention ;
support s'additionne.

\textbf{Garde.} Lisser peut déborder de la fenêtre ; il faut une marge
ou un cutoff.

\textbf{Ancrage.} Documents 67-68 et noyau du corpus ; définition locale
explicitée dans cette fiche ; définition pédagogique ; voir portée
précise dans ancrage et garde.

\hypertarget{mollificateur}{%
\subsection{Mollificateur}\label{mollificateur}}

\textbf{Comme un premier jeu.} Un pinceau de largeur δ.

\textbf{Symbole.} η\_δ(x)=δ\^{}(-1)η(x/δ), ∫η=1

\textbf{Sens et domaine.} Approximation de l'identité ; η lisse à
support compact.

\textbf{Garde.} Un contrôle en H¹ ou dans la métrique R doit être
prouvé, pas déduit de L² seul.

\textbf{Ancrage.} Documents 67-68 et noyau du corpus ; définition locale
explicitée dans cette fiche ; définition pédagogique ; voir portée
précise dans ancrage et garde.

\hypertarget{fonction-de-coupure}{%
\subsection{Fonction de coupure}\label{fonction-de-coupure}}

\textbf{Comme un premier jeu.} On efface doucement l'extérieur d'une
boîte.

\textbf{Symbole.} χ f

\textbf{Sens et domaine.} χ∈Cc∞, bornes de dérivées déclarées ;
multiplicateur local.

\textbf{Garde.} Les dérivées de χ ajoutent des termes au budget.

\textbf{Ancrage.} Documents 67-68 et noyau du corpus ; définition locale
explicitée dans cette fiche ; définition pédagogique ; voir portée
précise dans ancrage et garde.

\hypertarget{deux-moments}{%
\subsection{Deux moments}\label{deux-moments}}

\textbf{Comme un premier jeu.} Deux balances doivent être à zéro.

\textbf{Symbole.} ℓ±(f)=∫f(x)e\^{}(±x/2)dx

\textbf{Sens et domaine.} Fonctionnelles linéaires continues à support
fixé.

\textbf{Garde.} Elles ne sont ni l'intégrale seule ni une condition de
signe.

\textbf{Ancrage.} Documents 67-68 et noyau du corpus ; définition locale
explicitée dans cette fiche ; définition pédagogique ; voir portée
précise dans ancrage et garde.

\hypertarget{espace-admissible-e_s}{%
\subsection{Espace admissible E\_S}\label{espace-admissible-e_s}}

\textbf{Comme un premier jeu.} Les dessins passent deux portes et
restent dans la boîte.

\textbf{Symbole.} E\_S=H₀¹(-S,S)∩kerℓ₊∩kerℓ₋

\textbf{Sens et domaine.} Espace fermé de codimension 2 ; domaine H¹ de
la forme complète.

\textbf{Garde.} Changer la fenêtre ou la norme impose un nouveau
contrat.

\textbf{Ancrage.} Documents 67-68 et noyau du corpus ; définition locale
explicitée dans cette fiche ; définition pédagogique ; voir portée
précise dans ancrage et garde.

\hypertarget{moule-admissible-p}{%
\subsection{Moule admissible P}\label{moule-admissible-p}}

\textbf{Comme un premier jeu.} Un moule ferme les deux ouvertures.

\textbf{Symbole.} P=D²-1/4

\textbf{Sens et domaine.} Isomorphisme H₀³(-S,S)→E\_S ; annulations
exactes par intégration par parties.

\textbf{Garde.} P impose les moments, pas le signe de Q\_W. Ce P n'est
pas un polynôme zêta.

\textbf{Ancrage.} Document 68 ; V644\_LANGUE\_ET\_DOMAINES.md, L1-L7 ;
preuve auxiliaire interne ; portée et conditions dans L1-L7.

\hypertarget{inverse-du-moule}{%
\subsection{Inverse du moule}\label{inverse-du-moule}}

\textbf{Comme un premier jeu.} On retrouve le dessin avant le moule.

\textbf{Symbole.} ψ(x)=∫₋Sˣ 2sinh((x-y)/2)f(y)dy

\textbf{Sens et domaine.} Inverse de P pour f∈E\_S, deux moments ferment
les traces à droite.

\textbf{Garde.} Sans ces moments, le résultat ne reste pas dans H₀³.

\textbf{Ancrage.} Document 68 ; V644\_LANGUE\_ET\_DOMAINES.md, L1-L7 ;
preuve auxiliaire interne ; portée et conditions dans L1-L7.

\hypertarget{graine-polynomiale}{%
\subsection{Graine polynomiale}\label{graine-polynomiale}}

\textbf{Comme un premier jeu.} Une pâte qui colle bien aux bords.

\textbf{Symbole.} ψS,k=(1-x²/S²)\^{}4(x/S)\^{}k 1\_{[}-S,S{]}

\textbf{Sens et domaine.} Graine H₀⁴, C³ sur ℝ ; fS,k=PψS,k∈H₀² ;
famille indépendante et dense dans E\_S.

\textbf{Garde.} Pas C∞ en général ; la continuité H¹ de Q justifie son
usage.

\textbf{Ancrage.} Document 68 ; V644\_LANGUE\_ET\_DOMAINES.md, L1-L7 ;
preuve auxiliaire interne ; portée et conditions dans L1-L7.

\hypertarget{modulation}{%
\subsection{Modulation}\label{modulation}}

\textbf{Comme un premier jeu.} On ajoute une rotation à chaque place.

\textbf{Symbole.} M\_τ f=e\^{}(iτx)f

\textbf{Sens et domaine.} Unitaire L² ; déplace la fréquence Fourier.

\textbf{Garde.} Brise généralement les deux moments ; P M\_τ ψ reste
admissible mais change Q.

\textbf{Ancrage.} Documents 67-68 et noyau du corpus ; définition locale
explicitée dans cette fiche ; définition pédagogique ; voir portée
précise dans ancrage et garde.

\hypertarget{phase-constante}{%
\subsection{Phase constante}\label{phase-constante}}

\textbf{Comme un premier jeu.} Tout le dessin tourne du même angle.

\textbf{Symbole.} f↦e\^{}(iθ)f

\textbf{Sens et domaine.} Moments nuls, autocorrélation, R et Q\_W
conservés.

\textbf{Garde.} Une phase variable e\^{}(iτx) est une autre opération.

\textbf{Ancrage.} Documents 67-68 et noyau du corpus ; définition locale
explicitée dans cette fiche ; définition pédagogique ; voir portée
précise dans ancrage et garde.

\hypertarget{changement-duxe9chelle}{%
\subsection{Changement d'échelle}\label{changement-duxe9chelle}}

\textbf{Comme un premier jeu.} On agrandit la boîte et le dessin.

\textbf{Symbole.} f(x)↦f(x/S)

\textbf{Sens et domaine.} Le support et les dérivées changent ; P et les
points log(p\^{}m) ne se transportent pas inchangés.

\textbf{Garde.} Ce changement n'est pas une invariance de Q\_W.

\textbf{Ancrage.} Documents 67-68 et noyau du corpus ; définition locale
explicitée dans cette fiche ; définition pédagogique ; voir portée
précise dans ancrage et garde.

\hypertarget{duxe9rivuxe9e-faible}{%
\subsection{Dérivée faible}\label{duxe9rivuxe9e-faible}}

\textbf{Comme un premier jeu.} La pente est reconnue par de petits
dessins tests.

\textbf{Symbole.} ∫fφ′=-∫gφ

\textbf{Sens et domaine.} g est la dérivée distributionnelle de f
lorsqu'elle appartient à l'espace exigé.

\textbf{Garde.} La pente classique peut manquer aux coins ; les traces
restent importantes.

\textbf{Ancrage.} Documents 67-68 et noyau du corpus ; définition locale
explicitée dans cette fiche ; définition pédagogique ; voir portée
précise dans ancrage et garde.

\hypertarget{trace-au-bord}{%
\subsection{Trace au bord}\label{trace-au-bord}}

\textbf{Comme un premier jeu.} Le bord garde les attaches du dessin.

\textbf{Symbole.} f(±S), f′(±S),\ldots{}

\textbf{Sens et domaine.} Valeurs limites définies pour les Sobolev
suffisamment réguliers en dimension un.

\textbf{Garde.} Un prolongement par zéro sans traces nulles peut créer
des masses de Dirac.

\textbf{Ancrage.} Documents 67-68 et noyau du corpus ; définition locale
explicitée dans cette fiche ; définition pédagogique ; voir portée
précise dans ancrage et garde.

\hypertarget{espace-de-schwartz}{%
\subsection{Espace de Schwartz}\label{espace-de-schwartz}}

\textbf{Comme un premier jeu.} Le dessin et ses pentes s'effacent très
vite.

\textbf{Symbole.} sup \textbar x\^{}mD\^{}nf(x)\textbar\textless∞

\textbf{Sens et domaine.} Fonctions lisses à décroissance rapide, cœur
fréquentiel pour les approximations.

\textbf{Garde.} Décroissance rapide n'est pas support compact.

\textbf{Ancrage.} Documents 67-68 et noyau du corpus ; définition locale
explicitée dans cette fiche ; définition pédagogique ; voir portée
précise dans ancrage et garde.

\hypertarget{changement-de-mesure}{%
\subsection{Changement de mesure}\label{changement-de-mesure}}

\textbf{Comme un premier jeu.} La règle change l'épaisseur des bandes.

\textbf{Symbole.} dt=e\^{}(-x)dx

\textbf{Sens et domaine.} Le facteur racine dans J et ℒ conserve les
normes L².

\textbf{Garde.} Oublier la racine du jacobien détruit l'isométrie.

\textbf{Ancrage.} Documents 67-68 et noyau du corpus ; définition locale
explicitée dans cette fiche ; définition pédagogique ; voir portée
précise dans ancrage et garde.

\hypertarget{injection-j_s}{%
\subsection{Injection j\_S}\label{injection-j_s}}

\textbf{Comme un premier jeu.} On range un dessin dans une boîte plus
large.

\textbf{Symbole.} j\_S:X\_S→L²

\textbf{Sens et domaine.} Application identité sur les fonctions mais
changement de norme ; compacte à S fixé.

\textbf{Garde.} Même fonction ne signifie pas même topologie.

\textbf{Ancrage.} Documents 67-68 et noyau du corpus ; définition locale
explicitée dans cette fiche ; définition pédagogique ; voir portée
précise dans ancrage et garde.

\hypertarget{contraction-du-support}{%
\subsection{Contraction du support}\label{contraction-du-support}}

\textbf{Comme un premier jeu.} On rapproche le dessin du centre avant de
lisser.

\textbf{Symbole.} D\_λf(x)=λ\^{}(-1/2)f(x/λ)

\textbf{Sens et domaine.} Pour 0\textless λ\textless1, réduit le support
; bornée et continue à λ=1 dans R.

\textbf{Garde.} Ne préserve pas les moments ; on les corrige après le
lissage.

\textbf{Ancrage.} Documents 67-68 et noyau du corpus ; définition locale
explicitée dans cette fiche ; définition pédagogique ; voir portée
précise dans ancrage et garde.

\hypertarget{convergence-absolue}{%
\subsection{Convergence absolue}\label{convergence-absolue}}

\textbf{Comme un premier jeu.} Toutes les briques tiennent même sans
leurs signes.

\textbf{Symbole.} Σ\textbar a\_n\textbar\textless∞

\textbf{Sens et domaine.} Autorise réarrangement d'une série et facilite
les échanges.

\textbf{Garde.} Ne pas présumer cette propriété dans la bande critique
d'une série de Dirichlet.

\textbf{Ancrage.} Documents 67-68 et noyau du corpus ; définition locale
explicitée dans cette fiche ; définition pédagogique ; voir portée
précise dans ancrage et garde.

\hypertarget{intuxe9gration-par-parties}{%
\subsection{Intégration par parties}\label{intuxe9gration-par-parties}}

\textbf{Comme un premier jeu.} On déplace la pente d'un dessin vers
l'autre.

\textbf{Symbole.} ∫ψ″e\^{}(±x/2)=¼∫ψe\^{}(±x/2)

\textbf{Sens et domaine.} Identité sous traces ψ=ψ′=0, explique ℓ±Pψ=0.

\textbf{Garde.} Sans conditions de bord, des termes supplémentaires
restent.

\textbf{Ancrage.} Documents 67-68 et noyau du corpus ; définition locale
explicitée dans cette fiche ; définition pédagogique ; voir portée
précise dans ancrage et garde.

\hypertarget{distribution}{%
\subsection{Distribution}\label{distribution}}

\textbf{Comme un premier jeu.} Un dessin peut être reconnu par son effet
sur les tests.

\textbf{Symbole.} δ\_a et Df au sens faible

\textbf{Sens et domaine.} Forme linéaire continue sur Cc∞ ; dérivées
définies par dualité.

\textbf{Garde.} Une masse de bord change l'équation Pψ=f ; il faut
annuler les traces.

\textbf{Ancrage.} Documents 67-68 et noyau du corpus ; définition locale
explicitée dans cette fiche ; définition pédagogique ; voir portée
précise dans ancrage et garde.

\hypertarget{diamuxe8tre-du-support}{%
\subsection{Diamètre du support}\label{diamuxe8tre-du-support}}

\textbf{Comme un premier jeu.} On mesure la largeur de la valise.

\textbf{Symbole.} D=b-a

\textbf{Sens et domaine.} Pour supp f⊂{[}a,b{]}, h\_f(t)=0 si
\textbar t\textbar\textgreater D.

\textbf{Garde.} Le centre ne change pas le plafond e\^{}D ; une fenêtre
plus large peut le surestimer.

\textbf{Ancrage.} Documents 67-68 et noyau du corpus ; définition locale
explicitée dans cette fiche ; définition pédagogique ; voir portée
précise dans ancrage et garde.

\hypertarget{jet-au-bord}{%
\subsection{Jet au bord}\label{jet-au-bord}}

\textbf{Comme un premier jeu.} Les attaches du bord disparaissent
jusqu'à quatre niveaux.

\textbf{Symbole.} ψ\^{}(j)(±S)=0, j≤3

\textbf{Sens et domaine.} Annulation d'ordre 4 de w\_S ; assure
prolongement C³ et dérivées faibles L² d'ordre 4.

\textbf{Garde.} La dérivée d'ordre 4 peut sauter ; ne pas annoncer C⁴ ou
C∞.

\textbf{Ancrage.} Documents 67-68 et noyau du corpus ; définition locale
explicitée dans cette fiche ; définition pédagogique ; voir portée
précise dans ancrage et garde.

\hypertarget{ruxe9flexion}{%
\subsection{Réflexion}\label{ruxe9flexion}}

\textbf{Comme un premier jeu.} Le dessin passe devant le miroir.

\textbf{Symbole.} Rf(x)=f(-x)

\textbf{Sens et domaine.} Unitaire sur métriques ; Q invariant car
corrélation symétrique.

\textbf{Garde.} Ce R de réflexion est distinct de la métrique R(f).

\textbf{Ancrage.} Documents 67-68 et noyau du corpus ; définition locale
explicitée dans cette fiche ; définition pédagogique ; voir portée
précise dans ancrage et garde.

\hypertarget{alguxe8bre-de-convolution}{%
\subsection{Algèbre de convolution}\label{alguxe8bre-de-convolution}}

\textbf{Comme un premier jeu.} Deux dessins s'assemblent en glissant.

\textbf{Symbole.} A=Cc∞(ℝ), opération *

\textbf{Sens et domaine.} Algèbre commutative associative, sans unité
parmi les fonctions ; ℒ\_exp(f*g)=ℒ\_exp(f)ℒ\_exp(g).

\textbf{Garde.} Le produit ordinaire est une autre multiplication.

\textbf{Ancrage.}
3\_Preuves\_et\_calculs/preuves/V644\_LANGUE\_ET\_DOMAINES.md ; L6-L8 ;
relation démontrée dans V644, L6-L8.

\hypertarget{portuxe9e-exponentielle}{%
\subsection{Portée exponentielle}\label{portuxe9e-exponentielle}}

\textbf{Comme un premier jeu.} Le dessin devient une partition de notes.

\textbf{Symbole.} ℒ\_exp f(z)=∫e\^{}(zx)f(x)dx

\textbf{Sens et domaine.} Fonction entière pour f compact ; variable z
distincte de s Mellin, z=s-1/2.

\textbf{Garde.} Les zéros en ±1/2 sont contraintes auxiliaires, pas
zéros de ζ.

\textbf{Ancrage.}
3\_Preuves\_et\_calculs/preuves/V644\_LANGUE\_ET\_DOMAINES.md ; L6-L8 ;
relation démontrée dans V644, L6-L8.

\hypertarget{produit-qui-perd-les-moments}{%
\subsection{Produit qui perd les
moments}\label{produit-qui-perd-les-moments}}

\textbf{Comme un premier jeu.} Deux dessins autorisés peuvent fabriquer
un dessin interdit.

\textbf{Symbole.} ℓ±(f²)\textgreater0 pour f réel non nul

\textbf{Sens et domaine.} Si f∈I₁ réel non nul, f²≥0 et poids
exponentiels positifs : produit hors I₁.

\textbf{Garde.} La convolution f*f reste admissible et renforce les
annulations.

\textbf{Ancrage.}
3\_Preuves\_et\_calculs/preuves/V644\_LANGUE\_ET\_DOMAINES.md ; L6-L8 ;
relation démontrée dans V644, L6-L8.

\hypertarget{niveau-4---les-formes-opuxe9rations-et-invariants}{%
\section{Niveau 4 - Les formes, opérations et
invariants}\label{niveau-4---les-formes-opuxe9rations-et-invariants}}

\hypertarget{adjoint}{%
\subsection{Adjoint}\label{adjoint}}

\textbf{Comme un premier jeu.} La machine se regarde dans la bonne
balance.

\textbf{Symbole.} ⟨Af,g⟩=⟨f,A*g⟩

\textbf{Sens et domaine.} Opérateur adjoint relativement aux domaines et
à la métrique déclarés.

\textbf{Garde.} A* dépend du produit scalaire ; les opérateurs non
bornés exigent leur domaine.

\textbf{Ancrage.} Documents 67-68 et noyau du corpus ; définition locale
explicitée dans cette fiche ; définition pédagogique ; voir portée
précise dans ancrage et garde.

\hypertarget{corruxe9lation}{%
\subsection{Corrélation}\label{corruxe9lation}}

\textbf{Comme un premier jeu.} On glisse deux dessins pour voir où ils
se recouvrent.

\textbf{Symbole.} h\_ij(t)=∫f\_i(y)f̄\_j(y-t)dy

\textbf{Sens et domaine.} Corrélation croisée ; h\_ij(-t)=conj h\_ji(t).

\textbf{Garde.} Une corrélation croisée pair/impair peut être non nulle.

\textbf{Ancrage.} Documents 67-68 et noyau du corpus ; définition locale
explicitée dans cette fiche ; définition pédagogique ; voir portée
précise dans ancrage et garde.

\hypertarget{corruxe9lation-symuxe9trisuxe9e}{%
\subsection{Corrélation
symétrisée}\label{corruxe9lation-symuxe9trisuxe9e}}

\textbf{Comme un premier jeu.} Deux glissements opposés partagent la
balance.

\textbf{Symbole.} H\_ij=(h\_ij+h\_ji)/2 pour tests réels

\textbf{Sens et domaine.} La partie réelle symétrique gouverne les
matrices réelles de Q\_W.

\textbf{Garde.} Seule cette symétrisation s'annule entre parités
opposées.

\textbf{Ancrage.} Documents 67-68 et noyau du corpus ; définition locale
explicitée dans cette fiche ; définition pédagogique ; voir portée
précise dans ancrage et garde.

\hypertarget{polarisation}{%
\subsection{Polarisation}\label{polarisation}}

\textbf{Comme un premier jeu.} On ouvre une balance pour voir les
attaches entre jouets.

\textbf{Symbole.} q(f,g) reconstruit depuis Q(f)

\textbf{Sens et domaine.} Identité de polarisation d'une forme
quadratique hermitienne.

\textbf{Garde.} Le couplage est indispensable ; des diagonales positives
ne suffisent pas.

\textbf{Ancrage.} Documents 67-68 et noyau du corpus ; définition locale
explicitée dans cette fiche ; définition pédagogique ; voir portée
précise dans ancrage et garde.

\hypertarget{forme-compluxe8te-de-weil}{%
\subsection{Forme complète de Weil}\label{forme-compluxe8te-de-weil}}

\textbf{Comme un premier jeu.} Une balance compte tous ses plateaux.

\textbf{Symbole.}
Q\_W=E\_arch-C₀\textbar\textbar f\textbar\textbar₂²-Q\_pr

\textbf{Sens et domaine.}
E\_arch=∫₀∞w(x)\textbar\textbar U\_xf-f\textbar\textbar₂²dx,
C₀=log(8π)+γ\_E+π/2 ; Q\_pr=2Σ Λ(n)n\^{}(-1/2)Re h\_f(log n).

\textbf{Garde.} Toujours utiliser la partie archimédienne ET les
premiers ; le signe global reste ouvert.

\textbf{Ancrage.} Documents 67-68 et noyau du corpus ; définition locale
explicitée dans cette fiche ; définition pédagogique ; voir portée
précise dans ancrage et garde.

\hypertarget{noyau-archimuxe9dien}{%
\subsection{Noyau archimédien}\label{noyau-archimuxe9dien}}

\textbf{Comme un premier jeu.} Les petits déplacements pèsent plus
lourd.

\textbf{Symbole.} w(x)=e\^{}(-x/2)/(1-e\^{}(-2x))

\textbf{Sens et domaine.} Positif pour x\textgreater0 ;
\textasciitilde1/(2x) près de 0, décroît à l'infini.

\textbf{Garde.} Sa singularité exige un contrôle de U\_x f-f.

\textbf{Ancrage.} Documents 67-68 et noyau du corpus ; définition locale
explicitée dans cette fiche ; définition pédagogique ; voir portée
précise dans ancrage et garde.

\hypertarget{forme-des-premiers}{%
\subsection{Forme des premiers}\label{forme-des-premiers}}

\textbf{Comme un premier jeu.} Les tours premières mettent des poids sur
la balance.

\textbf{Symbole.} Q\_pr=2Σₚ,m log p·p\^{}(-m/2)Re h\_f(m log p)

\textbf{Sens et domaine.} Somme finie à support de diamètre D :
p\^{}m≤e\^{}D.

\textbf{Garde.} Chaque corrélation peut changer de signe ; cette partie
n'est pas positive en général.

\textbf{Ancrage.} Documents 67-68 et noyau du corpus ; définition locale
explicitée dans cette fiche ; définition pédagogique ; voir portée
précise dans ancrage et garde.

\hypertarget{muxe9trique-logarithmique-r}{%
\subsection{Métrique logarithmique
R}\label{muxe9trique-logarithmique-r}}

\textbf{Comme un premier jeu.} On mesure la taille et les petits
glissements.

\textbf{Symbole.} R(f)=\textbar\textbar f\textbar\textbar₂²+E\_arch(f)

\textbf{Sens et domaine.} Forme positive fermée de multiplicateur
1+a(τ), indépendante des premiers et de RH.

\textbf{Garde.} Domaine X\_S plus large que H¹ ; le Riesz change quand
on change de métrique.

\textbf{Ancrage.} Document 68 ; V644\_LANGUE\_ET\_DOMAINES.md, L1-L7 ;
preuve auxiliaire interne ; portée et conditions dans L1-L7.

\hypertarget{poids-logarithmique-a}{%
\subsection{Poids logarithmique a}\label{poids-logarithmique-a}}

\textbf{Comme un premier jeu.} Les rythmes rapides finissent par peser
lourd.

\textbf{Symbole.} a(τ)=2Σₙ≥₀ τ²/{[}β\_n(β\_n²+τ²){]}, β\_n=2n+1/2

\textbf{Sens et domaine.} Positif, monotone en \textbar τ\textbar, tend
vers l'infini ; multiplicateur de E\_arch.

\textbf{Garde.} Il contrôle une queue fréquentielle, pas le signe
arithmétique à lui seul.

\textbf{Ancrage.} Document 68 ; V644\_LANGUE\_ET\_DOMAINES.md, L1-L7 ;
preuve auxiliaire interne ; portée et conditions dans L1-L7.

\hypertarget{domaine-logarithmique-x_s}{%
\subsection{Domaine logarithmique
X\_S}\label{domaine-logarithmique-x_s}}

\textbf{Comme un premier jeu.} Une boîte dont la règle surveille les
glissements.

\textbf{Symbole.} supp f⊂{[}-S,S{]}, R(f)\textless∞, ℓ±f=0

\textbf{Sens et domaine.} Hilbert avec métrique R ; cœur admissible Cc∞
dense ; injection compacte dans L².

\textbf{Garde.} La compactité vaut à S fixé ; ses constantes ne sont pas
uniformes en S.

\textbf{Ancrage.} Document 68 ; V644\_LANGUE\_ET\_DOMAINES.md, L1-L7 ;
preuve auxiliaire interne ; portée et conditions dans L1-L7.

\hypertarget{riesz-dans-huxb9}{%
\subsection{Riesz dans H¹}\label{riesz-dans-huxb9}}

\textbf{Comme un premier jeu.} La balance devient une machine sur les
dessins.

\textbf{Symbole.} Q\_W(f,g)=⟨T\_Sf,g⟩H¹

\textbf{Sens et domaine.} T\_S signé, borné, autoadjoint et compact sur
E\_S.

\textbf{Garde.} Aucune marge c\textgreater0 dans H¹ sur tout E\_S, même
si la forme est positive.

\textbf{Ancrage.} Document 68 ; V644\_LANGUE\_ET\_DOMAINES.md, L1-L7 ;
preuve auxiliaire interne ; portée et conditions dans L1-L7.

\hypertarget{riesz-dans-r}{%
\subsection{Riesz dans R}\label{riesz-dans-r}}

\textbf{Comme un premier jeu.} Le même dessin est l'identité moins une
correction.

\textbf{Symbole.} Q\_W(f,g)=⟨(I-K\_S)f,g⟩R

\textbf{Sens et domaine.} K\_S compact autoadjoint sur X\_S ; finitement
de directions négatives à S fixé.

\textbf{Garde.} Cela ne dit pas que leur nombre est zéro.

\textbf{Ancrage.} Document 68 ; V644\_LANGUE\_ET\_DOMAINES.md, L1-L7 ;
preuve auxiliaire interne ; portée et conditions dans L1-L7.

\hypertarget{opuxe9rateur-compact}{%
\subsection{Opérateur compact}\label{opuxe9rateur-compact}}

\textbf{Comme un premier jeu.} Beaucoup de dessins tiennent dans
quelques petites boîtes.

\textbf{Symbole.} A(ball) relativement compact

\textbf{Sens et domaine.} Les images de suites bornées ont des
sous-suites convergentes.

\textbf{Garde.} Compact ne signifie pas positif ; ses valeurs propres
peuvent s'approcher de zéro.

\textbf{Ancrage.} Document 68 ; V644\_LANGUE\_ET\_DOMAINES.md, L1-L7 ;
preuve auxiliaire interne ; portée et conditions dans L1-L7.

\hypertarget{coercivituxe9}{%
\subsection{Coercivité}\label{coercivituxe9}}

\textbf{Comme un premier jeu.} La balance garde toujours une marge
minimale.

\textbf{Symbole.} Q(f)≥c\textbar\textbar f\textbar\textbar²,
c\textgreater0

\textbf{Sens et domaine.} Minoration uniforme dans une métrique donnée.

\textbf{Garde.} La positivité Q≥0 est plus faible ; la marge H¹ uniforme
est exclue ici.

\textbf{Ancrage.} Documents 67-68 et noyau du corpus ; définition locale
explicitée dans cette fiche ; définition pédagogique ; voir portée
précise dans ancrage et garde.

\hypertarget{matrice-de-gram}{%
\subsection{Matrice de Gram}\label{matrice-de-gram}}

\textbf{Comme un premier jeu.} Une table dit comment les dessins se
chevauchent.

\textbf{Symbole.} B\_ij=⟨f\_i,f\_j⟩H¹

\textbf{Sens et domaine.} B positive définie pour une famille
indépendante ; ici rationnelle exacte.

\textbf{Garde.} B ne doit pas être identifiée à G\_W, la forme dont le
signe est recherché.

\textbf{Ancrage.} Documents 67-68 et noyau du corpus ; définition locale
explicitée dans cette fiche ; définition pédagogique ; voir portée
précise dans ancrage et garde.

\hypertarget{compression-de-weil}{%
\subsection{Compression de Weil}\label{compression-de-weil}}

\textbf{Comme un premier jeu.} On prend une petite fenêtre dans une
grande balance.

\textbf{Symbole.} (G\_W)\_ij=Q\_W(f\_i,f\_j)

\textbf{Sens et domaine.} Matrice hermitienne finie sur un sous-espace
admissible.

\textbf{Garde.} Un résultat fini ne vaut que sur ce sous-espace.

\textbf{Ancrage.} Documents 67-68 et noyau du corpus ; définition locale
explicitée dans cette fiche ; définition pédagogique ; voir portée
précise dans ancrage et garde.

\hypertarget{paire-de-matrices}{%
\subsection{Paire de matrices}\label{paire-de-matrices}}

\textbf{Comme un premier jeu.} Deux balances suivent les mêmes
coordonnées.

\textbf{Symbole.} G\_Wv=λBv

\textbf{Sens et domaine.} Problème spectral généralisé avec
B\textgreater0.

\textbf{Garde.} Transformer G\_W exige de transformer B simultanément.

\textbf{Ancrage.} Documents 67-68 et noyau du corpus ; définition locale
explicitée dans cette fiche ; définition pédagogique ; voir portée
précise dans ancrage et garde.

\hypertarget{changement-de-base}{%
\subsection{Changement de base}\label{changement-de-base}}

\textbf{Comme un premier jeu.} On renomme les pièces du puzzle.

\textbf{Symbole.} G↦C*GC, B↦C*BC

\textbf{Sens et domaine.} C inversible ; inertie de G et spectre
généralisé du couple conservés.

\textbf{Garde.} Entrées, valeurs propres ordinaires et pivots ne sont
pas tous conservés.

\textbf{Ancrage.} Documents 67-68 et noyau du corpus ; définition locale
explicitée dans cette fiche ; définition pédagogique ; voir portée
précise dans ancrage et garde.

\hypertarget{cholesky}{%
\subsection{Cholesky}\label{cholesky}}

\textbf{Comme un premier jeu.} On remet la règle de mesure à plat.

\textbf{Symbole.} B=LL*

\textbf{Sens et domaine.} Factorisation pour B hermitienne positive
définie ; normalise la paire.

\textbf{Garde.} Une base très mal conditionnée demande une précision
contrôlée.

\textbf{Ancrage.} Documents 67-68 et noyau du corpus ; définition locale
explicitée dans cette fiche ; définition pédagogique ; voir portée
précise dans ancrage et garde.

\hypertarget{valeur-propre}{%
\subsection{Valeur propre}\label{valeur-propre}}

\textbf{Comme un premier jeu.} Un dessin ressort dans sa propre
direction.

\textbf{Symbole.} Av=λv

\textbf{Sens et domaine.} λ indique la dilatation sur une direction ;
pour hermitien λ réel.

\textbf{Garde.} Une valeur calculée sans intervalle d'erreur reste
numérique.

\textbf{Ancrage.} Documents 67-68 et noyau du corpus ; définition locale
explicitée dans cette fiche ; définition pédagogique ; voir portée
précise dans ancrage et garde.

\hypertarget{inertie}{%
\subsection{Inertie}\label{inertie}}

\textbf{Comme un premier jeu.} On compte les directions qui montent,
descendent ou restent à zéro.

\textbf{Symbole.} (n₊,n₋,n₀)

\textbf{Sens et domaine.} Signature d'une forme hermitienne ; invariante
sous congruence inversible.

\textbf{Garde.} Le nombre de directions négatives fini n'est pas leur
absence.

\textbf{Ancrage.} Documents 67-68 et noyau du corpus ; définition locale
explicitée dans cette fiche ; définition pédagogique ; voir portée
précise dans ancrage et garde.

\hypertarget{parituxe9}{%
\subsection{Parité}\label{parituxe9}}

\textbf{Comme un premier jeu.} Deux dessins : miroir pareil ou miroir
opposé.

\textbf{Symbole.} f(-x)=±f(x)

\textbf{Sens et domaine.} Q\_W et les métriques respectent la réflexion
; blocs pair/impair découplés.

\textbf{Garde.} Un changement de base mêlant ces secteurs peut créer des
entrées croisées.

\textbf{Ancrage.} Document 68 ; V644\_LANGUE\_ET\_DOMAINES.md, L1-L7 ;
preuve auxiliaire interne ; portée et conditions dans L1-L7.

\hypertarget{compluxe9ment-de-schur}{%
\subsection{Complément de Schur}\label{compluxe9ment-de-schur}}

\textbf{Comme un premier jeu.} Le nouveau jouet tire sur les attaches
des anciens.

\textbf{Symbole.} δ=a-c*G\^{}(-1)c

\textbf{Sens et domaine.} Si G\textgreater0, δ est le minimum résiduel ;
sinon identité algébrique sous inversibilité.

\textbf{Garde.} Si G est singulière positive, il faut c∈Ran G et G†.

\textbf{Ancrage.} Documents 67-68 et noyau du corpus ; définition locale
explicitée dans cette fiche ; définition pédagogique ; voir portée
précise dans ancrage et garde.

\hypertarget{pseudo-inverse}{%
\subsection{Pseudo-inverse}\label{pseudo-inverse}}

\textbf{Comme un premier jeu.} La machine inverse ce qu'elle sait voir.

\textbf{Symbole.} G†

\textbf{Sens et domaine.} Inverse sur les valeurs propres non nulles,
zéro sur le noyau.

\textbf{Garde.} Ne répare pas un couplage hors de l'image.

\textbf{Ancrage.} Documents 67-68 et noyau du corpus ; définition locale
explicitée dans cette fiche ; définition pédagogique ; voir portée
précise dans ancrage et garde.

\hypertarget{drapeau-de-sous-espaces}{%
\subsection{Drapeau de sous-espaces}\label{drapeau-de-sous-espaces}}

\textbf{Comme un premier jeu.} On ouvre des boîtes dans un ordre choisi.

\textbf{Symbole.} V₁⊂V₂⊂\ldots{}

\textbf{Sens et domaine.} Une chaîne fixe sert à définir les dettes
successives.

\textbf{Garde.} Changer l'ordre change les pivots ; l'invariant de dette
exige le même drapeau.

\textbf{Ancrage.} Documents 67-68 et noyau du corpus ; définition locale
explicitée dans cette fiche ; définition pédagogique ; voir portée
précise dans ancrage et garde.

\hypertarget{dette-normalisuxe9e}{%
\subsection{Dette normalisée}\label{dette-normalisuxe9e}}

\textbf{Comme un premier jeu.} La nouvelle charge est mesurée avec la
même règle.

\textbf{Symbole.} δ\_G/δ\_B

\textbf{Sens et domaine.} Invariant sous changement de base ancien et
v↦αv+w ancien, pour drapeau fixé et inverses valides.

\textbf{Garde.} Ce quotient n'est pas invariant sous changement
arbitraire du drapeau.

\textbf{Ancrage.} Document 68 ; V644\_LANGUE\_ET\_DOMAINES.md, L1-L7 ;
preuve auxiliaire interne ; portée et conditions dans L1-L7.

\hypertarget{image-et-noyau}{%
\subsection{Image et noyau}\label{image-et-noyau}}

\textbf{Comme un premier jeu.} Ce que la machine peut produire, et ce
qu'elle efface.

\textbf{Symbole.} Ran G, ker G

\textbf{Sens et domaine.} Sous-espaces ; pour G hermitienne finie, Ran
G=(ker G)⊥.

\textbf{Garde.} Dans Schur singulier, c hors de Ran G crée une direction
négative.

\textbf{Ancrage.} Documents 67-68 et noyau du corpus ; définition locale
explicitée dans cette fiche ; définition pédagogique ; voir portée
précise dans ancrage et garde.

\hypertarget{duxe9terminant}{%
\subsection{Déterminant}\label{duxe9terminant}}

\textbf{Comme un premier jeu.} On mesure si un puzzle aplati perd une
direction.

\textbf{Symbole.} det G

\textbf{Sens et domaine.} Produit des valeurs propres en dimension finie
; zéro si matrice singulière.

\textbf{Garde.} Un déterminant positif seul ne garantit pas une forme
positive.

\textbf{Ancrage.} Documents 67-68 et noyau du corpus ; définition locale
explicitée dans cette fiche ; définition pédagogique ; voir portée
précise dans ancrage et garde.

\hypertarget{duxe9composition-ldl}{%
\subsection{Décomposition LDL*}\label{duxe9composition-ldl}}

\textbf{Comme un premier jeu.} On ouvre la balance pivot par pivot.

\textbf{Symbole.} G=LDL*

\textbf{Sens et domaine.} Factorisation avec pivotage adapté ; D porte
l'inertie.

\textbf{Garde.} Les petits pivots et les permutations doivent être
contrôlés ; ils dépendent du drapeau.

\textbf{Ancrage.} Documents 67-68 et noyau du corpus ; définition locale
explicitée dans cette fiche ; définition pédagogique ; voir portée
précise dans ancrage et garde.

\hypertarget{plafond-premier-exact}{%
\subsection{Plafond premier exact}\label{plafond-premier-exact}}

\textbf{Comme un premier jeu.} On ferme la liste de briques sans en
oublier une.

\textbf{Symbole.} X=⌊e\^{}(2S)⌋

\textbf{Sens et domaine.} Encadrement rationnel Taylor
inférieur/supérieur avec même partie entière.

\textbf{Garde.} Le plafond exact ne certifie pas les poids
logarithmiques ni la quadrature archimédienne.

\textbf{Ancrage.} Documents 67-68 et noyau du corpus ; définition locale
explicitée dans cette fiche ; définition pédagogique ; voir portée
précise dans ancrage et garde.

\hypertarget{forme-hermitienne}{%
\subsection{Forme hermitienne}\label{forme-hermitienne}}

\textbf{Comme un premier jeu.} Le miroir répond dans l'autre sens.

\textbf{Symbole.} q(g,f)=conj q(f,g)

\textbf{Sens et domaine.} Garantit Q(f)=q(f,f) réel ; symétrie de la
forme complète.

\textbf{Garde.} Réel ne signifie pas positif.

\textbf{Ancrage.} Documents 67-68 et noyau du corpus ; définition locale
explicitée dans cette fiche ; définition pédagogique ; voir portée
précise dans ancrage et garde.

\hypertarget{norme-dopuxe9rateur}{%
\subsection{Norme d'opérateur}\label{norme-dopuxe9rateur}}

\textbf{Comme un premier jeu.} Le plus grand agrandissement possible.

\textbf{Symbole.}
\textbar\textbar A\textbar\textbar=sup\_\textbar\textbar f\textbar\textbar=1\textbar\textbar Af\textbar\textbar{}

\textbf{Sens et domaine.} Dépend des normes de départ et d'arrivée ;
contrôlée par une borne effective.

\textbf{Garde.} Une petite entrée matricielle ne donne pas une petite
norme globale sans estimation.

\textbf{Ancrage.} Documents 67-68 et noyau du corpus ; définition locale
explicitée dans cette fiche ; définition pédagogique ; voir portée
précise dans ancrage et garde.

\hypertarget{reste-de-taylor-positif}{%
\subsection{Reste de Taylor positif}\label{reste-de-taylor-positif}}

\textbf{Comme un premier jeu.} On compte les petites piles encore
oubliées.

\textbf{Symbole.} e\^{}x=Σₖ≤N x\^{}k/k!+reste

\textbf{Sens et domaine.} Pour x≥0, termes positifs et borne géométrique
si x\textless N+2.

\textbf{Garde.} Le reste rationnel n'est pas une certification de toutes
les autres fonctions transcendantes.

\textbf{Ancrage.} Documents 67-68 et noyau du corpus ; définition locale
explicitée dans cette fiche ; définition pédagogique ; voir portée
précise dans ancrage et garde.

\hypertarget{calcul-complexe-des-formes}{%
\subsection{Calcul complexe des
formes}\label{calcul-complexe-des-formes}}

\textbf{Comme un premier jeu.} Chaque flèche tire sur toutes les autres.

\textbf{Symbole.} Q(Σc\_i f\_i)=c*G c

\textbf{Sens et domaine.} Une base réelle avec matrice réelle symétrique
se complexifie hermitiennement ; critère PSD identique.

\textbf{Garde.} Cela couvre les combinaisons complexes de cette famille,
pas tout paramètre A d'autres constructions du corpus.

\textbf{Ancrage.} Documents 67-68 et noyau du corpus ; définition locale
explicitée dans cette fiche ; définition pédagogique ; voir portée
précise dans ancrage et garde.

\hypertarget{parties-ruxe9elle-et-imaginaire}{%
\subsection{Parties réelle et
imaginaire}\label{parties-ruxe9elle-et-imaginaire}}

\textbf{Comme un premier jeu.} Deux dessins réels portent une flèche
complexe.

\textbf{Symbole.} f=u+iv, Q(f)=Q(u)+Q(v)

\textbf{Sens et domaine.} Pour une forme à coefficients réels
symétriques, les termes mixtes s'annulent.

\textbf{Garde.} La propriété dépend de la forme réelle ; les phases
d'autres paramètres peuvent rester dans ses coefficients.

\textbf{Ancrage.} Documents 67-68 et noyau du corpus ; définition locale
explicitée dans cette fiche ; définition pédagogique ; voir portée
précise dans ancrage et garde.

\hypertarget{duxe9faut-de-translation}{%
\subsection{Défaut de translation}\label{duxe9faut-de-translation}}

\textbf{Comme un premier jeu.} On compare le dessin à son petit
déplacement.

\textbf{Symbole.}
\textbar\textbar U\_xf-f\textbar\textbar²=2\textbar\textbar f\textbar\textbar²-2Re
h\_f(x)

\textbf{Sens et domaine.} Identité exacte L² ; pour H¹,
défaut≤x²\textbar\textbar f′\textbar\textbar².

\textbf{Garde.} Le défaut est positif, mais Q\_W soustrait aussi une
masse et les premiers.

\textbf{Ancrage.} Documents 67-68 et noyau du corpus ; définition locale
explicitée dans cette fiche ; définition pédagogique ; voir portée
précise dans ancrage et garde.

\hypertarget{matrice-de-correction}{%
\subsection{Matrice de correction}\label{matrice-de-correction}}

\textbf{Comme un premier jeu.} Deux retouches règlent deux balances.

\textbf{Symbole.} {[}{[}A,-A/2{]},{[}A,A/2{]}{]}

\textbf{Sens et domaine.} Pour φ paire et A=∫φe\^{}(x/2)\textgreater0,
moments de φ,φ′ donnent déterminant A².

\textbf{Garde.} L'ordre des lignes ± et les signes doivent rester
annoncés.

\textbf{Ancrage.} Documents 67-68 et noyau du corpus ; définition locale
explicitée dans cette fiche ; définition pédagogique ; voir portée
précise dans ancrage et garde.

\hypertarget{uxe9nergie-de-duxe9placement}{%
\subsection{Énergie de déplacement}\label{uxe9nergie-de-duxe9placement}}

\textbf{Comme un premier jeu.} Les petits glissements paient une
contribution positive.

\textbf{Symbole.} E\_arch=∫w\textbar\textbar U\_xf-f\textbar\textbar²

\textbf{Sens et domaine.} Forme positive, fermée par multiplicateur
Fourier a(τ).

\textbf{Garde.} E\_arch seule n'est pas Q\_W ; soustractions
arithmétiques restent présentes.

\textbf{Ancrage.} Documents 67-68 et noyau du corpus ; définition locale
explicitée dans cette fiche ; définition pédagogique ; voir portée
précise dans ancrage et garde.

\hypertarget{terme-croisuxe9}{%
\subsection{Terme croisé}\label{terme-croisuxe9}}

\textbf{Comme un premier jeu.} Une attache relie deux jouets.

\textbf{Symbole.} 2Re(c̄\_i c\_j q(f\_i,f\_j))

\textbf{Sens et domaine.} Partie indispensable d'une forme sur une
somme.

\textbf{Garde.} Sa phase peut rendre négative une somme à diagonales
positives.

\textbf{Ancrage.} Documents 67-68 et noyau du corpus ; définition locale
explicitée dans cette fiche ; définition pédagogique ; voir portée
précise dans ancrage et garde.

\hypertarget{duxe9coupage-tuxeate-et-queue}{%
\subsection{Découpage tête et
queue}\label{duxe9coupage-tuxeate-et-queue}}

\textbf{Comme un premier jeu.} Une boîte petite et tout ce qui reste.

\textbf{Symbole.} X\_S=V⊕\_R V⊥

\textbf{Sens et domaine.} Décomposition orthogonale R ; blocs de Q et
couplages définis.

\textbf{Garde.} La tête et la queue sont liées par G₀₁ ; elles ne sont
pas deux tests indépendants.

\textbf{Ancrage.} Documents 67-68 et noyau du corpus ; définition locale
explicitée dans cette fiche ; définition pédagogique ; voir portée
précise dans ancrage et garde.

\hypertarget{transposuxe9e-conjuguuxe9e}{%
\subsection{Transposée conjuguée}\label{transposuxe9e-conjuguuxe9e}}

\textbf{Comme un premier jeu.} On échange lignes et colonnes avec leur
miroir.

\textbf{Symbole.} C*=conj Cᵀ

\textbf{Sens et domaine.} Adjoint matriciel pour métrique euclidienne ;
congruence utilise C*GC.

\textbf{Garde.} Pour métrique B, l'adjoint opérateur est B\^{}(-1)C*B si
inverses valides.

\textbf{Ancrage.} Documents 67-68 et noyau du corpus ; définition locale
explicitée dans cette fiche ; définition pédagogique ; voir portée
précise dans ancrage et garde.

\hypertarget{codimension}{%
\subsection{Codimension}\label{codimension}}

\textbf{Comme un premier jeu.} Deux balances retirent deux libertés.

\textbf{Symbole.} codim E\_S=2

\textbf{Sens et domaine.} Deux moments indépendants continus coupent un
sous-espace fermé.

\textbf{Garde.} Codimension finie n'implique pas dimension finie du
domaine.

\textbf{Ancrage.} Documents 67-68 et noyau du corpus ; définition locale
explicitée dans cette fiche ; définition pédagogique ; voir portée
précise dans ancrage et garde.

\hypertarget{dualituxe9}{%
\subsection{Dualité}\label{dualituxe9}}

\textbf{Comme un premier jeu.} Un petit capteur lit un dessin.

\textbf{Symbole.} ℓ(f)

\textbf{Sens et domaine.} Fonctionnelle linéaire continue sur l'espace
choisi ; normes et adjoints associés.

\textbf{Garde.} Un capteur qui ne lit rien peut laisser des modes
invisibles.

\textbf{Ancrage.} Documents 67-68 et noyau du corpus ; définition locale
explicitée dans cette fiche ; définition pédagogique ; voir portée
précise dans ancrage et garde.

\hypertarget{suxe9rie-du-noyau}{%
\subsection{Série du noyau}\label{suxe9rie-du-noyau}}

\textbf{Comme un premier jeu.} Le poids est une pile de petites
décroissances.

\textbf{Symbole.} w(x)=Σₙ≥₀e\^{}(-β\_nx)

\textbf{Sens et domaine.} Série positive pour x\textgreater0 ; Tonelli
produit la série de a(τ).

\textbf{Garde.} Le calcul spectral emploie des termes positifs du noyau,
pas des zéros de ζ.

\textbf{Ancrage.} Documents 67-68 et noyau du corpus ; définition locale
explicitée dans cette fiche ; définition pédagogique ; voir portée
précise dans ancrage et garde.

\hypertarget{polynuxf4me-de-corruxe9lation}{%
\subsection{Polynôme de
corrélation}\label{polynuxf4me-de-corruxe9lation}}

\textbf{Comme un premier jeu.} Deux dessins polynomiaux se recouvrent
exactement.

\textbf{Symbole.} h\_ij(Sz)=polynôme rationnel, 0≤z≤2

\textbf{Sens et domaine.} Intégrale symbolique sur {[}z-1,1{]}, calcul
par fractions rationnelles.

\textbf{Garde.} Ce polynôme vaut sur l'intervalle indiqué ; hors
recouvrement la corrélation est zéro.

\textbf{Ancrage.} Documents 67-68 et noyau du corpus ; définition locale
explicitée dans cette fiche ; définition pédagogique ; voir portée
précise dans ancrage et garde.

\hypertarget{identituxe9-de-muxe9trique-du-moule}{%
\subsection{Identité de métrique du
moule}\label{identituxe9-de-muxe9trique-du-moule}}

\textbf{Comme un premier jeu.} Le moule donne aussi une règle exacte.

\textbf{Symbole.}
\textbar\textbar Pψ\textbar\textbar H¹²=\textbar\textbar ψ‴\textbar\textbar²+3/2\textbar\textbar ψ″\textbar\textbar²+9/16\textbar\textbar ψ′\textbar\textbar²+1/16\textbar\textbar ψ\textbar\textbar²

\textbf{Sens et domaine.} Intégration par parties dans H₀³ ; constantes
d'équivalence indépendantes de S.

\textbf{Garde.} Cette identité concerne la métrique B, pas Q\_W.

\textbf{Ancrage.} Document 68 ; V644\_LANGUE\_ET\_DOMAINES.md, L1-L7 ;
preuve auxiliaire interne ; portée et conditions dans L1-L7.

\hypertarget{conditionnement}{%
\subsection{Conditionnement}\label{conditionnement}}

\textbf{Comme un premier jeu.} Une règle presque écrasée rend les
lectures fragiles.

\textbf{Symbole.} κ(B)=λ\_max(B)/λ\_min(B)

\textbf{Sens et domaine.} Mesure spectrale pour B\textgreater0 euclidien
; amplifie les erreurs de normalisation.

\textbf{Garde.} Diffère des autres κ du corpus ; sa valeur est propre à
la base.

\textbf{Ancrage.} Documents 67-68 et noyau du corpus ; définition locale
explicitée dans cette fiche ; définition pédagogique ; voir portée
précise dans ancrage et garde.

\hypertarget{psd-et-marge}{%
\subsection{PSD et marge}\label{psd-et-marge}}

\textbf{Comme un premier jeu.} La balance tient ou garde une marge.

\textbf{Symbole.} G⪰0 versus G⪰cB

\textbf{Sens et domaine.} Première propriété admet directions nulles ;
seconde impose c\textgreater0 dans la métrique B.

\textbf{Garde.} La seconde est impossible pour toute E\_S avec B\_H¹ ;
la première reste ouverte.

\textbf{Ancrage.} Documents 67-68 et noyau du corpus ; définition locale
explicitée dans cette fiche ; définition pédagogique ; voir portée
précise dans ancrage et garde.

\hypertarget{uxe9quivalence-des-normes-du-moule}{%
\subsection{Équivalence des normes du
moule}\label{uxe9quivalence-des-normes-du-moule}}

\textbf{Comme un premier jeu.} Le moule fournit deux règles qui ne se
contredisent pas.

\textbf{Symbole.}
1/16\textbar\textbar ψ\textbar\textbar H³²≤\textbar\textbar Pψ\textbar\textbar H¹²≤3/2\textbar\textbar ψ\textbar\textbar H³²

\textbf{Sens et domaine.} Constantes exactes et indépendantes de S avec
norme H³ somme des dérivées.

\textbf{Garde.} Ces constantes ne sont pas des marges positives de Q\_W.

\textbf{Ancrage.} Document 68 ; V644\_LANGUE\_ET\_DOMAINES.md, L1-L7 ;
preuve auxiliaire interne ; portée et conditions dans L1-L7.

\hypertarget{iduxe9aux-de-jets-admissibles}{%
\subsection{Idéaux de jets
admissibles}\label{iduxe9aux-de-jets-admissibles}}

\textbf{Comme un premier jeu.} Plus de balances sont attachées aux deux
portes.

\textbf{Symbole.} I\_r=\{f∈A : ∫x\^{}j e\^{}(±x/2)f=0, j\textless r\}

\textbf{Sens et domaine.} Idéaux de convolution ; I₀=A ;
I\_r*I\_s⊂I\_(r+s), somme dans I\_min(r,s).

\textbf{Garde.} Les unités d'amplitude se multiplient ; convolution
n'implique pas positivité de Q\_W.

\textbf{Ancrage.}
3\_Preuves\_et\_calculs/preuves/V644\_LANGUE\_ET\_DOMAINES.md ; L6-L8 ;
relation démontrée dans V644, L6-L8.

\hypertarget{arithmuxe9tique-des-fonctions}{%
\subsection{Arithmétique des
fonctions}\label{arithmuxe9tique-des-fonctions}}

\textbf{Comme un premier jeu.} Quatre gestes ont chacun leur règle de
montage.

\textbf{Symbole.} addition, scalaire, produit, convolution

\textbf{Sens et domaine.} Somme/scalaire conservent I\_r ; convolution
additionne les ordres garantis ; produit ordinaire peut perdre I₁.

\textbf{Garde.} Ni D\_x ni U\_a ne sont des homomorphismes de
convolution dans le sens naïf.

\textbf{Ancrage.}
3\_Preuves\_et\_calculs/preuves/V644\_LANGUE\_ET\_DOMAINES.md ; L6-L8 ;
relation démontrée dans V644, L6-L8.

\hypertarget{duxe9rivation-pour-convolution}{%
\subsection{Dérivation pour
convolution}\label{duxe9rivation-pour-convolution}}

\textbf{Comme un premier jeu.} Le bouton position se partage entre deux
dessins.

\textbf{Symbole.} x(f*g)=(xf)*g+f*(xg)

\textbf{Sens et domaine.} Multiplication par x est une dérivation de
l'algèbre de convolution ; D\_x(f*g)=Df*g=f*Dg.

\textbf{Garde.} D\_x est dérivation du produit ordinaire, pas de la
convolution.

\textbf{Ancrage.}
3\_Preuves\_et\_calculs/preuves/V644\_LANGUE\_ET\_DOMAINES.md ; L6-L8 ;
relation démontrée dans V644, L6-L8.

\hypertarget{uxe9chelle-de-jets-par-p}{%
\subsection{Échelle de jets par P}\label{uxe9chelle-de-jets-par-p}}

\textbf{Comme un premier jeu.} Le moule ferme une balance de plus à
chaque porte.

\textbf{Symbole.} ℒ\_exp(Pf)=(z²-1/4)ℒ\_exp f

\textbf{Sens et domaine.} P:I\_r→I\_(r+1) sur graines Cc∞ ;
x:I\_(r+1)→I\_r ; D\_x:I\_r→I\_r.

\textbf{Garde.} Ordre garanti au moins ; certaines annulations
supplémentaires peuvent apparaître.

\textbf{Ancrage.}
3\_Preuves\_et\_calculs/preuves/V644\_LANGUE\_ET\_DOMAINES.md ; L6-L8 ;
relation démontrée dans V644, L6-L8.

\hypertarget{jets-ordinaires-en-zuxe9ro}{%
\subsection{Jets ordinaires en zéro}\label{jets-ordinaires-en-zuxe9ro}}

\textbf{Comme un premier jeu.} Une nouvelle balance lit les rythmes
lents.

\textbf{Symbole.} ℒ\_exp\^{}(j)f(0)=∫x\^{}j f(x)dx

\textbf{Sens et domaine.} L'annulation d'ordre q en 0 définit la queue
Y\_(q-1) ; D\_x gagne un ordre en 0 et P le conserve.

\textbf{Garde.} Ce point du test n'est pas un zéro critique de ζ.

\textbf{Ancrage.}
3\_Preuves\_et\_calculs/preuves/V644\_LANGUE\_ET\_DOMAINES.md ; L6-L8 ;
relation démontrée dans V644, L6-L8.

\hypertarget{quotient-de-jets}{%
\subsection{Quotient de jets}\label{quotient-de-jets}}

\textbf{Comme un premier jeu.} Deux petits carnets gardent les premiers
traits du dessin.

\textbf{Symbole.} A/I\_r ≃ ℂ{[}ε{]}/(ε\^{}r) ⊕ ℂ{[}ε{]}/(ε\^{}r)

\textbf{Sens et domaine.} Jets normalisés ℒ\_exp\^{}(j)f(±1/2)/j! ;
produit de coefficients tronqué ; codimension 2r.

\textbf{Garde.} Le quotient oublie tous les autres traits ;
r\textgreater1 ne couvre pas tous les tests de I₁.

\textbf{Ancrage.}
3\_Preuves\_et\_calculs/preuves/V644\_LANGUE\_ET\_DOMAINES.md ; L6-L8 ;
relation démontrée dans V644, L6-L8.

\hypertarget{iduxe9al-de-convolution}{%
\subsection{Idéal de convolution}\label{iduxe9al-de-convolution}}

\textbf{Comme un premier jeu.} Une porte fermée reste fermée quand on
colle un autre dessin.

\textbf{Symbole.} I\_r*A⊂I\_r

\textbf{Sens et domaine.} Sous-espace stable sous convolution par toute
graine ; somme conserve l'ordre minimal garanti.

\textbf{Garde.} Idéal ici se rapporte à convolution, pas au produit
ordinaire.

\textbf{Ancrage.}
3\_Preuves\_et\_calculs/preuves/V644\_LANGUE\_ET\_DOMAINES.md ; L6-L8 ;
relation démontrée dans V644, L6-L8.

\hypertarget{polynuxf4me-tronquuxe9}{%
\subsection{Polynôme tronqué}\label{polynuxf4me-tronquuxe9}}

\textbf{Comme un premier jeu.} Un carnet efface tous les traits trop
hauts.

\textbf{Symbole.} ℂ{[}ε{]}/(ε\^{}r)

\textbf{Sens et domaine.} Addition coefficient par coefficient et
produit Cauchy avec ε\^{}r=0 ; représentation exacte d'un jet d'ordre
r-1.

\textbf{Garde.} Ce ε formel est distinct de l'erreur ε\_N(T).

\textbf{Ancrage.}
3\_Preuves\_et\_calculs/preuves/V644\_LANGUE\_ET\_DOMAINES.md ; L6-L8 ;
relation démontrée dans V644, L6-L8.

\hypertarget{deux-moules-compatibles}{%
\subsection{Deux moules compatibles}\label{deux-moules-compatibles}}

\textbf{Comme un premier jeu.} Deux clés ferment les portes et les
rythmes lents.

\textbf{Symbole.} D\_x\^{}q P\^{}r ψ∈I\_r∩Y\_(q-1)

\textbf{Sens et domaine.} D\_x et P commutent ; ℒ\_exp multiplie par
(-z)\^{}q(z²-1/4)\^{}r ; support contenu dans celui de la graine.

\textbf{Garde.} La construction peut être dans une queue positive sans
couvrir la tête où le signe est inconnu.

\textbf{Ancrage.}
3\_Preuves\_et\_calculs/preuves/V644\_LANGUE\_ET\_DOMAINES.md ; L6-L8 ;
relation démontrée dans V644, L6-L8.

\hypertarget{niveau-5---les-contrats-de-recherche-et-leurs-limites}{%
\section{Niveau 5 - Les contrats de recherche et leurs
limites}\label{niveau-5---les-contrats-de-recherche-et-leurs-limites}}

\hypertarget{queue-de-lopuxe9rateur}{%
\subsection{Queue de l'opérateur}\label{queue-de-lopuxe9rateur}}

\textbf{Comme un premier jeu.} Les jouets restants ne peuvent plus faire
basculer seuls la balance.

\textbf{Symbole.}
\textbar\textbar K\_S\textbar V⊥\textbar\textbar\textless1

\textbf{Sens et domaine.} Une queue coercive explicite est construite
par moments ordinaires ; le taux pour la queue orthogonale de la base
brute f\_k reste à établir.

\textbf{Garde.} Ne pas identifier la queue de moments avec le complément
de span(f\_k).

\textbf{Ancrage.} Documents 67-68 et noyau du corpus ; définition locale
explicitée dans cette fiche ; définition pédagogique ; voir portée
précise dans ancrage et garde.

\hypertarget{ruxe9duction-par-queue-et-schur}{%
\subsection{Réduction par queue et
Schur}\label{ruxe9duction-par-queue-et-schur}}

\textbf{Comme un premier jeu.} On range l'infini et l'on garde ses
attaches exactes.

\textbf{Symbole.} G\_eff=G₀₀-G₀₁G₁₁\^{}(-1)G₁₀

\textbf{Sens et domaine.} Si le bloc de queue est uniformément positif,
signe global équivalent au Schur fini.

\textbf{Garde.} Une compression finie positive sans borne de queue ne
suffit pas.

\textbf{Ancrage.} Document 68 ; V644\_LANGUE\_ET\_DOMAINES.md, L1-L7 ;
preuve auxiliaire interne ; portée et conditions dans L1-L7.

\hypertarget{densituxe9-dune-famille}{%
\subsection{Densité d'une famille}\label{densituxe9-dune-famille}}

\textbf{Comme un premier jeu.} Les petits dessins peuvent approcher tous
les dessins.

\textbf{Symbole.} closure(∪V\_n)=E\_S

\textbf{Sens et domaine.} Ici densité de P{[}w\_S·polynômes{]} dans H¹ ;
puis dans X\_S par le cœur.

\textbf{Garde.} Dense ne signifie pas que le signe est positif sur cette
famille.

\textbf{Ancrage.} Document 68 ; V644\_LANGUE\_ET\_DOMAINES.md, L1-L7 ;
preuve auxiliaire interne ; portée et conditions dans L1-L7.

\hypertarget{quantificateur-global}{%
\subsection{Quantificateur global}\label{quantificateur-global}}

\textbf{Comme un premier jeu.} Chaque boîte et chaque taille doivent
être couvertes.

\textbf{Symbole.} ∀S\textgreater0, ∀n≥1

\textbf{Sens et domaine.} Une preuve globale exige tous les supports et
tous les tests via cœur et continuité.

\textbf{Garde.} Une extension de plafond couvre quelques indices
supplémentaires seulement.

\textbf{Ancrage.} Documents 67-68 et noyau du corpus ; définition locale
explicitée dans cette fiche ; définition pédagogique ; voir portée
précise dans ancrage et garde.

\hypertarget{bellman-des-contrats}{%
\subsection{Bellman des contrats}\label{bellman-des-contrats}}

\textbf{Comme un premier jeu.} On choisit un coup qui ouvre plusieurs
portes.

\textbf{Symbole.} V(s)=max\_a E{[}r(s,a)+γV(s′){]}

\textbf{Sens et domaine.} s porte objets, obligations et preuves ;
actions ont résultats conditionnels et coûts ; ici horizon fini γ=1.

\textbf{Garde.} Le score de recherche n'est ni une probabilité de RH ni
la forme de Weil.

\textbf{Ancrage.} Documents 67-68 et noyau du corpus ; définition locale
explicitée dans cette fiche ; définition pédagogique ; voir portée
précise dans ancrage et garde.

\hypertarget{portes-et-et-ou}{%
\subsection{Portes ET et OU}\label{portes-et-et-ou}}

\textbf{Comme un premier jeu.} Deux clés pour une porte, ou deux portes
possibles.

\textbf{Symbole.} ET: tous ; OU: un chemin valide

\textbf{Sens et domaine.} Graphe d'obligations avec complémentarités ;
utile pour les coups doubles.

\textbf{Garde.} Une flèche logique est autorisée seulement si son
contrat est prouvé.

\textbf{Ancrage.} Documents 67-68 et noyau du corpus ; définition locale
explicitée dans cette fiche ; définition pédagogique ; voir portée
précise dans ancrage et garde.

\hypertarget{jet}{%
\subsection{Jet}\label{jet}}

\textbf{Comme un premier jeu.} On garde le dessin, la pente et les
courbures.

\textbf{Symbole.} J\^{}mF(z)=(F,F′,\ldots,F\^{}(m))(z)

\textbf{Sens et domaine.} Données locales jusqu'à un ordre fini.

\textbf{Garde.} Un jet local n'indique pas à lui seul la branche globale
suivie.

\textbf{Ancrage.} Documents 67-68 et noyau du corpus ; définition locale
explicitée dans cette fiche ; définition pédagogique ; voir portée
précise dans ancrage et garde.

\hypertarget{branche}{%
\subsection{Branche}\label{branche}}

\textbf{Comme un premier jeu.} Un chemin local à travers une forêt.

\textbf{Symbole.} F(z,t)=0, z=z(t)

\textbf{Sens et domaine.} Si ∂\_zF≠0, fonction implicite locale ; près
d'un zéro multiple il faut un autre modèle.

\textbf{Garde.} Continuer à travers une collision nécessite conditions
et choix ; sujet encore ouvert du corpus.

\textbf{Ancrage.} Documents 67-68 et noyau du corpus ; définition locale
explicitée dans cette fiche ; définition pédagogique ; voir portée
précise dans ancrage et garde.

\hypertarget{collision-de-zuxe9ros}{%
\subsection{Collision de zéros}\label{collision-de-zuxe9ros}}

\textbf{Comme un premier jeu.} Deux chemins se touchent.

\textbf{Symbole.} F=∂\_zF=0

\textbf{Sens et domaine.} Système local de dégénérescence ; jets,
discriminant ou préparation locale à spécifier.

\textbf{Garde.} Une collision certifiée locale ne contrôle pas tous les
passages globaux.

\textbf{Ancrage.} Documents 67-68 et noyau du corpus ; définition locale
explicitée dans cette fiche ; définition pédagogique ; voir portée
précise dans ancrage et garde.

\hypertarget{discriminant}{%
\subsection{Discriminant}\label{discriminant}}

\textbf{Comme un premier jeu.} Une alarme quand deux racines se
rejoignent.

\textbf{Symbole.} Disc(P)=0

\textbf{Sens et domaine.} Pour un polynôme, détecte racines multiples ;
extension locale à préciser.

\textbf{Garde.} Ne pas appliquer la formule polynomiale directement à
une fonction entière quelconque.

\textbf{Ancrage.} Documents 67-68 et noyau du corpus ; définition locale
explicitée dans cette fiche ; définition pédagogique ; voir portée
précise dans ancrage et garde.

\hypertarget{coefficient-complexe-a}{%
\subsection{Coefficient complexe A}\label{coefficient-complexe-a}}

\textbf{Comme un premier jeu.} Une flèche peut tourner et grandir.

\textbf{Symbole.} A∈ℂ

\textbf{Sens et domaine.} Paramètre de certaines familles du corpus ;
\textbar A\textbar²≥0, mais les termes croisés gardent la phase.

\textbf{Garde.} Un argument limité à A réel positif ne couvre pas A
complexe arbitraire.

\textbf{Ancrage.} Documents 67-68 et noyau du corpus ; définition locale
explicitée dans cette fiche ; définition pédagogique ; voir portée
précise dans ancrage et garde.

\hypertarget{trace-et-correspondances-de-weil}{%
\subsection{Trace et correspondances de
Weil}\label{trace-et-correspondances-de-weil}}

\textbf{Comme un premier jeu.} Une règle relie les allers et les
retours.

\textbf{Symbole.} σ(ξξ′)\textgreater0 dans le cadre géométrique

\textbf{Sens et domaine.} Positivité des correspondances dans la preuve
sur les courbes sur corps finis ; source Wei48.

\textbf{Garde.} Aucun transfert de cette géométrie à ζ des entiers n'est
démontré.

\textbf{Ancrage.} Documents 67-68 et noyau du corpus ; définition locale
explicitée dans cette fiche ; définition pédagogique ; voir portée
précise dans ancrage et garde.

\hypertarget{frobenius}{%
\subsection{Frobenius}\label{frobenius}}

\textbf{Comme un premier jeu.} On répète un geste exact dans un monde
fini.

\textbf{Symbole.} Frob:x↦x\^{}q

\textbf{Sens et domaine.} Morphisme en caractéristique p ; action
cohomologique avec trace dans le cadre approprié.

\textbf{Garde.} Un opérateur virtuel aux mêmes valeurs propres n'est pas
automatiquement un morphisme.

\textbf{Ancrage.} Documents 67-68 et noyau du corpus ; définition locale
explicitée dans cette fiche ; définition pédagogique ; voir portée
précise dans ancrage et garde.

\hypertarget{riemannroch}{%
\subsection{Riemann--Roch}\label{riemannroch}}

\textbf{Comme un premier jeu.} On compte les dessins autorisés et les
contraintes.

\textbf{Symbole.} ℓ(D)-ℓ(K-D)=deg D+1-g

\textbf{Sens et domaine.} Théorème pour un diviseur sur une courbe lisse
projective ; géométrie nécessaire.

\textbf{Garde.} Un bilan formel de dimensions ne construit pas une
courbe pour ζ.

\textbf{Ancrage.} Documents 67-68 et noyau du corpus ; définition locale
explicitée dans cette fiche ; définition pédagogique ; voir portée
précise dans ancrage et garde.

\hypertarget{marge-spectrale}{%
\subsection{Marge spectrale}\label{marge-spectrale}}

\textbf{Comme un premier jeu.} On cherche le plus petit espace sous la
balance.

\textbf{Symbole.} inf spec(A)

\textbf{Sens et domaine.} Dépend de l'opérateur et de sa métrique ;
existence ou borne certifiée à distinguer.

\textbf{Garde.} La marge R éventuelle peut dépendre de S ; la marge H¹
strictement positive est impossible ici.

\textbf{Ancrage.} Documents 67-68 et noyau du corpus ; définition locale
explicitée dans cette fiche ; définition pédagogique ; voir portée
précise dans ancrage et garde.

\hypertarget{certificat}{%
\subsection{Certificat}\label{certificat}}

\textbf{Comme un premier jeu.} On garde une règle qui encadre le
résultat.

\textbf{Symbole.} borne d'erreur + décision

\textbf{Sens et domaine.} Calcul rationnel exact ou intervalle validé et
toutes conditions contrôlées.

\textbf{Garde.} Deux précisions concordantes ne remplacent pas une borne
d'erreur.

\textbf{Ancrage.} Documents 67-68 et noyau du corpus ; définition locale
explicitée dans cette fiche ; définition pédagogique ; voir portée
précise dans ancrage et garde.

\hypertarget{calcul-par-intervalles}{%
\subsection{Calcul par intervalles}\label{calcul-par-intervalles}}

\textbf{Comme un premier jeu.} Le résultat habite une petite boîte
fermée.

\textbf{Symbole.} x∈{[}a,b{]}

\textbf{Sens et domaine.} Arithmétique à arrondi dirigé et propagation
rigoureuse des inclusions.

\textbf{Garde.} Une valeur à beaucoup de chiffres n'est pas un
intervalle validé.

\textbf{Ancrage.} Documents 67-68 et noyau du corpus ; définition locale
explicitée dans cette fiche ; définition pédagogique ; voir portée
précise dans ancrage et garde.

\hypertarget{stabilituxe9-dun-signe}{%
\subsection{Stabilité d'un signe}\label{stabilituxe9-dun-signe}}

\textbf{Comme un premier jeu.} La balance ne bascule pas malgré une
petite secousse.

\textbf{Symbole.} \textbar erreur\textbar\textless marge

\textbf{Sens et domaine.} Perturbation contrôlée d'une forme ou d'un
spectre.

\textbf{Garde.} Le couplage et le conditionnement amplifient l'erreur ;
il faut les borner.

\textbf{Ancrage.} Documents 67-68 et noyau du corpus ; définition locale
explicitée dans cette fiche ; définition pédagogique ; voir portée
précise dans ancrage et garde.

\hypertarget{invariant-symbolique}{%
\subsection{Invariant symbolique}\label{invariant-symbolique}}

\textbf{Comme un premier jeu.} Un signe survit à la traduction.

\textbf{Symbole.} I(Tx)=I(x)

\textbf{Sens et domaine.} Une quantité conservée par une transformation
définie avec ses conditions.

\textbf{Garde.} Un même symbole écrit deux fois n'est pas un invariant.

\textbf{Ancrage.} Documents 67-68 et noyau du corpus ; définition locale
explicitée dans cette fiche ; définition pédagogique ; voir portée
précise dans ancrage et garde.

\hypertarget{invariant-relationnel}{%
\subsection{Invariant relationnel}\label{invariant-relationnel}}

\textbf{Comme un premier jeu.} Les attaches survivent même si les pièces
changent.

\textbf{Symbole.} R(x,y)⇒R(Tx,Ty)

\textbf{Sens et domaine.} Moments nuls, incidence de sous-espaces,
congruence simultanée, ordre d'un drapeau.

\textbf{Garde.} Une relation peut être conservée sans conserver les
valeurs numériques.

\textbf{Ancrage.} Documents 67-68 et noyau du corpus ; définition locale
explicitée dans cette fiche ; définition pédagogique ; voir portée
précise dans ancrage et garde.

\hypertarget{grammaire-des-passages}{%
\subsection{Grammaire des passages}\label{grammaire-des-passages}}

\textbf{Comme un premier jeu.} Les mots savent quelles phrases ils
peuvent former.

\textbf{Symbole.} objet : domaine → codomaine ; garde ; coût

\textbf{Sens et domaine.} Contrat de traduction typé, preuve et perte
séparées ; validateur programmé.

\textbf{Garde.} La validation vérifie le contrat codé, pas une preuve
automatique de toutes les mathématiques.

\textbf{Ancrage.} Documents 67-68 et noyau du corpus ; définition locale
explicitée dans cette fiche ; définition pédagogique ; voir portée
précise dans ancrage et garde.

\hypertarget{utilituxe9-de-recherche}{%
\subsection{Utilité de recherche}\label{utilituxe9-de-recherche}}

\textbf{Comme un premier jeu.} On compte les portes réellement ouvertes.

\textbf{Symbole.} U\_decision≠U\_a

\textbf{Sens et domaine.} Choix éditorial de gains sur obligations ;
aucune unité physique.

\textbf{Garde.} Ne pas confondre U\_decision avec la translation U\_a.

\textbf{Ancrage.} Documents 67-68 et noyau du corpus ; définition locale
explicitée dans cette fiche ; définition pédagogique ; voir portée
précise dans ancrage et garde.

\hypertarget{les-diffuxe9rents-g}{%
\subsection{Les différents G}\label{les-diffuxe9rents-g}}

\textbf{Comme un premier jeu.} Quatre boîtes ne portent plus la même
étiquette.

\textbf{Symbole.} G\_log(s), G\_NB, G\_W, B

\textbf{Sens et domaine.} G\_log désigne un logarithme local selon le
document ; G\_NB un Gram ; G\_W une compression ; B la métrique externe.

\textbf{Garde.} La désambiguïsation exige le document ; le symbole G
seul ne donne pas le domaine.

\textbf{Ancrage.} Documents 67-68 et noyau du corpus ; définition locale
explicitée dans cette fiche ; définition pédagogique ; voir portée
précise dans ancrage et garde.

\hypertarget{les-diffuxe9rents-d}{%
\subsection{Les différents D}\label{les-diffuxe9rents-d}}

\textbf{Comme un premier jeu.} Trois gestes portent des noms complets.

\textbf{Symbole.} D\_x, D\_mult, D\_Frob

\textbf{Sens et domaine.} D\_x dérive ; D\_mult mesure un défaut
multiplicatif ; un logarithme de Frobenius exige une branche et un
domaine.

\textbf{Garde.} Ne pas composer ces symboles comme s'ils étaient la même
machine.

\textbf{Ancrage.} Documents 67-68 et noyau du corpus ; définition locale
explicitée dans cette fiche ; définition pédagogique ; voir portée
précise dans ancrage et garde.

\hypertarget{projection}{%
\subsection{Projection}\label{projection}}

\textbf{Comme un premier jeu.} On garde la partie qui entre dans une
boîte.

\textbf{Symbole.} Π\_V f

\textbf{Sens et domaine.} Projection orthogonale dans une métrique
annoncée ; Π\_V²=Π\_V.

\textbf{Garde.} Ce P/Π n'est pas le moule D²-1/4 ; projection H¹ et
projection R diffèrent.

\textbf{Ancrage.} Documents 67-68 et noyau du corpus ; définition locale
explicitée dans cette fiche ; définition pédagogique ; voir portée
précise dans ancrage et garde.

\hypertarget{polynuxf4me}{%
\subsection{Polynôme}\label{polynuxf4me}}

\textbf{Comme un premier jeu.} Un dessin fabriqué par quelques briques
de puissances.

\textbf{Symbole.} P\_poly(u)=Σc\_k u\^{}k

\textbf{Sens et domaine.} Fonction algébrique de variable déclarée ;
zêta polynomiale en corps fini selon le contexte.

\textbf{Garde.} Ce P n'est pas l'opérateur différentiel.

\textbf{Ancrage.} Documents 67-68 et noyau du corpus ; définition locale
explicitée dans cette fiche ; définition pédagogique ; voir portée
précise dans ancrage et garde.

\hypertarget{les-diffuxe9rentes-ux3ba}{%
\subsection{Les différentes κ}\label{les-diffuxe9rentes-ux3ba}}

\textbf{Comme un premier jeu.} Deux petites étiquettes doivent garder
leur boîte.

\textbf{Symbole.} κ\_Li, κ\_pli

\textbf{Sens et domaine.} Coefficient ou paramètre selon le document ;
définition locale obligatoire.

\textbf{Garde.} Aucune égalité entre les κ n'est fournie par leur seule
lettre.

\textbf{Ancrage.} Documents 67-68 et noyau du corpus ; définition locale
explicitée dans cette fiche ; définition pédagogique ; voir portée
précise dans ancrage et garde.

\hypertarget{tonelli-et-uxe9change-positif}{%
\subsection{Tonelli et échange
positif}\label{tonelli-et-uxe9change-positif}}

\textbf{Comme un premier jeu.} On compte les mêmes briques dans deux
ordres.

\textbf{Symbole.} ∫Σ f\_n=Σ∫f\_n si f\_n≥0

\textbf{Sens et domaine.} Autorise l'échange pour fonctions mesurables
non négatives, même si valeur infinie.

\textbf{Garde.} Pour des termes signés, il faut d'autres hypothèses de
domination/convergence.

\textbf{Ancrage.} Documents 67-68 et noyau du corpus ; définition locale
explicitée dans cette fiche ; définition pédagogique ; voir portée
précise dans ancrage et garde.

\hypertarget{plancherel}{%
\subsection{Plancherel}\label{plancherel}}

\textbf{Comme un premier jeu.} La même chanson a la même taille sur deux
portées.

\textbf{Symbole.}
\textbar\textbar f\textbar\textbar₂=\textbar\textbar f̂\textbar\textbar₂

\textbf{Sens et domaine.} Isométrie Fourier L² avec normalisation
unitaire.

\textbf{Garde.} L'égalité L² seule ne donne pas les domaines des formes
plus fortes.

\textbf{Ancrage.} Documents 67-68 et noyau du corpus ; définition locale
explicitée dans cette fiche ; définition pédagogique ; voir portée
précise dans ancrage et garde.

\hypertarget{commutateur}{%
\subsection{Commutateur}\label{commutateur}}

\textbf{Comme un premier jeu.} Deux gestes changent selon leur ordre.

\textbf{Symbole.} {[}P,M\_τ{]}=M\_τ(2iτD-τ²)

\textbf{Sens et domaine.} P M\_τψ=M\_τ(Pψ+2iτψ′-τ²ψ).

\textbf{Garde.} Il explique pourquoi moduler la graine puis appliquer P
n'est pas moduler le test seul.

\textbf{Ancrage.} Documents 67-68 et noyau du corpus ; définition locale
explicitée dans cette fiche ; définition pédagogique ; voir portée
précise dans ancrage et garde.

\hypertarget{pont-ruxe9gularisuxe9}{%
\subsection{Pont régularisé}\label{pont-ruxe9gularisuxe9}}

\textbf{Comme un premier jeu.} On déplace, lisse, range et passe dans le
moule.

\textbf{Symbole.} T=P{[}χ(η\_δ*Γ\_Sg){]}

\textbf{Sens et domaine.} Γ\_S est le décalage de l'entrée prolongée par
zéro ; sortie admissible ; coût δ\^{}(-3) en H¹ sous cutoffs fixés.

\textbf{Garde.} Le coût et les limites doivent être uniformisés avant de
conclure globalement.

\textbf{Ancrage.} Documents 67-68 et noyau du corpus ; définition locale
explicitée dans cette fiche ; définition pédagogique ; voir portée
précise dans ancrage et garde.

\hypertarget{correction-des-moments}{%
\subsection{Correction des moments}\label{correction-des-moments}}

\textbf{Comme un premier jeu.} Deux petites retouches remettent les
balances à zéro.

\textbf{Symbole.} f↦f-aφ-bφ′

\textbf{Sens et domaine.} Matrice des moments de φ et φ′ inversible pour
φ paire non nulle positive ; correction continue.

\textbf{Garde.} Il faut un support intérieur commun et des coefficients
calculés, pas une correction abstraite gratuite.

\textbf{Ancrage.} Documents 67-68 et noyau du corpus ; définition locale
explicitée dans cette fiche ; définition pédagogique ; voir portée
précise dans ancrage et garde.

\hypertarget{rabais-bellman-ux3b3}{%
\subsection{Rabais Bellman γ}\label{rabais-bellman-ux3b3}}

\textbf{Comme un premier jeu.} Les coups éloignés peuvent peser moins
lourd.

\textbf{Symbole.} 0≤γ≤1

\textbf{Sens et domaine.} Paramètre de décision, ici γ=1 à horizon fini
pour les contrats.

\textbf{Garde.} γ\_decision n'est ni Γ fonction spéciale ni γ\_E
constante d'Euler.

\textbf{Ancrage.} Documents 67-68 et noyau du corpus ; définition locale
explicitée dans cette fiche ; définition pédagogique ; voir portée
précise dans ancrage et garde.

\hypertarget{principe-variationnel}{%
\subsection{Principe variationnel}\label{principe-variationnel}}

\textbf{Comme un premier jeu.} On cherche le dessin qui fait descendre
le plus la balance.

\textbf{Symbole.} λ\_min=inf Q(f)/R(f)

\textbf{Sens et domaine.} Quotient de Rayleigh dans la métrique annoncée
; restriction finie donne un diagnostic.

\textbf{Garde.} Une infimum sur un petit sous-espace peut manquer une
direction négative ailleurs.

\textbf{Ancrage.} Documents 67-68 et noyau du corpus ; définition locale
explicitée dans cette fiche ; définition pédagogique ; voir portée
précise dans ancrage et garde.

\hypertarget{erreur-de-famille}{%
\subsection{Erreur de famille}\label{erreur-de-famille}}

\textbf{Comme un premier jeu.} Les petites dettes disparaissent dans
toutes les boîtes.

\textbf{Symbole.} G\_n≥-ε\_n B\_n, ε\_n→0

\textbf{Sens et domaine.} Avec famille imbriquée dense et continuité,
cette borne pour tout n implique Q\_W≥0.

\textbf{Garde.} Une suite de bornes non certifiées ou sans tous les n ne
remplit pas le contrat.

\textbf{Ancrage.} Documents 67-68 et noyau du corpus ; définition locale
explicitée dans cette fiche ; définition pédagogique ; voir portée
précise dans ancrage et garde.

\hypertarget{duxe9faut-de-transport}{%
\subsection{Défaut de transport}\label{duxe9faut-de-transport}}

\textbf{Comme un premier jeu.} On mesure ce qu'une traduction abîme.

\textbf{Symbole.} \textbar q(Tf,Tg)-q(f,g)\textbar{}

\textbf{Sens et domaine.} Différence de formes sur domaines compatibles,
bornée par normes et erreurs de T.

\textbf{Garde.} Une transformation non isométrique exige cette
comptabilité.

\textbf{Ancrage.} Documents 67-68 et noyau du corpus ; définition locale
explicitée dans cette fiche ; définition pédagogique ; voir portée
précise dans ancrage et garde.

\hypertarget{indice-nuxe9gatif}{%
\subsection{Indice négatif}\label{indice-nuxe9gatif}}

\textbf{Comme un premier jeu.} Combien de directions font pencher vers
le bas.

\textbf{Symbole.} ind₋(Q)

\textbf{Sens et domaine.} Dimension maximale d'un sous-espace où
Q\textless0 ; finie à S fixé via I-compact.

\textbf{Garde.} Fini ne signifie pas zéro et ne couvre pas tous les
supports.

\textbf{Ancrage.} Documents 67-68 et noyau du corpus ; définition locale
explicitée dans cette fiche ; définition pédagogique ; voir portée
précise dans ancrage et garde.

\hypertarget{projection-fruxe9quentielle}{%
\subsection{Projection
fréquentielle}\label{projection-fruxe9quentielle}}

\textbf{Comme un premier jeu.} On écoute seulement les rythmes assez
lents.

\textbf{Symbole.} Π\_T=ℱ\^{}(-1)1\_{[}-T,T{]}ℱ

\textbf{Sens et domaine.} Projection orthogonale L² ; queue contrôlée
par 1+a(T) dans R.

\textbf{Garde.} Elle ne conserve pas le support compact ni les moments ;
c'est un outil d'estimation, pas un test admissible automatique.

\textbf{Ancrage.} Documents 67-68 et noyau du corpus ; définition locale
explicitée dans cette fiche ; définition pédagogique ; voir portée
précise dans ancrage et garde.

\hypertarget{support-et-temps-logarithmique}{%
\subsection{Support et temps
logarithmique}\label{support-et-temps-logarithmique}}

\textbf{Comme un premier jeu.} On transforme les tailles en positions.

\textbf{Symbole.} x=log u, D=diam supp f

\textbf{Sens et domaine.} Nombre de premiers en jeu lié à e\^{}D ;
support centré n'est qu'une convention.

\textbf{Garde.} Ne pas identifier ce x à un temps physique ni réduire le
diamètre à la position du centre.

\textbf{Ancrage.} Documents 67-68 et noyau du corpus ; définition locale
explicitée dans cette fiche ; définition pédagogique ; voir portée
précise dans ancrage et garde.

\hypertarget{rigidituxe9-de-hamburger}{%
\subsection{Rigidité de Hamburger}\label{rigidituxe9-de-hamburger}}

\textbf{Comme un premier jeu.} Les règles complètes forcent une chanson
précise.

\textbf{Symbole.} équation fonctionnelle + série + croissance

\textbf{Sens et domaine.} Dans le cadre Ham21, les hypothèses complètes
imposent F=cζ ; voir source corrigée document 67.

\textbf{Garde.} L'équation fonctionnelle isolée ne suffit pas.

\textbf{Ancrage.} Documents 67-68 et noyau du corpus ; définition locale
explicitée dans cette fiche ; définition pédagogique ; voir portée
précise dans ancrage et garde.

\hypertarget{contre-exemple-local}{%
\subsection{Contre-exemple local}\label{contre-exemple-local}}

\textbf{Comme un premier jeu.} Un seul jouet peut montrer qu'une règle
casse.

\textbf{Symbole.} ∃x : propriété fausse

\textbf{Sens et domaine.} Réfute la proposition avec ses hypothèses
exactes ; un témoin numérique demande une certification.

\textbf{Garde.} Fermer un candidat ne ferme pas RH ni toutes les autres
routes.

\textbf{Ancrage.} Documents 67-68 et noyau du corpus ; définition locale
explicitée dans cette fiche ; définition pédagogique ; voir portée
précise dans ancrage et garde.

\hypertarget{uxe9tat-uxe9pistuxe9mique}{%
\subsection{État épistémique}\label{uxe9tat-uxe9pistuxe9mique}}

\textbf{Comme un premier jeu.} La carte garde ce qui a été construit.

\textbf{Symbole.} s=(objets, preuves, obligations)

\textbf{Sens et domaine.} État du protocole, résultats et limites,
horodatage avant les expériences.

\textbf{Garde.} Un résultat appris ensuite ne doit pas être présenté
comme une prévision préalable.

\textbf{Ancrage.} Documents 67-68 et noyau du corpus ; définition locale
explicitée dans cette fiche ; définition pédagogique ; voir portée
précise dans ancrage et garde.

\hypertarget{coup-double}{%
\subsection{Coup double}\label{coup-double}}

\textbf{Comme un premier jeu.} Une clé ouvre deux portes liées.

\textbf{Symbole.} action→deux contrats vérifiés

\textbf{Sens et domaine.} Ici P : moments+norme ; parité :
réduction+contre-épreuve ; R : domaine+queue finie.

\textbf{Garde.} Deux noms nouveaux ne comptent pas comme deux preuves.

\textbf{Ancrage.} Documents 67-68 et noyau du corpus ; définition locale
explicitée dans cette fiche ; définition pédagogique ; voir portée
précise dans ancrage et garde.

\hypertarget{crituxe8re-nymanbeurling}{%
\subsection{Critère Nyman--Beurling}\label{crituxe8re-nymanbeurling}}

\textbf{Comme un premier jeu.} Approcher un dessin simple suffit dans le
bon monde.

\textbf{Symbole.} 1∈closure span\{ρ\_θ avec moments\}

\textbf{Sens et domaine.} Équivalence à RH sous l'énoncé et la norme
exacts de Beurling ; voir document 67.

\textbf{Garde.} La densité de toute entrée L² après régularisation n'est
pas la densité des entrées NB.

\textbf{Ancrage.} Documents 67-68 et noyau du corpus ; définition locale
explicitée dans cette fiche ; définition pédagogique ; voir portée
précise dans ancrage et garde.

\hypertarget{erreur-dapproximation}{%
\subsection{Erreur d'approximation}\label{erreur-dapproximation}}

\textbf{Comme un premier jeu.} Les petites boîtes manquent encore une
partie du dessin.

\textbf{Symbole.} \textbar\textbar(I-Π\_N)j\_S\textbar\textbar{}

\textbf{Sens et domaine.} Norme de l'injection compacte après projection
finie dans la métrique cible.

\textbf{Garde.} Existence d'une limite zéro ne fournit pas son taux
certifié.

\textbf{Ancrage.} Documents 67-68 et noyau du corpus ; définition locale
explicitée dans cette fiche ; définition pédagogique ; voir portée
précise dans ancrage et garde.

\hypertarget{budget-arithmuxe9tique-luxb2}{%
\subsection{Budget arithmétique L²}\label{budget-arithmuxe9tique-luxb2}}

\textbf{Comme un premier jeu.} On compte tout le poids qui peut tirer la
balance.

\textbf{Symbole.} M\_S=C₀+1+2Σₚᵐ≤e²ˢ log p·p\^{}(-m/2)

\textbf{Sens et domaine.} Borne de
\textbar(C₀+1)⟨f,g⟩+q\_pr(f,g)\textbar≤M\_S\textbar\textbar f\textbar\textbar\textbar\textbar g\textbar\textbar.

\textbf{Garde.} Budget valable à S fixé, croît avec le support ; ne
donne pas le signe.

\textbf{Ancrage.} Document 68 ; V644\_LANGUE\_ET\_DOMAINES.md, L1-L7 ;
preuve auxiliaire interne ; portée et conditions dans L1-L7.

\hypertarget{porte-de-queue-positive}{%
\subsection{Porte de queue positive}\label{porte-de-queue-positive}}

\textbf{Comme un premier jeu.} Les jouets restants ne tirent plus assez
fort.

\textbf{Symbole.} \textbar\textbar j\_S\textbar V⊥\textbar\textbar²
M\_S\textless1

\textbf{Sens et domaine.} La queue de moments Y\_N satisfait la
condition par la borne ε\_N(T) ; existence et paramètres explicites
prouvés.

\textbf{Garde.} Queue positive seule ne rend pas la tête positive ; le
Schur complet demeure nécessaire.

\textbf{Ancrage.} Documents 67-68 et noyau du corpus ; définition locale
explicitée dans cette fiche ; définition pédagogique ; voir portée
précise dans ancrage et garde.

\hypertarget{obligation-globale-de-signe}{%
\subsection{Obligation globale de
signe}\label{obligation-globale-de-signe}}

\textbf{Comme un premier jeu.} Toutes les dernières balances doivent
tenir.

\textbf{Symbole.} ∀S, G\_eff,S⪰0

\textbf{Sens et domaine.} Après queue certifiée, résoudre le Schur fini
pour chaque S, ou une loi uniforme couvrante.

\textbf{Garde.} Ni l'existence du Schur ni les expériences finies
n'acquittent cette obligation.

\textbf{Ancrage.} Documents 67-68 et noyau du corpus ; définition locale
explicitée dans cette fiche ; définition pédagogique ; voir portée
précise dans ancrage et garde.

\hypertarget{ruxe9solvante}{%
\subsection{Résolvante}\label{ruxe9solvante}}

\textbf{Comme un premier jeu.} On cherche comment une machine répond à
une sollicitation.

\textbf{Symbole.} (A-zI)\^{}(-1)

\textbf{Sens et domaine.} Définie pour z hors du spectre avec domaines
fixés ; utile au contrôle de queue.

\textbf{Garde.} Une résolvante postulée sans exclusion spectrale est
circulaire.

\textbf{Ancrage.} Documents 67-68 et noyau du corpus ; définition locale
explicitée dans cette fiche ; définition pédagogique ; voir portée
précise dans ancrage et garde.

\hypertarget{spectre-essentiel}{%
\subsection{Spectre essentiel}\label{spectre-essentiel}}

\textbf{Comme un premier jeu.} La grande foule des rythmes reste autour
d'un.

\textbf{Symbole.} spec\_ess(I-K)=\{1\}

\textbf{Sens et domaine.} Pour K compact autoadjoint sur un Hilbert
infini ; seules exceptions discrètes ailleurs.

\textbf{Garde.} Cela autorise des valeurs négatives isolées ; aucun
signe global automatique.

\textbf{Ancrage.} Documents 67-68 et noyau du corpus ; définition locale
explicitée dans cette fiche ; définition pédagogique ; voir portée
précise dans ancrage et garde.

\hypertarget{rang-fini}{%
\subsection{Rang fini}\label{rang-fini}}

\textbf{Comme un premier jeu.} Quelques gestes suffisent pour cette
machine.

\textbf{Symbole.} A\_N=Σₖ≤N u\_k⟨v\_k,·⟩

\textbf{Sens et domaine.} Image de dimension finie ; approxime en norme
les compacts sur Hilbert.

\textbf{Garde.} La possibilité d'approximer ne fournit pas les
coefficients ni leur erreur certifiée.

\textbf{Ancrage.} Documents 67-68 et noyau du corpus ; définition locale
explicitée dans cette fiche ; définition pédagogique ; voir portée
précise dans ancrage et garde.

\hypertarget{suxe9rie-de-neumann}{%
\subsection{Série de Neumann}\label{suxe9rie-de-neumann}}

\textbf{Comme un premier jeu.} On corrige le dessin par petits retours
de plus en plus faibles.

\textbf{Symbole.} (I-K)\^{}(-1)=Σₙ≥₀K\^{}n si
\textbar\textbar K\textbar\textbar\textless1

\textbf{Sens et domaine.} Inverse borné avec contrôle
1/(1-\textbar\textbar K\textbar\textbar).

\textbf{Garde.} Utilisable pour le bloc de queue seulement après une
borne stricte.

\textbf{Ancrage.} Document 68 ; V644\_LANGUE\_ET\_DOMAINES.md, L1-L7 ;
preuve auxiliaire interne ; portée et conditions dans L1-L7.

\hypertarget{muxe9trique-induxe9pendante}{%
\subsection{Métrique indépendante}\label{muxe9trique-induxe9pendante}}

\textbf{Comme un premier jeu.} La règle doit exister avant la décision
de la balance.

\textbf{Symbole.} R ou B définies sans Q≥0

\textbf{Sens et domaine.} Norme externe positive démontrée sans
hypothèse RH ; permet un Riesz signé.

\textbf{Garde.} Définir la norme comme √Q\_W avant d'avoir prouvé Q\_W≥0
serait circulaire.

\textbf{Ancrage.} Documents 67-68 et noyau du corpus ; définition locale
explicitée dans cette fiche ; définition pédagogique ; voir portée
précise dans ancrage et garde.

\hypertarget{drapeau-par-secteurs}{%
\subsection{Drapeau par secteurs}\label{drapeau-par-secteurs}}

\textbf{Comme un premier jeu.} On range d'abord selon les deux miroirs.

\textbf{Symbole.} V\_pair⊕V\_impair

\textbf{Sens et domaine.} Schur peut être calculé par bloc de parité ;
pivot interprété dans son drapeau.

\textbf{Garde.} Comparer à un ordre mélangé nécessite garder la même
extension ou signaler un nouveau drapeau.

\textbf{Ancrage.} Documents 67-68 et noyau du corpus ; définition locale
explicitée dans cette fiche ; définition pédagogique ; voir portée
précise dans ancrage et garde.

\hypertarget{noyau-de-basse-fruxe9quence}{%
\subsection{Noyau de basse
fréquence}\label{noyau-de-basse-fruxe9quence}}

\textbf{Comme un premier jeu.} Les rythmes lents se recouvrent dans la
boîte.

\textbf{Symbole.} sin(T(x-y))/(π(x-y))

\textbf{Sens et domaine.} Noyau de 1\_{[}-S,S{]}Π\_T1\_{[}-S,S{]},
compact positif à support fixé.

\textbf{Garde.} Borne de queue effective à construire avec approximation
contrôlée du noyau et contraintes.

\textbf{Ancrage.} Documents 67-68 et noyau du corpus ; définition locale
explicitée dans cette fiche ; définition pédagogique ; voir portée
précise dans ancrage et garde.

\hypertarget{minimum-de-schur}{%
\subsection{Minimum de Schur}\label{minimum-de-schur}}

\textbf{Comme un premier jeu.} On règle les anciens jouets pour soulager
le nouveau.

\textbf{Symbole.} inf\_z Q(z+v)=δ

\textbf{Sens et domaine.} Valable si ancien bloc \textgreater0 ; formule
minimisation avec z=-G\^{}(-1)c.

\textbf{Garde.} Si le bloc est indéfini, ce minimum vaut -∞ même si la
formule algébrique existe.

\textbf{Ancrage.} Documents 67-68 et noyau du corpus ; définition locale
explicitée dans cette fiche ; définition pédagogique ; voir portée
précise dans ancrage et garde.

\hypertarget{preuve-de-phase}{%
\subsection{Preuve de phase}\label{preuve-de-phase}}

\textbf{Comme un premier jeu.} Deux rotations identiques se compensent.

\textbf{Symbole.} h\_(e\^{}(iθ)f)=h\_f

\textbf{Sens et domaine.} Facteurs e\^{}(iθ) et conj e\^{}(iθ)
s'annulent dans corrélation et normes.

\textbf{Garde.} La compensation n'a pas lieu pour une phase variable aux
deux positions.

\textbf{Ancrage.} Documents 67-68 et noyau du corpus ; définition locale
explicitée dans cette fiche ; définition pédagogique ; voir portée
précise dans ancrage et garde.

\hypertarget{une-queue-existe}{%
\subsection{Une queue existe}\label{une-queue-existe}}

\textbf{Comme un premier jeu.} Les pièces oubliées deviennent petites
après la règle douce.

\textbf{Symbole.} η\_N→0

\textbf{Sens et domaine.} Compactité j\_S et exhaustion R-orthogonale
dense donnent convergence en norme.

\textbf{Garde.} La vitesse et les indices de réussite doivent encore
être calculés.

\textbf{Ancrage.} Documents 67-68 et noyau du corpus ; définition locale
explicitée dans cette fiche ; définition pédagogique ; voir portée
précise dans ancrage et garde.

\hypertarget{r-notation-qualifiuxe9e}{%
\subsection{R notation qualifiée}\label{r-notation-qualifiuxe9e}}

\textbf{Comme un premier jeu.} Deux boîtes séparent une lettre partagée.

\textbf{Symbole.} R\_metric, R\_reflection

\textbf{Sens et domaine.} La métrique logarithmique et l'opérateur
miroir ont des rôles différents.

\textbf{Garde.} L'identifiant stable empêche une composition ambiguë.

\textbf{Ancrage.} Documents 67-68 et noyau du corpus ; définition locale
explicitée dans cette fiche ; définition pédagogique ; voir portée
précise dans ancrage et garde.

\hypertarget{croissance-logarithmique}{%
\subsection{Croissance logarithmique}\label{croissance-logarithmique}}

\textbf{Comme un premier jeu.} Les rythmes rapides prennent un poids qui
grandit doucement.

\textbf{Symbole.} a(τ) comparable à log(2+\textbar τ\textbar)

\textbf{Sens et domaine.} La série positive donne divergence ;
comparaisons élémentaires donnent ordre logarithmique.

\textbf{Garde.} Les constantes d'une estimation effective doivent être
écrites, pas remplacées par le mot comparable.

\textbf{Ancrage.} Documents 67-68 et noyau du corpus ; définition locale
explicitée dans cette fiche ; définition pédagogique ; voir portée
précise dans ancrage et garde.

\hypertarget{budget-de-contraction}{%
\subsection{Budget de contraction}\label{budget-de-contraction}}

\textbf{Comme un premier jeu.} Rétrécir peut rendre les rythmes plus
rapides, mais leur coût est borné.

\textbf{Symbole.} a(τ/λ)≤λ\^{}(-2)a(τ)

\textbf{Sens et domaine.} Termes positifs de la série de a donnent
l'inégalité pour 0\textless λ≤1 ; continuité forte à1.

\textbf{Garde.} Ce budget ne préserve pas Q\_W ni les moments.

\textbf{Ancrage.} Documents 67-68 et noyau du corpus ; définition locale
explicitée dans cette fiche ; définition pédagogique ; voir portée
précise dans ancrage et garde.

\hypertarget{transport-dopuxe9rateur}{%
\subsection{Transport d'opérateur}\label{transport-dopuxe9rateur}}

\textbf{Comme un premier jeu.} La même machine change de pièce avec sa
règle.

\textbf{Symbole.} A↦UAU\^{}(-1) si U unitaire

\textbf{Sens et domaine.} Spectre et propriétés unitaires conservés pour
domaines correctement transportés.

\textbf{Garde.} Changer la métrique de Riesz H¹→R n'est pas
automatiquement cette opération.

\textbf{Ancrage.} Documents 67-68 et noyau du corpus ; définition locale
explicitée dans cette fiche ; définition pédagogique ; voir portée
précise dans ancrage et garde.

\hypertarget{fonction-logarithmique-gs}{%
\subsection{Fonction logarithmique
G(s)}\label{fonction-logarithmique-gs}}

\textbf{Comme un premier jeu.} Un carnet mesure la rotation d'une
fonction.

\textbf{Symbole.} G(s)=Log F(s) selon archive

\textbf{Sens et domaine.} Définition dépendant d'un domaine sans zéros
et d'une branche ; voir source historique locale.

\textbf{Garde.} Ne pas l'identifier au Gram ni la considérer globalement
définie à travers tous les zéros.

\textbf{Ancrage.} Documents 67-68 et noyau du corpus ; définition locale
explicitée dans cette fiche ; définition pédagogique ; voir portée
précise dans ancrage et garde.

\hypertarget{surface-virtuelle}{%
\subsection{Surface virtuelle}\label{surface-virtuelle}}

\textbf{Comme un premier jeu.} Un dessin ressemble à un terrain avant
d'en être un.

\textbf{Symbole.} axiomes formels du corpus

\textbf{Sens et domaine.} Construction comparative/virtuelle des
documents46-48, dont domaine et axiomes doivent être repris.

\textbf{Garde.} Aucune variété géométrique de ζ construite par cette
seule analogie.

\textbf{Ancrage.} Documents 67-68 et noyau du corpus ; définition locale
explicitée dans cette fiche ; définition pédagogique ; voir portée
précise dans ancrage et garde.

\hypertarget{zuxe9ro-critique-double}{%
\subsection{Zéro critique double}\label{zuxe9ro-critique-double}}

\textbf{Comme un premier jeu.} La note touche la ligne deux fois.

\textbf{Symbole.} F(ρ)=F′(ρ)=0, Re ρ=1/2

\textbf{Sens et domaine.} Question de multiplicité/collision avec jets
certifiés locaux selon la famille.

\textbf{Garde.} RH n'impose pas la simplicité des zéros ; ne pas
confondre ces deux buts.

\textbf{Ancrage.} Documents 67-68 et noyau du corpus ; définition locale
explicitée dans cette fiche ; définition pédagogique ; voir portée
précise dans ancrage et garde.

\hypertarget{logique-des-domaines}{%
\subsection{Logique des domaines}\label{logique-des-domaines}}

\textbf{Comme un premier jeu.} La phrase dit quelle boîte le dessin peut
atteindre.

\textbf{Symbole.} Γ⊢f : A(S,r,u)

\textbf{Sens et domaine.} S support, r ordre garanti de jets, u degré
d'unité A ; règles sound dans A, preuves L8.

\textbf{Garde.} Pas un système complet décidant toute identité ou la
vérité de RH.

\textbf{Ancrage.}
3\_Preuves\_et\_calculs/preuves/V644\_LANGUE\_ET\_DOMAINES.md ; L6-L8 ;
relation démontrée dans V644, L6-L8.

\hypertarget{queue-de-moments-ordinaires}{%
\subsection{Queue de moments
ordinaires}\label{queue-de-moments-ordinaires}}

\textbf{Comme un premier jeu.} Les rythmes lents deviennent minuscules
après plusieurs balances.

\textbf{Symbole.} Y\_N=\{∫x\^{}j f=0, 0≤j≤N\}

\textbf{Sens et domaine.} Codimension N+1 dans X\_S ; queue positive
effective via Taylor Fourier et budget L².

\textbf{Garde.} Distincte des idéaux I\_r, qui utilisent les poids
e\^{}(±x/2).

\textbf{Ancrage.}
3\_Preuves\_et\_calculs/preuves/V644\_LANGUE\_ET\_DOMAINES.md ; L6-L8 ;
relation démontrée dans V644, L6-L8.

\hypertarget{borne-explicite-de-queue}{%
\subsection{Borne explicite de queue}\label{borne-explicite-de-queue}}

\textbf{Comme un premier jeu.} On encadre les rythmes lents dans une
petite boîte.

\textbf{Symbole.} ε\_N(T)=2(ST)\^{}(2N+3)/{[}π(2N+3)²((N+1)!)²{]}

\textbf{Sens et domaine.} Pour Y\_N, énergie Fourier dans
{[}-T,T{]}≤ε\_N\textbar\textbar f\textbar\textbar₂² ; si ε\textless1,
\textbar\textbar f\textbar\textbar₂²≤R/{[}(1-ε)r(T){]}.

\textbf{Garde.} Borne prouvée, mais choisir des paramètres grossiers
peut donner une tête astronomique.

\textbf{Ancrage.}
3\_Preuves\_et\_calculs/preuves/V644\_LANGUE\_ET\_DOMAINES.md ; L6-L8 ;
relation démontrée dans V644, L6-L8.

\hypertarget{contrat-de-queue-entier}{%
\subsection{Contrat de queue entier}\label{contrat-de-queue-entier}}

\textbf{Comme un premier jeu.} Une recette ferme la queue avec des
comptes entiers.

\textbf{Symbole.} Lbar=8+2M, T=2\^{}(8Lbar), q=ceil(6ST)

\textbf{Sens et domaine.} Pour Y\_(q-1), Q\_W≥R/(1+2Lbar) ; M nombre
exact de puissances premières actives.

\textbf{Garde.} La tête a dimension q ; elle n'a pas été calculée dans
cette expérience.

\textbf{Ancrage.}
3\_Preuves\_et\_calculs/preuves/V644\_LANGUE\_ET\_DOMAINES.md ; L6-L8 ;
relation démontrée dans V644, L6-L8.

\hypertarget{erreur-de-dette-par-neumann}{%
\subsection{Erreur de dette par
Neumann}\label{erreur-de-dette-par-neumann}}

\textbf{Comme un premier jeu.} Les corrections restantes ont un ticket
de taille.

\textbf{Symbole.} reste≤η\^{}(m+1)/(1-η)

\textbf{Sens et domaine.} η=L\_S/{[}(1-ε)r(T){]}\textless1 borne le
compact comprimé sur la queue ; budget inverse et Schur prouvés.

\textbf{Garde.} Le solveur des actions/projections doit encore être
certifié.

\textbf{Ancrage.}
3\_Preuves\_et\_calculs/preuves/V644\_LANGUE\_ET\_DOMAINES.md ; L6-L8 ;
relation démontrée dans V644, L6-L8.

\hypertarget{repruxe9sentant-riesz-des-moments}{%
\subsection{Représentant Riesz des
moments}\label{repruxe9sentant-riesz-des-moments}}

\textbf{Comme un premier jeu.} Un capteur devient un dessin dans la
bonne règle.

\textbf{Symbole.} ⟨h\_j,f⟩R=∫x\^{}j f(x)dx

\textbf{Sens et domaine.} Les h\_j engendrent F\_N=Y\_N⊥\_R ;
codimension N+1 ; solveur et résidus à certifier.

\textbf{Garde.} Cette tête diffère du préfixe polynomial brut utilisé
par les matrices exploratoires.

\textbf{Ancrage.}
3\_Preuves\_et\_calculs/preuves/V644\_LANGUE\_ET\_DOMAINES.md ; L6-L8 ;
relation démontrée dans V644, L6-L8.

\hypertarget{annexe---tous-les-noms-du-registre-canonique}{%
\section{Annexe - Tous les noms du registre
canonique}\label{annexe---tous-les-noms-du-registre-canonique}}

Les 95 entrées suivantes conservent les dernières lignes nommées du
registre 20. Les images historiques peuvent rester plus techniques que
celles du noyau. Leur absence éventuelle de statut explicite est
signalée.

\hypertarget{miroir-de-muxf6bius}{%
\subsection{Miroir de Möbius}\label{miroir-de-muxf6bius}}

\textbf{Image héritée.} zéros et cancellation sont deux faces d'un même
paysage

\textbf{Ancrage exact.} \(1/\zeta(s)=\sum \mu(n)n^{-s}\), sommes de
Möbius

\textbf{Garde héritée.} assimiler \(\mu\) à des tirages indépendants
exacts

\textbf{Correspondance autorisée.} traduire localisation spectrale en
fluctuation arithmétique

\textbf{Provenance.} 1\_Presentation/20\_Registre\_metaphores\_V6.md:11
; 1. Tableau canonique. \textbf{Statut :} hérité du registre : MÉCANISME
RETENU.

\hypertarget{filtre-antipuxf4le}{%
\subsection{Filtre antipôle}\label{filtre-antipuxf4le}}

\textbf{Image héritée.} filtre notch supprimant un mode cohérent

\textbf{Ancrage exact.} \(W_N(1)=0\)

\textbf{Garde héritée.} présenter l'annulation du pôle comme un
mécanisme suffisant pour RH

\textbf{Correspondance autorisée.} isoler un mode universel avant
d'étudier le résidu

\textbf{Provenance.} 1\_Presentation/20\_Registre\_metaphores\_V6.md:12
; 1. Tableau canonique. \textbf{Statut :} hérité du registre : BRANCHE
CLOSE / OUTIL RETENU.

\hypertarget{thermomuxe8tre-dhorizon}{%
\subsection{Thermomètre d'horizon}\label{thermomuxe8tre-dhorizon}}

\textbf{Image héritée.} thermomètre qui mesure sans produire la
température

\textbf{Ancrage exact.} énergie \(Q\), exposant relié à
\(\Theta=\sup\Re\rho\)

\textbf{Garde héritée.} demander à une norme positive de fabriquer la
positivité de Weil

\textbf{Correspondance autorisée.} utiliser \(Q\) comme diagnostic de
croissance

\textbf{Provenance.} 1\_Presentation/20\_Registre\_metaphores\_V6.md:13
; 1. Tableau canonique. \textbf{Statut :} hérité du registre :
DIAGNOSTIC EXACT.

\hypertarget{effaceur-orthogonal}{%
\subsection{Effaceur orthogonal}\label{effaceur-orthogonal}}

\textbf{Image héritée.} filtre adaptatif retirant une composante
structurée

\textbf{Ancrage exact.} projection Nyman-Beurling / Gram

\textbf{Garde héritée.} confondre forte corrélation finie et convergence
asymptotique

\textbf{Correspondance autorisée.} quantifier l'énergie capturée par une
direction

\textbf{Provenance.} 1\_Presentation/20\_Registre\_metaphores\_V6.md:14
; 1. Tableau canonique. \textbf{Statut :} hérité du registre :
DIAGNOSTIC.

\hypertarget{positivituxe9-porteuse}{%
\subsection{Positivité porteuse}\label{positivituxe9-porteuse}}

\textbf{Image héritée.} arche dont le signe porte une contrainte

\textbf{Ancrage exact.} forme hermitienne de Weil, signature, inertie

\textbf{Garde héritée.} appeler « information » toute quantité \(\ge0\)
par définition

\textbf{Correspondance autorisée.} distinguer une positivité
arithmétique d'une somme de carrés automatique

\textbf{Provenance.} 1\_Presentation/20\_Registre\_metaphores\_V6.md:15
; 1. Tableau canonique. \textbf{Statut :} hérité du registre : CENTRAL.

\hypertarget{census-spectral}{%
\subsection{Census spectral}\label{census-spectral}}

\textbf{Image héritée.} recenser plutôt que localiser chaque individu

\textbf{Ancrage exact.} rang, trace, inertie, proportions

\textbf{Garde héritée.} transformer une proportion en RH complète

\textbf{Correspondance autorisée.} obtenir des résultats proportionnels
sans localisation totale

\textbf{Provenance.} 1\_Presentation/20\_Registre\_metaphores\_V6.md:16
; 1. Tableau canonique. \textbf{Statut :} hérité du registre : MÉTHODE
CENTRALE.

\hypertarget{microscope-de-population}{%
\subsection{Microscope de population}\label{microscope-de-population}}

\textbf{Image héritée.} épidémiologie plutôt qu'autopsie individuelle

\textbf{Ancrage exact.} familles de fonctions \(L\), moyennes

\textbf{Garde héritée.} attribuer automatiquement une propriété moyenne
à \(\zeta\) seule

\textbf{Correspondance autorisée.} faire émerger des symétries
collectives

\textbf{Provenance.} 1\_Presentation/20\_Registre\_metaphores\_V6.md:17
; 1. Tableau canonique. \textbf{Statut :} hérité du registre : PARADIGME
EXTERNE.

\hypertarget{thermostat-critique}{%
\subsection{Thermostat critique}\label{thermostat-critique}}

\textbf{Image héritée.} seuil de transition

\textbf{Ancrage exact.} flot de de Bruijn-Newman, \(\Lambda\)

\textbf{Garde héritée.} interpréter \(\Lambda\) comme température
physique

\textbf{Correspondance autorisée.} reformuler RH comme criticité exacte
\(\Lambda=0\)

\textbf{Provenance.} 1\_Presentation/20\_Registre\_metaphores\_V6.md:18
; 1. Tableau canonique. \textbf{Statut :} hérité du registre : RÉSERVE
CONCEPTUELLE.

\hypertarget{scuxe8ne-manquante}{%
\subsection{Scène manquante}\label{scuxe8ne-manquante}}

\textbf{Image héritée.} chorégraphie visible, théâtre absent

\textbf{Ancrage exact.} géométrie/cohomologie/trace

\textbf{Garde héritée.} postuler une géométrie parce qu'une formule lui
ressemble

\textbf{Correspondance autorisée.} chercher une structure expliquant
trace, dualité, positivité

\textbf{Provenance.} 1\_Presentation/20\_Registre\_metaphores\_V6.md:19
; 1. Tableau canonique. \textbf{Statut :} hérité du registre : OUVERT.

\hypertarget{jumeau-de-crash-test}{%
\subsection{Jumeau de crash-test}\label{jumeau-de-crash-test}}

\textbf{Image héritée.} même carrosserie analytique, moteur différent

\textbf{Ancrage exact.} Davenport-Heilbronn et analogues

\textbf{Garde héritée.} traiter DH comme un modèle complet de \(\zeta\)

\textbf{Correspondance autorisée.} rejeter les critères ne voyant que la
symétrie fonctionnelle

\textbf{Provenance.} 1\_Presentation/20\_Registre\_metaphores\_V6.md:20
; 1. Tableau canonique. \textbf{Statut :} hérité du registre :
INSTRUMENT ACTIF.

\hypertarget{modes-sombres}{%
\subsection{Modes sombres}\label{modes-sombres}}

\textbf{Image héritée.} mode réel que le capteur ne voit pas

\textbf{Ancrage exact.} observabilité des opérateurs / complétude

\textbf{Garde héritée.} baptiser « sombre » toute inconnue non comprise

\textbf{Correspondance autorisée.} formuler explicitement les pertes
d'observabilité

\textbf{Provenance.} 1\_Presentation/20\_Registre\_metaphores\_V6.md:21
; 1. Tableau canonique. \textbf{Statut :} hérité du registre : CADRE
D'OBSERVABILITÉ.

\hypertarget{compression-fertile}{%
\subsection{Compression fertile}\label{compression-fertile}}

\textbf{Image héritée.} biopsie informative

\textbf{Ancrage exact.} compressions finies d'une forme/opérateur infini

\textbf{Garde héritée.} croire qu'une discrétisation finie prouve à elle
seule une propriété asymptotique

\textbf{Correspondance autorisée.} transférer rang, signe, inertie vers
un calcul fini

\textbf{Provenance.} 1\_Presentation/20\_Registre\_metaphores\_V6.md:22
; 1. Tableau canonique. \textbf{Statut :} hérité du registre : ACTIF.

\hypertarget{secteur-sombre-arithmuxe9tique-dark-zeta}{%
\subsection{Secteur sombre arithmétique /
DARK-ZETA}\label{secteur-sombre-arithmuxe9tique-dark-zeta}}

\textbf{Image héritée.} cartographier l'invisible par ses déformations
du visible

\textbf{Ancrage exact.} information perdue par norme, moyenne,
covariance

\textbf{Garde héritée.} introduire une entité ou « énergie noire »
physique

\textbf{Correspondance autorisée.} trianguler plusieurs empreintes
indépendantes

\textbf{Provenance.} 1\_Presentation/20\_Registre\_metaphores\_V6.md:23
; 1. Tableau canonique. \textbf{Statut :} hérité du registre : CADRE
TRANSVERSAL.

\hypertarget{tomographie-de-cohuxe9rence-tomo-2}{%
\subsection{Tomographie de cohérence
TOMO-2}\label{tomographie-de-cohuxe9rence-tomo-2}}

\textbf{Image héritée.} scanner inter-échelles

\textbf{Ancrage exact.} Gram des translations logarithmiques

\textbf{Garde héritée.} prétendre voir la phase Fourier

\textbf{Correspondance autorisée.} mesurer le spectre de
puissance/covariance

\textbf{Provenance.} 1\_Presentation/20\_Registre\_metaphores\_V6.md:24
; 1. Tableau canonique. \textbf{Statut :} hérité du registre :
DIAGNOSTIC SURVIVANT.

\hypertarget{tomographie-de-phase-tomo-3}{%
\subsection{Tomographie de phase
TOMO-3}\label{tomographie-de-phase-tomo-3}}

\textbf{Image héritée.} scanner de fermeture de phase

\textbf{Ancrage exact.} tricorrélation / bispectre

\textbf{Garde héritée.} identifier automatiquement bicoherence et
multiplicativité

\textbf{Correspondance autorisée.} tester des couplages ternaires

\textbf{Provenance.} 1\_Presentation/20\_Registre\_metaphores\_V6.md:25
; 1. Tableau canonique. \textbf{Statut :} hérité du registre :
RÉTROGRADÉE EN BRUT.

\hypertarget{fermeture-multiplicative}{%
\subsection{Fermeture multiplicative}\label{fermeture-multiplicative}}

\textbf{Image héritée.} une relation locale qui se referme sur elle-même

\textbf{Ancrage exact.} \(\Delta_F(m,n)=a(mn)-a(m)a(n)\), lois sur
\(p^k\)

\textbf{Garde héritée.} confondre multiplicativité et RH

\textbf{Correspondance autorisée.} mesurer la cohérence eulérienne

\textbf{Provenance.} 1\_Presentation/20\_Registre\_metaphores\_V6.md:26
; 1. Tableau canonique. \textbf{Statut :} hérité du registre :
FORMALISÉE.

\hypertarget{flux-dinertie}{%
\subsection{Flux d'inertie}\label{flux-dinertie}}

\textbf{Image héritée.} suivre les directions de signe sous déformation

\textbf{Ancrage exact.} \((n_+,n_-,n_0)\) d'une famille hermitienne

\textbf{Garde héritée.} interpréter un saut fini comme la criticité
asymptotique sans limite contrôlée

\textbf{Correspondance autorisée.} remplacer le suivi de chaque zéro par
un invariant de signature

\textbf{Provenance.} 1\_Presentation/20\_Registre\_metaphores\_V6.md:27
; 1. Tableau canonique. \textbf{Statut :} hérité du registre : RÉSERVE.

\hypertarget{guxe9omuxe9trie-par-empreintes}{%
\subsection{Géométrie par
empreintes}\label{guxe9omuxe9trie-par-empreintes}}

\textbf{Image héritée.} reconstruire l'animal par ses traces

\textbf{Ancrage exact.} signatures/indices/dualités distinguant
\(\zeta\) de faux jumeaux

\textbf{Garde héritée.} prendre une liste d'analogies pour une géométrie
existante

\textbf{Correspondance autorisée.} inférer des axiomes nécessaires avant
de construire la géométrie

\textbf{Provenance.} 1\_Presentation/20\_Registre\_metaphores\_V6.md:28
; 1. Tableau canonique. \textbf{Statut :} hérité du registre : ACTIF
CONCEPTUEL.

\hypertarget{census-de-muxf6bius}{%
\subsection{Census de Möbius}\label{census-de-muxf6bius}}

\textbf{Image héritée.} recenser la structure d'une population de blocs

\textbf{Ancrage exact.} covariances/rangs de blocs de \(\mu\)

\textbf{Garde héritée.} déduire les zéros individuels d'une covariance
d'ordre 2

\textbf{Correspondance autorisée.} étudier où la cancellation apparaît
statistiquement

\textbf{Provenance.} 1\_Presentation/20\_Registre\_metaphores\_V6.md:29
; 1. Tableau canonique. \textbf{Statut :} hérité du registre :
SECONDAIRE.

\hypertarget{prime-lock-verrou-premier}{%
\subsection{PRIME-LOCK / verrou
premier}\label{prime-lock-verrou-premier}}

\textbf{Image héritée.} serrure locale du produit d'Euler

\textbf{Ancrage exact.} \(-F'/F\), support sur puissances de premiers,
défaut de tour

\textbf{Garde héritée.} appeler « RH » une propriété satisfaite aussi
par Steinhaus

\textbf{Correspondance autorisée.} détecter la structure eulérienne
locale

\textbf{Provenance.} 1\_Presentation/20\_Registre\_metaphores\_V6.md:30
; 1. Tableau canonique. \textbf{Statut :} hérité du registre : TESTÉ /
VALIDÉ COMME DÉTECTEUR LOCAL.

\hypertarget{dual-lock-serrure-uxe0-deux-cluxe9s}{%
\subsection{DUAL-LOCK / serrure à deux
clés}\label{dual-lock-serrure-uxe0-deux-cluxe9s}}

\textbf{Image héritée.} une serrure ne s'ouvre que si local et global
sont compatibles

\textbf{Ancrage exact.} compression de la forme explicite : archimédien
+ conducteur/pôle + local

\textbf{Garde héritée.} juxtaposer deux scores indépendants et appeler
cela interaction

\textbf{Correspondance autorisée.} tester intrinsèquement la
compatibilité entre données aux premiers et complétion globale

\textbf{Provenance.} 1\_Presentation/20\_Registre\_metaphores\_V6.md:31
; 1. Tableau canonique. \textbf{Statut :} hérité du registre : TESTÉ /
SURVIT COMME DÉTECTEUR D'INERTIE.

\hypertarget{ruxe9sidu-sombre}{%
\subsection{Résidu sombre}\label{ruxe9sidu-sombre}}

\textbf{Image héritée.} information visible seulement après avoir retiré
presque tout le signal dominant

\textbf{Ancrage exact.} petite valeur finale après cancellation des
contributions local/global

\textbf{Garde héritée.} assimiler le résidu à une substance ou conclure
qu'un petit nombre est automatiquement profond

\textbf{Correspondance autorisée.} quantifier où signe/inertie sont
portés par un résidu de cancellation

\textbf{Provenance.} 1\_Presentation/20\_Registre\_metaphores\_V6.md:32
; 1. Tableau canonique. \textbf{Statut :} hérité du registre : NOUVEAU -
TESTÉ V6.3.

\hypertarget{le-thermomuxe8tre}{%
\subsection{Le Thermomètre}\label{le-thermomuxe8tre}}

\textbf{Image héritée.} Q(X) prend la « fièvre transverse » : sa vitesse
de croissance dit de combien le zéro le plus éloigné s'écarte de la
droite.

\textbf{Ancrage exact.} Théorème A : ½ limsup log(1+Q)/log X = Θ - ½ ;
RH ⇔ Q = X\^{}o(1). (LEMMES, L6)

\textbf{Garde héritée.} Lire la fièvre ne la fait pas baisser : c'est
une reformulation de RH.

\textbf{Provenance.} 1\_Presentation/20\_Registre\_metaphores\_V6.md:98
; Annexe B - Lexique des résultats qui ont porté (fusion de la v6.0,
Claude). \textbf{Statut :} hérité du registre : statut non explicité
dans cette ligne historique.

\hypertarget{la-loupe-double}{%
\subsection{La Loupe double}\label{la-loupe-double}}

\textbf{Image héritée.} c est fabriqué pour que chaque zéro devienne un
pôle double : chaque zéro est grossi d'un facteur log.

\textbf{Ancrage exact.} G = m(m+1)/(s-ρ)² + \ldots{} (Lemme 5) ; Q ∼
(S/3)log³X sous RH.

\textbf{Garde héritée.} Une loupe grossit aussi les défauts : les
oscillations du terme suivant (les Battements).

\textbf{Provenance.} 1\_Presentation/20\_Registre\_metaphores\_V6.md:99
; Annexe B - Lexique des résultats qui ont porté (fusion de la v6.0,
Claude). \textbf{Statut :} hérité du registre : statut non explicité
dans cette ligne historique.

\hypertarget{le-contrepoids}{%
\subsection{Le Contrepoids}\label{le-contrepoids}}

\textbf{Image héritée.} Le terme 2γΛ équilibre exactement le pôle de ζ
en s = 1.

\textbf{Ancrage exact.} G est régulière en 1, G(1) = 2γ² + 6γ₁ (Lemme
4).

\textbf{Garde héritée.} L'équilibre ne tient qu'en un point.

\textbf{Provenance.} 1\_Presentation/20\_Registre\_metaphores\_V6.md:100
; Annexe B - Lexique des résultats qui ont porté (fusion de la v6.0,
Claude). \textbf{Statut :} hérité du registre : statut non explicité
dans cette ligne historique.

\hypertarget{le-tamis-uxe0-deux}{%
\subsection{Le Tamis à deux}\label{le-tamis-uxe0-deux}}

\textbf{Image héritée.} c ne retient que les entiers qui ont au plus
deux facteurs premiers distincts.

\textbf{Ancrage exact.} Lemme 3 : c(n) = 0 dès trois facteurs premiers
distincts.

\textbf{Garde héritée.} Aucun argument n'a encore utilisé cette
parcimonie.

\textbf{Provenance.} 1\_Presentation/20\_Registre\_metaphores\_V6.md:101
; Annexe B - Lexique des résultats qui ont porté (fusion de la v6.0,
Claude). \textbf{Statut :} hérité du registre : statut non explicité
dans cette ligne historique.

\hypertarget{la-note-fondamentale}{%
\subsection{La Note fondamentale}\label{la-note-fondamentale}}

\textbf{Image héritée.} κ = 2 + γ - log 4π est la note que donnent tous
les zéros ensemble ; elle revient partout.

\textbf{Ancrage exact.} κ = Σ 1/(ρ(1-ρ)) = 2λ₁ (Li) ; dans Q (Théorème
B), dans d\_N² (Burnol, BCF).

\textbf{Garde héritée.} C'est une constante, pas un mécanisme ; Σ 1/∣ρ∣²
= κ seulement sous RH.

\textbf{Provenance.} 1\_Presentation/20\_Registre\_metaphores\_V6.md:107
; Annexe B - Lexique des résultats qui ont porté (fusion de la v6.0,
Claude). \textbf{Statut :} hérité du registre : statut non explicité
dans cette ligne historique.

\hypertarget{le-diapason}{%
\subsection{Le Diapason}\label{le-diapason}}

\textbf{Image héritée.} Le son de référence quand tout est accordé.

\textbf{Ancrage exact.} Théorème B : sous RH, Q ∼ (4κ/3) log³X (zéros
simples).

\textbf{Garde héritée.} Le son pur n'apparaît qu'à l'infini : à X ≤
4·10\^{}6, Q n'est qu'à 44-72 \% du diapason.

\textbf{Provenance.} 1\_Presentation/20\_Registre\_metaphores\_V6.md:108
; Annexe B - Lexique des résultats qui ont porté (fusion de la v6.0,
Claude). \textbf{Statut :} hérité du registre : statut non explicité
dans cette ligne historique.

\hypertarget{les-battements}{%
\subsection{Les Battements}\label{les-battements}}

\textbf{Image héritée.} Deux cordes presque accordées battent ; deux
zéros voisins font battre le terme en L².

\textbf{Ancrage exact.} Proposition C : coefficient α + Φ(L), Φ = somme
de termes e\^{}i(γ-γ′)L/(γ-γ′) ; corrélation 0,996.

\textbf{Garde héritée.} Les paires trop proches ne sont pas résolues et
brouillent le signal.

\textbf{Provenance.} 1\_Presentation/20\_Registre\_metaphores\_V6.md:109
; Annexe B - Lexique des résultats qui ont porté (fusion de la v6.0,
Claude). \textbf{Statut :} hérité du registre : statut non explicité
dans cette ligne historique.

\hypertarget{le-carrousel}{%
\subsection{Le Carrousel}\label{le-carrousel}}

\textbf{Image héritée.} Chaque zéro sur la droite est une nacelle qui
tourne sans monter ; un zéro hors de la droite serait une nacelle qui
s'élève.

\textbf{Ancrage exact.} Les modes e\^{}iγu sont les vecteurs propres de
la translation ; sous RH, orthogonaux en moyenne (doc. 14).

\textbf{Garde héritée.} La rotation pure est équivalente à RH : on ne
voit pas le moteur qui l'impose.

\textbf{Provenance.} 1\_Presentation/20\_Registre\_metaphores\_V6.md:110
; Annexe B - Lexique des résultats qui ont porté (fusion de la v6.0,
Claude). \textbf{Statut :} hérité du registre : statut non explicité
dans cette ligne historique.

\hypertarget{la-gare-centrale}{%
\subsection{La Gare centrale}\label{la-gare-centrale}}

\textbf{Image héritée.} Toutes les lignes essayées arrivent à Q.

\textbf{Ancrage exact.} (★) ⇔ Q ≪ log³X ; e² = β²Q + T\_N ; RH ⇔ Q =
X\^{}o(1).

\textbf{Garde héritée.} Ce n'est pas une destination : c'est le
carrefour où RH se cache.

\textbf{Provenance.} 1\_Presentation/20\_Registre\_metaphores\_V6.md:116
; Annexe B - Lexique des résultats qui ont porté (fusion de la v6.0,
Claude). \textbf{Statut :} hérité du registre : statut non explicité
dans cette ligne historique.

\hypertarget{le-bouchon-et-la-fuite}{%
\subsection{Le Bouchon et la Fuite}\label{le-bouchon-et-la-fuite}}

\textbf{Image héritée.} Boucher le pôle en s = 1 ; ce qui fuit encore
est exactement Q, plus un filet (la queue).

\textbf{Ancrage exact.} β = -S\_N/A\_N ⇒ W\_N(1) = 0 ; e² = β²Q(N) +
T\_N ; b\_N → -1 - 2γ/log N (doc. 12).

\textbf{Garde héritée.} Boucher n'arrête pas la fuite principale.

\textbf{Provenance.} 1\_Presentation/20\_Registre\_metaphores\_V6.md:117
; Annexe B - Lexique des résultats qui ont porté (fusion de la v6.0,
Claude). \textbf{Statut :} hérité du registre : statut non explicité
dans cette ligne historique.

\hypertarget{le-palier}{%
\subsection{Le Palier}\label{le-palier}}

\textbf{Image héritée.} La famille à deux directions s'arrête sur un
palier, un étage au-dessus du sol.

\textbf{Ancrage exact.} \{μ, a\} ≈ 1,3 κ/log N de N = 50 à 10\^{}5 (doc.
10, 12).

\textbf{Garde héritée.} Sous RH le palier devrait monter vers 4/3 : ce
n'est pas un sol.

\textbf{Provenance.} 1\_Presentation/20\_Registre\_metaphores\_V6.md:118
; Annexe B - Lexique des résultats qui ont porté (fusion de la v6.0,
Claude). \textbf{Statut :} hérité du registre : statut non explicité
dans cette ligne historique.

\hypertarget{le-dernier-milliuxe8me}{%
\subsection{Le Dernier millième}\label{le-dernier-milliuxe8me}}

\textbf{Image héritée.} Quatre coups de pinceau font 90 \% du tableau ;
les derniers détails en demandent toujours plus.

\textbf{Ancrage exact.} 90 \% en 3-4 directions propres ; 99,9 \% en 52,
81, 121 directions (N = 200, 500, 998) (doc. 14).

\textbf{Garde héritée.} N ≤ 998 : la loi de croissance n'est
qu'observée.

\textbf{Provenance.} 1\_Presentation/20\_Registre\_metaphores\_V6.md:119
; Annexe B - Lexique des résultats qui ont porté (fusion de la v6.0,
Claude). \textbf{Statut :} hérité du registre : statut non explicité
dans cette ligne historique.

\hypertarget{luxe9cho-des-pgcd}{%
\subsection{L'Écho des pgcd}\label{luxe9cho-des-pgcd}}

\textbf{Image héritée.} Deux dents de scie de périodes a et b se
répondent avec l'intensité pgcd²/(12ab).

\textbf{Ancrage exact.} Franel-Landau (L8) ; V\_N = (1/12) Σ
J₂(q)(Σ\_q∣d w\_d/d)².

\textbf{Garde héritée.} L'écho est une moyenne : il gouverne la queue,
pas les zéros.

\textbf{Provenance.} 1\_Presentation/20\_Registre\_metaphores\_V6.md:120
; Annexe B - Lexique des résultats qui ont porté (fusion de la v6.0,
Claude). \textbf{Statut :} hérité du registre : statut non explicité
dans cette ligne historique.

\hypertarget{le-sosie}{%
\subsection{Le Sosie}\label{le-sosie}}

\textbf{Image héritée.} Même visage que ζ (même équation fonctionnelle)
mais sans produit d'Euler, et avec des zéros hors de la droite.

\textbf{Ancrage exact.} Davenport-Heilbronn : zéro 0,8085 + 85,699i ;
pente de Q 0,56 contre 0,33 (doc. DH).

\textbf{Garde héritée.} -

\textbf{Provenance.} 1\_Presentation/20\_Registre\_metaphores\_V6.md:126
; Annexe B - Lexique des résultats qui ont porté (fusion de la v6.0,
Claude). \textbf{Statut :} hérité du registre : statut non explicité
dans cette ligne historique.

\hypertarget{le-moduxe8le-ruxe9duit}{%
\subsection{Le Modèle réduit}\label{le-moduxe8le-ruxe9duit}}

\textbf{Image héritée.} Une maquette où RH est vraie, grâce à une
géométrie que l'original n'a pas.

\textbf{Ancrage exact.} Courbes sur F\_q (Weil) ; test F\_101 : max ∣a∣
= 20 ≤ 2√101 (mémo P2.2).

\textbf{Garde héritée.} La maquette a une dimension de plus : Spec ℤ ×
Spec ℤ = Spec ℤ.

\textbf{Provenance.} 1\_Presentation/20\_Registre\_metaphores\_V6.md:127
; Annexe B - Lexique des résultats qui ont porté (fusion de la v6.0,
Claude). \textbf{Statut :} hérité du registre : statut non explicité
dans cette ligne historique.

\hypertarget{le-duxe9tecteur}{%
\subsection{Le Détecteur}\label{le-duxe9tecteur}}

\textbf{Image héritée.} K(L) monte doucement pour ζ et s'emballe pour le
Sosie.

\textbf{Ancrage exact.} K(15) = 0,098 pour ζ ; K(12) = 0,885 pour DH
(doc. 09).

\textbf{Garde héritée.} Dépistage, pas preuve.

\textbf{Provenance.} 1\_Presentation/20\_Registre\_metaphores\_V6.md:128
; Annexe B - Lexique des résultats qui ont porté (fusion de la v6.0,
Claude). \textbf{Statut :} hérité du registre : statut non explicité
dans cette ligne historique.

\hypertarget{la-distance-de-suxe9curituxe9}{%
\subsection{La Distance de
sécurité}\label{la-distance-de-suxe9curituxe9}}

\textbf{Image héritée.} Les zéros ne doivent pas se coller : une
répulsion sous l'espacement moyen.

\textbf{Ancrage exact.} (H3⋆) (doc. H3) ; plus petit écart normalisé des
300 premiers zéros : 0,29.

\textbf{Garde héritée.} Non démontrée, hors de portée de RH seule.

\textbf{Provenance.} 1\_Presentation/20\_Registre\_metaphores\_V6.md:134
; Annexe B - Lexique des résultats qui ont porté (fusion de la v6.0,
Claude). \textbf{Statut :} hérité du registre : statut non explicité
dans cette ligne historique.

\hypertarget{le-point-critique}{%
\subsection{Le Point critique}\label{le-point-critique}}

\textbf{Image héritée.} ζ serait exactement au seuil d'une transition,
au bord de l'eau.

\textbf{Ancrage exact.} de Bruijn-Newman : 0 ≤ Λ ≤ 0,22 ; RH ⇔ Λ = 0
(Rodgers-Tao ; Polymath15).

\textbf{Garde héritée.} Λ n'est pas une température physique.

\textbf{Provenance.} 1\_Presentation/20\_Registre\_metaphores\_V6.md:135
; Annexe B - Lexique des résultats qui ont porté (fusion de la v6.0,
Claude). \textbf{Statut :} hérité du registre : statut non explicité
dans cette ligne historique.

\hypertarget{le-compteur}{%
\subsection{Le Compteur}\label{le-compteur}}

\textbf{Image héritée.} Au lieu de localiser chaque zéro (GPS), on
compte les valeurs propres positives d'une forme de Weil comprimée.

\textbf{Ancrage exact.} 2026 : ≥ 2/3 des zéros simples et sur la droite
(rang-trace, inertie de Sylvester) (doc. 13).

\textbf{Garde héritée.} Ne donne que des proportions.

\textbf{Provenance.} 1\_Presentation/20\_Registre\_metaphores\_V6.md:136
; Annexe B - Lexique des résultats qui ont porté (fusion de la v6.0,
Claude). \textbf{Statut :} hérité du registre : statut non explicité
dans cette ligne historique.

\hypertarget{le-musuxe9e-des-erreurs}{%
\subsection{Le Musée des erreurs}\label{le-musuxe9e-des-erreurs}}

\textbf{Image héritée.} Des erreurs réelles exposées, avec des leurres,
pour entraîner l'œil.

\textbf{Ancrage exact.} Banc d'audit : 21 erreurs, 7 leurres, script de
notation.

\textbf{Garde héritée.} Petit échantillon, un seul projet.

\textbf{Provenance.} 1\_Presentation/20\_Registre\_metaphores\_V6.md:142
; Annexe B - Lexique des résultats qui ont porté (fusion de la v6.0,
Claude). \textbf{Statut :} hérité du registre : statut non explicité
dans cette ligne historique.

\hypertarget{les-quatre-coins}{%
\subsection{Les Quatre coins}\label{les-quatre-coins}}

\textbf{Image héritée.} Une table à deux entrées dont les quatre cases
sont remplies.

\textbf{Ancrage exact.} ζ, L(s, χ₅) / E2 / F6 / DH : verrou local (fuite
nulle) et signe global (zéros sur la droite) sont indépendants.

\textbf{Garde héritée.} Les témoins sont fabriqués.

\textbf{Provenance.} 1\_Presentation/20\_Registre\_metaphores\_V6.md:152
; Mise à jour V6.3d/e - cinq noms proposés. \textbf{Statut :} hérité du
registre : statut non explicité dans cette ligne historique.

\hypertarget{le-peigne}{%
\subsection{Le Peigne}\label{le-peigne}}

\textbf{Image héritée.} Des dents régulières posées sur la droite
critique.

\textbf{Ancrage exact.} Zéros de 1 + 6\^{}(1/2-s) : pas 2π/log 6 ; leur
transformée de Fourier est la fuite Λ\_F(6\^{}k).

\textbf{Garde héritée.} Le peigne s'ajoute aux zéros de ζ, il ne les
déplace pas.

\textbf{Provenance.} 1\_Presentation/20\_Registre\_metaphores\_V6.md:153
; Mise à jour V6.3d/e - cinq noms proposés. \textbf{Statut :} hérité du
registre : statut non explicité dans cette ligne historique.

\hypertarget{lempreinte-en-creux}{%
\subsection{L'Empreinte en creux}\label{lempreinte-en-creux}}

\textbf{Image héritée.} Les zéros de ζ laissent leur moule dans les
écarts entre zéros des produits d'Euler.

\textbf{Ancrage exact.} Déficit des différences de zéros près des zéros
de ζ : 30 \% (ζ), 17 \% (L(s, χ₅)), 5 \% (DH).

\textbf{Garde héritée.} Phénomène connu (Pérez-Marco ; Ford-Zaharescu) ;
seule la comparaison avec le Sosie est propre au corpus.

\textbf{Provenance.} 1\_Presentation/20\_Registre\_metaphores\_V6.md:154
; Mise à jour V6.3d/e - cinq noms proposés. \textbf{Statut :} hérité du
registre : statut non explicité dans cette ligne historique.

\hypertarget{le-tour-du-propriuxe9taire}{%
\subsection{Le Tour du propriétaire}\label{le-tour-du-propriuxe9taire}}

\textbf{Image héritée.} Faire le tour de la parcelle et compter les
tours de l'aiguille.

\textbf{Ancrage exact.} Principe de l'argument : nombre exact de zéros
dans un rectangle ; a révélé deux zéros de DH manquants.

\textbf{Garde héritée.} Compte, ne localise pas.

\textbf{Provenance.} 1\_Presentation/20\_Registre\_metaphores\_V6.md:155
; Mise à jour V6.3d/e - cinq noms proposés. \textbf{Statut :} hérité du
registre : statut non explicité dans cette ligne historique.

\hypertarget{le-pied-uxe0-coulisse}{%
\subsection{Le Pied à coulisse}\label{le-pied-uxe0-coulisse}}

\textbf{Image héritée.} Les zéros hauts servent de règle pour mesurer
les zéros bas.

\textbf{Ancrage exact.} Demi-largeur du creux = 1 - β (heuristique) ;
pour ζ, 0,73 aux basses hauteurs, 0,51 ± 0,04 entre 1200 et 1500, pour
0,5 attendu (doc 24, 7.4).

\textbf{Garde héritée.} Reformulation de RH, probablement circulaire ;
ne mesure que des zéros déjà connus sur la droite.

\textbf{Provenance.} 1\_Presentation/20\_Registre\_metaphores\_V6.md:156
; Mise à jour V6.3d/e - cinq noms proposés. \textbf{Statut :} hérité du
registre : statut non explicité dans cette ligne historique.

\hypertarget{la-portuxe9e}{%
\subsection{La Portée}\label{la-portuxe9e}}

\textbf{Image héritée.} Jusqu'où la lunette porte avant de voir
l'intrus.

\textbf{Ancrage exact.} Largeur W du support (donc nombre de premiers) à
partir de laquelle une fenêtre voit un zéro hors droite : entre 1,5 et
2,0 pour le zéro de DH à 85,7.

\textbf{Garde héritée.} Propriété de l'instrument, pas de l'objet.

\textbf{Provenance.} 1\_Presentation/20\_Registre\_metaphores\_V6.md:170
; Mise à jour V6.4 - après le test de l'Atlas (document 26).
\textbf{Statut :} hérité du registre : statut non explicité dans cette
ligne historique.

\hypertarget{le-point-deuler}{%
\subsection{Le Point d'Euler}\label{le-point-deuler}}

\textbf{Image héritée.} Le seul endroit de la route où tout est aligné.

\textbf{Ancrage exact.} θ = 0 dans F\_θ : seul point où 100 \% des zéros
sont sur la droite et où l'empreinte de ζ est maximale.

\textbf{Garde héritée.} Une seule famille testée.

\textbf{Provenance.} 1\_Presentation/20\_Registre\_metaphores\_V6.md:180
; Mise à jour V6.5 - la famille modulo 7 (document 27). \textbf{Statut
:} hérité du registre : statut non explicité dans cette ligne
historique.

\hypertarget{la-plage-de-survie}{%
\subsection{La Plage de survie}\label{la-plage-de-survie}}

\textbf{Image héritée.} La bande de terrain sec qui rétrécit quand la
marée monte.

\textbf{Ancrage exact.} Intervalle de θ où tous les zéros jusqu'à T
restent sur la droite : {[}-0,0096 ; +0,072{]} pour T = 1000.

\textbf{Garde héritée.} Observée jusqu'à 1000 seulement.

\textbf{Provenance.} 1\_Presentation/20\_Registre\_metaphores\_V6.md:181
; Mise à jour V6.5 - la famille modulo 7 (document 27). \textbf{Statut
:} hérité du registre : statut non explicité dans cette ligne
historique.

\hypertarget{la-rencontre}{%
\subsection{La Rencontre}\label{la-rencontre}}

\textbf{Image héritée.} Deux danseurs se rejoignent sur la ligne et la
quittent ensemble.

\textbf{Ancrage exact.} Deux zéros voisins se touchent en θ\_c, puis
s'écartent de la droite comme √(θ - θ\_c).

\textbf{Garde héritée.} Rien n'est détruit : le nombre de zéros est
conservé.

\textbf{Provenance.} 1\_Presentation/20\_Registre\_metaphores\_V6.md:182
; Mise à jour V6.5 - la famille modulo 7 (document 27). \textbf{Statut
:} hérité du registre : statut non explicité dans cette ligne
historique.

\hypertarget{la-tour}{%
\subsection{La Tour}\label{la-tour}}

\textbf{Image héritée.} Les étages d'un immeuble, chacun contrôlé par la
même règle.

\textbf{Ancrage exact.} Les corps à qⁿ éléments ; à chaque étage n, la
forme de Castelnuovo Q\_n est positive.

\textbf{Garde héritée.} Pour ζ, on ne sait pas ce que sont les étages.

\textbf{Provenance.} 1\_Presentation/20\_Registre\_metaphores\_V6.md:192
; Mise à jour V6.6 - le Modèle réduit (document 28). \textbf{Statut :}
hérité du registre : statut non explicité dans cette ligne historique.

\hypertarget{la-positivituxe9-saturuxe9e}{%
\subsection{La Positivité saturée}\label{la-positivituxe9-saturuxe9e}}

\textbf{Image héritée.} Une poutre chargée exactement à sa limite, qui
tient.

\textbf{Ancrage exact.} Courbe y² = x⁵ + x³ + 1 sur F₃ (P(u) = 1 + 9u⁴)
: ρ\_4 = 1 exactement, 118 points sur F₈₁.

\textbf{Garde héritée.} Saturation à un étage, pas à tous.

\textbf{Provenance.} 1\_Presentation/20\_Registre\_metaphores\_V6.md:193
; Mise à jour V6.6 - le Modèle réduit (document 28). \textbf{Statut :}
hérité du registre : statut non explicité dans cette ligne historique.

\hypertarget{le-faux-jumeau-entier}{%
\subsection{Le Faux jumeau entier}\label{le-faux-jumeau-entier}}

\textbf{Image héritée.} Un faux billet avec un vrai filigrane.

\textbf{Ancrage exact.} Équation fonctionnelle, comptages de points
fermés entiers positifs, RH fausse (220 cas en genre 1 et 2 sur F₃, F₅)
; démasqué à l'étage 2 à 5.

\textbf{Garde héritée.} Ne vient d'aucune géométrie.

\textbf{Provenance.} 1\_Presentation/20\_Registre\_metaphores\_V6.md:194
; Mise à jour V6.6 - le Modèle réduit (document 28). \textbf{Statut :}
hérité du registre : statut non explicité dans cette ligne historique.

\hypertarget{le-rang}{%
\subsection{Le Rang}\label{le-rang}}

\textbf{Image héritée.} Plus d'inconnues que d'équations : il reste des
directions libres.

\textbf{Ancrage exact.} Fenêtre L \textgreater{} 2g : forme de Weil
d'une vraie courbe exactement singulière ; pour ζ, L et F6, chute
géométrique de λ\_min avec la dimension excédentaire.

\textbf{Garde héritée.} Pour ζ, la chute n'est jamais totale.

\textbf{Provenance.} 1\_Presentation/20\_Registre\_metaphores\_V6.md:204
; Mise à jour V6.7 - STABILITY-RADIUS (document 29). \textbf{Statut :}
hérité du registre : statut non explicité dans cette ligne historique.

\hypertarget{le-recensement}{%
\subsection{Le Recensement}\label{le-recensement}}

\textbf{Image héritée.} Un appel où chaque absent répond « absent ».

\textbf{Ancrage exact.} À la fenêtre 2g, nombre de valeurs propres
négatives = zéros hors droite / 2.

\textbf{Garde héritée.} Exact dans le modèle seulement.

\textbf{Provenance.} 1\_Presentation/20\_Registre\_metaphores\_V6.md:205
; Mise à jour V6.7 - STABILITY-RADIUS (document 29). \textbf{Statut :}
hérité du registre : statut non explicité dans cette ligne historique.

\hypertarget{la-rencontre-figuxe9e}{%
\subsection{La Rencontre figée}\label{la-rencontre-figuxe9e}}

\textbf{Image héritée.} Deux danseurs arrêtés au moment où ils se
rejoignent.

\textbf{Ancrage exact.} Vraie courbe à deux zéros confondus : marge de
l'objet nulle, RH démontrée.

\textbf{Garde héritée.} Aucune courbe voisine ne franchit.

\textbf{Provenance.} 1\_Presentation/20\_Registre\_metaphores\_V6.md:206
; Mise à jour V6.7 - STABILITY-RADIUS (document 29). \textbf{Statut :}
hérité du registre : statut non explicité dans cette ligne historique.

\hypertarget{le-passage-interdit}{%
\subsection{Le Passage interdit}\label{le-passage-interdit}}

\textbf{Image héritée.} Une porte au ras du mur, mais murée.

\textbf{Ancrage exact.} Marge nulle sans sortie : aucune perturbation
qui sort de RH n'est une courbe.

\textbf{Garde héritée.} Formulé dans le modèle seulement.

\textbf{Provenance.} 1\_Presentation/20\_Registre\_metaphores\_V6.md:207
; Mise à jour V6.7 - STABILITY-RADIUS (document 29). \textbf{Statut :}
hérité du registre : statut non explicité dans cette ligne historique.

\hypertarget{larchipel-deuler}{%
\subsection{L'Archipel d'Euler}\label{larchipel-deuler}}

\textbf{Image héritée.} Des îles dans une mer qui monte.

\textbf{Ancrage exact.} Dans V\_q, la région RH jusqu'à la hauteur T est
faite d'îles autour des points d'Euler, de rayon ≈ N\^{}(-0,77) ; à
l'infini, il ne reste que les points (Saias-Weingartner).

\textbf{Garde héritée.} La mer n'est pas uniforme : chenaux le long de
certains grands cercles.

\textbf{Provenance.} 1\_Presentation/20\_Registre\_metaphores\_V6.md:217
; Mise à jour V6.8 - L'Archipel d'Euler (document 30). \textbf{Statut :}
hérité du registre : statut non explicité dans cette ligne historique.

\hypertarget{luxeele-sans-mer}{%
\subsection{L'Île sans mer}\label{luxeele-sans-mer}}

\textbf{Image héritée.} Une île seule au monde.

\textbf{Ancrage exact.} Espaces de dimension 1 : ζ (Hamburger), L(s,
(·/3)), L(s, (·/5)) : aucune voisine qui garde l'équation fonctionnelle.

\textbf{Garde héritée.} Parmi les séries de Dirichlet seulement.

\textbf{Provenance.} 1\_Presentation/20\_Registre\_metaphores\_V6.md:218
; Mise à jour V6.8 - L'Archipel d'Euler (document 30). \textbf{Statut :}
hérité du registre : statut non explicité dans cette ligne historique.

\hypertarget{le-continent-bas}{%
\subsection{Le Continent bas}\label{le-continent-bas}}

\textbf{Image héritée.} À marée basse, les îles sont reliées.

\textbf{Ancrage exact.} À \textbar t\textbar{} ≤ 50, les trois points
d'Euler de V₇ impair sont dans une même région sans paire hors droite.

\textbf{Garde héritée.} Dépend de la carte et de la hauteur.

\textbf{Provenance.} 1\_Presentation/20\_Registre\_metaphores\_V6.md:219
; Mise à jour V6.8 - L'Archipel d'Euler (document 30). \textbf{Statut :}
hérité du registre : statut non explicité dans cette ligne historique.

\hypertarget{le-chauxeenon}{%
\subsection{Le Chaînon}\label{le-chauxeenon}}

\textbf{Image héritée.} Le pas discret est l'exponentielle d'un
générateur continu.

\textbf{Ancrage exact.} F = q\^{}D ; les zéros de ζ\_C sont le spectre
de D, répété avec la période 2πi / log q.

\textbf{Garde héritée.} Dans le Modèle réduit, D n'est défini qu'à cette
période près.

\textbf{Provenance.} 1\_Presentation/20\_Registre\_metaphores\_V6.md:229
; Mise à jour V6.9 - Le chaînon (document 32). \textbf{Statut :} hérité
du registre : statut non explicité dans cette ligne historique.

\hypertarget{le-spectre-des-longueurs}{%
\subsection{Le Spectre des longueurs}\label{le-spectre-des-longueurs}}

\textbf{Image héritée.} Les zéros chantent les longueurs des orbites.

\textbf{Ancrage exact.} S(u), calculé sur les zéros seuls, a des pics en
log n de hauteur c(n)/√n ; produit d'Euler ⟺ pics seulement en log
p\^{}k.

\textbf{Garde héritée.} Fenêtre finie ; fond lisse non retiré.

\textbf{Provenance.} 1\_Presentation/20\_Registre\_metaphores\_V6.md:230
; Mise à jour V6.9 - Le chaînon (document 32). \textbf{Statut :} hérité
du registre : statut non explicité dans cette ligne historique.

\hypertarget{laxe-par-zuxe9ro-distinct}{%
\subsection{L'Axe par zéro distinct}\label{laxe-par-zuxe9ro-distinct}}

\textbf{Image héritée.} Les nombres construisent un axe par note, pas
par instrument.

\textbf{Ancrage exact.} Le rang du Gram de Krylov est le nombre de zéros
distincts ; la multiplicité y est une masse.

\textbf{Garde héritée.} Pour ζ, la simplicité des zéros n'est que
conjecturée.

\textbf{Provenance.} 1\_Presentation/20\_Registre\_metaphores\_V6.md:231
; Mise à jour V6.9 - Le chaînon (document 32). \textbf{Statut :} hérité
du registre : statut non explicité dans cette ligne historique.

\hypertarget{la-muxe9trique-manquante}{%
\subsection{La Métrique manquante}\label{la-muxe9trique-manquante}}

\textbf{Image héritée.} La machine est là ; il manque la règle qui
mesure.

\textbf{Ancrage exact.} Les 220 faux jumeaux ont l'opérateur, la
symétrie et le réseau ; seul leur Gram est indéfini.

\textbf{Garde héritée.} Pour ζ, « là » veut dire dans une fenêtre finie.

\textbf{Provenance.} 1\_Presentation/20\_Registre\_metaphores\_V6.md:232
; Mise à jour V6.9 - Le chaînon (document 32). \textbf{Statut :} hérité
du registre : statut non explicité dans cette ligne historique.

\hypertarget{le-guxe9nuxe9rateur-des-premiers}{%
\subsection{Le Générateur des
premiers}\label{le-guxe9nuxe9rateur-des-premiers}}

\textbf{Image héritée.} Les nombres premiers jouent la partition ;
l'opérateur est l'instrument qu'on en déduit.

\textbf{Ancrage exact.} Avec les seuls premiers ≤ e\^{}(Nh+2a), le
faisceau de Toeplitz des traces donne e\^{}(iγh) pour les zéros de la
fenêtre (226 zéros de ζ avec les premiers ≤ 439).

\textbf{Garde héritée.} Fenêtre finie ; pas h choisi.

\textbf{Provenance.} 1\_Presentation/20\_Registre\_metaphores\_V6.md:242
; Mise à jour V6.10 - Le générateur des premiers (document 33).
\textbf{Statut :} hérité du registre : statut non explicité dans cette
ligne historique.

\hypertarget{la-masse-unituxe9}{%
\subsection{La Masse unité}\label{la-masse-unituxe9}}

\textbf{Image héritée.} Un vrai zéro pèse exactement son poids.

\textbf{Ancrage exact.} Une vraie valeur propre porte la masse ĝ(γ) ;
masse 1 = zéro simple ; test sans aucun zéro.

\textbf{Garde héritée.} Tolérance de 1 \% ; zéros proches confondus.

\textbf{Provenance.} 1\_Presentation/20\_Registre\_metaphores\_V6.md:243
; Mise à jour V6.10 - Le générateur des premiers (document 33).
\textbf{Statut :} hérité du registre : statut non explicité dans cette
ligne historique.

\hypertarget{le-repliement}{%
\subsection{Le Repliement}\label{le-repliement}}

\textbf{Image héritée.} Une roue filmée trop lentement semble tourner à
l'envers.

\textbf{Ancrage exact.} Au pas h, γ n'est vu qu'à 2π/h près ; au pas log
2, tous les zéros de ζ sont repliés.

\textbf{Garde héritée.} Le Modèle réduit, avec 2g zéros par période,
n'en souffre pas.

\textbf{Provenance.} 1\_Presentation/20\_Registre\_metaphores\_V6.md:244
; Mise à jour V6.10 - Le générateur des premiers (document 33).
\textbf{Statut :} hérité du registre : statut non explicité dans cette
ligne historique.

\hypertarget{la-muxe9trique-en-avance}{%
\subsection{La Métrique en avance}\label{la-muxe9trique-en-avance}}

\textbf{Image héritée.} La sentinelle entend avant de voir.

\textbf{Ancrage exact.} DH : le quadruplet de 240,4 crée deux directions
négatives avant d'être résolu par l'opérateur.

\textbf{Garde héritée.} Seuil de bruit (10\^{}-11 de la plus grande
valeur propre).

\textbf{Provenance.} 1\_Presentation/20\_Registre\_metaphores\_V6.md:245
; Mise à jour V6.10 - Le générateur des premiers (document 33).
\textbf{Statut :} hérité du registre : statut non explicité dans cette
ligne historique.

\hypertarget{le-pas-de-base}{%
\subsection{Le Pas de base}\label{le-pas-de-base}}

\textbf{Image héritée.} La base d'un corps fini est la cadence d'une
caméra qui filme un même mouvement.

\textbf{Ancrage exact.} F\_q² voit le même D avec une demi-période :
zéros = courbe ∪ tordue (525 sur 525) ; RH conservée.

\textbf{Garde héritée.} Pour ζ, il n'y a pas de cadence naturelle.

\textbf{Provenance.} 1\_Presentation/20\_Registre\_metaphores\_V6.md:255
; Mise à jour V6.11 - La base comme pas de temps (document 35).
\textbf{Statut :} hérité du registre : statut non explicité dans cette
ligne historique.

\hypertarget{la-descente}{%
\subsection{La Descente}\label{la-descente}}

\textbf{Image héritée.} Redescendre d'un étage, c'est deviner un signe
par pièce.

\textbf{Ancrage exact.} Les objets sur F\_q compatibles avec les
comptages sur F\_q² sont les choix de signe des facteurs irréductibles
(525 sur 525 ; 2 dans 440 cas, jusqu'à 12).

\textbf{Garde héritée.} Polynômes de Weil, pas forcément des courbes.

\textbf{Provenance.} 1\_Presentation/20\_Registre\_metaphores\_V6.md:256
; Mise à jour V6.11 - La base comme pas de temps (document 35).
\textbf{Statut :} hérité du registre : statut non explicité dans cette
ligne historique.

\hypertarget{la-muxe9trique-noyuxe9e}{%
\subsection{La Métrique noyée}\label{la-muxe9trique-noyuxe9e}}

\textbf{Image héritée.} Une voix seule couverte par une foule.

\textbf{Ancrage exact.} À nombres premiers fixés, DH perd ses directions
négatives quand le pas grandit (10 → 0 entre h = 0,01 et 0,1) ; à N
fixé, elle les garde.

\textbf{Garde héritée.} Une seule fonction test ; seuil de bruit fixé.

\textbf{Provenance.} 1\_Presentation/20\_Registre\_metaphores\_V6.md:257
; Mise à jour V6.11 - La base comme pas de temps (document 35).
\textbf{Statut :} hérité du registre : statut non explicité dans cette
ligne historique.

\hypertarget{linnovation}{%
\subsection{L'Innovation}\label{linnovation}}

\textbf{Image héritée.} Le GPS qui dit, à chaque pas, si la position est
encore possible.

\textbf{Ancrage exact.} Récurrence de Levinson sur les traces des
premiers : ζ et L restent dans le disque (\textbar k\_n\textbar{} ≤
0,79) ; DH sonne à la fenêtre 3,52 (premiers ≤ 36).

\textbf{Garde héritée.} Horizon numérique (double précision).

\textbf{Provenance.} 1\_Presentation/20\_Registre\_metaphores\_V6.md:267
; Mise à jour V6.12 - Les idées rejetées relues (document 36).
\textbf{Statut :} hérité du registre : statut non explicité dans cette
ligne historique.

\hypertarget{le-fantuxf4me-pesuxe9}{%
\subsection{Le Fantôme pesé}\label{le-fantuxf4me-pesuxe9}}

\textbf{Image héritée.} On reconnaît un fantôme à ce qu'il ne pèse rien.

\textbf{Ancrage exact.} Le mode fantôme de V4.3 (Im ≈ 73,8) porte une
masse de 1,3·10\^{}-9 ; les vrais zéros, 1.

\textbf{Garde héritée.} Critère par fenêtre, pas borne uniforme.

\textbf{Provenance.} 1\_Presentation/20\_Registre\_metaphores\_V6.md:268
; Mise à jour V6.12 - Les idées rejetées relues (document 36).
\textbf{Statut :} hérité du registre : statut non explicité dans cette
ligne historique.

\hypertarget{le-rayon-spectral}{%
\subsection{Le Rayon spectral}\label{le-rayon-spectral}}

\textbf{Image héritée.} La plus grande dérive se lit au bord du spectre.

\textbf{Ancrage exact.} Θ - ½ dans une fenêtre = max \textbar Im
γ\textbar{} du générateur : 0,308517 pour DH, \textless{} 10\^{}-6 pour
ζ et L.

\textbf{Garde héritée.} Borne inférieure : un quadruplet non résolu ne
compte pas.

\textbf{Provenance.} 1\_Presentation/20\_Registre\_metaphores\_V6.md:269
; Mise à jour V6.12 - Les idées rejetées relues (document 36).
\textbf{Statut :} hérité du registre : statut non explicité dans cette
ligne historique.

\hypertarget{le-temps-de-heisenberg}{%
\subsection{Le Temps de Heisenberg}\label{le-temps-de-heisenberg}}

\textbf{Image héritée.} Une note n'est entendue que lorsque l'oreille
sépare les notes voisines.

\textbf{Ancrage exact.} Un quadruplet hors droite à la hauteur γ₀ est
détecté à la fenêtre L\_c ≈ log(q·γ₀/2π) : rapport 0,95-1,08 pour ζ et
L(s, χ₅), 0,80-0,89 pour DH ; premiers jusqu'à ≈ q·γ₀/2π (conducteur
analytique).

\textbf{Garde héritée.} Poids \textless{} 10\^{}-5 (lobes secondaires de
ĝ) ; repliement si h \textgreater{} π/γ₀ ; le facteur de DH n'est pas
expliqué.

\textbf{Provenance.} 1\_Presentation/20\_Registre\_metaphores\_V6.md:279
; Mise à jour V6.14 - La loi de seuil (document 38). \textbf{Statut :}
hérité du registre : statut non explicité dans cette ligne historique.

\hypertarget{la-balance-dinertie}{%
\subsection{La Balance d'inertie}\label{la-balance-dinertie}}

\textbf{Image héritée.} Deux poids opposés sur un plateau qui est le
fond.

\textbf{Ancrage exact.} G = B + paire cesse d'être positive ⟺ C\^{}-1 +
U*B\^{}-1U est définie positive (noyau reproduisant du fond aux deux
points de la paire).

\textbf{Garde héritée.} B doit être définie positive.

\textbf{Provenance.} 1\_Presentation/20\_Registre\_metaphores\_V6.md:280
; Mise à jour V6.14 - La loi de seuil (document 38). \textbf{Statut :}
hérité du registre : statut non explicité dans cette ligne historique.

\hypertarget{le-signe-fragile}{%
\subsection{Le Signe fragile}\label{le-signe-fragile}}

\textbf{Image héritée.} L'aiguille d'une boussole cassée continue de
pointer avec assurance.

\textbf{Ancrage exact.} Fond positif par construction, innovation de
-2·10\^{}-3 à la taille 296 (λ\_min/λ\_max = 4·10\^{}-17) ; à 30
chiffres, une erreur de 0,45 sur la fenêtre.

\textbf{Garde héritée.} -

\textbf{Provenance.} 1\_Presentation/20\_Registre\_metaphores\_V6.md:281
; Mise à jour V6.14 - La loi de seuil (document 38). \textbf{Statut :}
hérité du registre : statut non explicité dans cette ligne historique.

\hypertarget{le-poids-archimuxe9dien}{%
\subsection{Le Poids archimédien}\label{le-poids-archimuxe9dien}}

\textbf{Image héritée.} La longueur d'une corde fixe la note qu'elle
peut entendre.

\textbf{Ancrage exact.} La longueur de détection d'un quadruplet est
Ω(γ₀) = 2π × densité des zéros : 1,02-1,08 pour ζ, L(s, χ₅) et L(Δ, s) ;
elle double en degré 2.

\textbf{Garde héritée.} Facteur 0,8-0,9 pour DH et 0,87-0,98 pour ζ·L.

\textbf{Provenance.} 1\_Presentation/20\_Registre\_metaphores\_V6.md:291
; Mise à jour V6.15 - Le degré 2 et le recensement (document 39).
\textbf{Statut :} hérité du registre : statut non explicité dans cette
ligne historique.

\hypertarget{le-quadruplet-repliuxe9}{%
\subsection{Le Quadruplet replié}\label{le-quadruplet-repliuxe9}}

\textbf{Image héritée.} Un passager qui revient par l'autre porte.

\textbf{Ancrage exact.} Au pas h, un quadruplet au-delà de π/h entre
dans la fenêtre en γ - 2π/h ; 4 quadruplets de DH (1 092 à 1 464) sont
comptés ainsi, avec un poids de 5·10\^{}-8 à 5·10\^{}-7.

\textbf{Garde héritée.} Le poids tombe vite dans les lobes secondaires.

\textbf{Provenance.} 1\_Presentation/20\_Registre\_metaphores\_V6.md:292
; Mise à jour V6.15 - Le degré 2 et le recensement (document 39).
\textbf{Statut :} hérité du registre : statut non explicité dans cette
ligne historique.

\hypertarget{la-somme-sans-ruxe9pulsion}{%
\subsection{La Somme sans répulsion}\label{la-somme-sans-ruxe9pulsion}}

\textbf{Image héritée.} Deux chorales qui chantent sans s'écouter.

\textbf{Ancrage exact.} ζ·L(s, χ₅) : allumage à 0,87-0,98 de Ω, contre
1,02-1,08 pour une seule fonction.

\textbf{Garde héritée.} Un seul produit testé, non pré-enregistré.

\textbf{Provenance.} 1\_Presentation/20\_Registre\_metaphores\_V6.md:293
; Mise à jour V6.15 - Le degré 2 et le recensement (document 39).
\textbf{Statut :} hérité du registre : statut non explicité dans cette
ligne historique.

\hypertarget{la-rencontre-uxe0-trois}{%
\subsection{La Rencontre à trois}\label{la-rencontre-uxe0-trois}}

\textbf{Image héritée.} Trois chemins qui se croisent au même carrefour
: le choc de deux zéros.

\textbf{Ancrage exact.} À la fenêtre de Heisenberg, la marge locale de
la métrique d'une paire de zéros d'écart normalisé s vaut 1 - sinc²(πs),
la densité de corrélation de paires de Montgomery ; elle s'annule à la
rencontre. Mesuré sur 1 070 zéros de ζ : 4 marges sous 0,1 (13 pour le
GUE, 320 pour Poisson).

\textbf{Garde héritée.} Modèle à deux atomes ; les voisins comptent pour
la métrique réelle.

\textbf{Provenance.} 1\_Presentation/20\_Registre\_metaphores\_V6.md:303
; Mise à jour V6.16 - La Rencontre à trois (document 40). \textbf{Statut
:} hérité du registre : statut non explicité dans cette ligne
historique.

\hypertarget{le-compteur-de-paires}{%
\subsection{Le Compteur de paires}\label{le-compteur-de-paires}}

\textbf{Image héritée.} Un compteur qui ne bouge qu'aux chocs.

\textbf{Ancrage exact.} β = 2Σ\textbar u\_γ\textbar² est lisse à 95 \%
(terme θ′) ; sa part d'écart 2Σ(Z″/Z′)²/\textbar ρ\textbar² est portée à
35 \% par les 10 \% de zéros les plus serrés ; la paire de Lehmer près
de 7 005 y pèse environ mille fois la médiane.

\textbf{Garde héritée.} Heuristique pour la finitude (répulsion
nécessaire).

\textbf{Provenance.} 1\_Presentation/20\_Registre\_metaphores\_V6.md:304
; Mise à jour V6.16 - La Rencontre à trois (document 40). \textbf{Statut
:} hérité du registre : statut non explicité dans cette ligne
historique.

\hypertarget{la-zone-dombre}{%
\subsection{La Zone d'ombre}\label{la-zone-dombre}}

\textbf{Image héritée.} Une case blanche sur la carte.

\textbf{Ancrage exact.} Case de la carte des croisements sans croisement
produit par le corpus, ou croisement sans pendant dans le Modèle réduit
; le navigateur peut la viser.

\textbf{Garde héritée.} Une case vide n'est pas une question : il faut
encore la formuler.

\textbf{Provenance.} 1\_Presentation/20\_Registre\_metaphores\_V6.md:305
; Mise à jour V6.16 - La Rencontre à trois (document 40). \textbf{Statut
:} hérité du registre : statut non explicité dans cette ligne
historique.

\hypertarget{le-contour-tournuxe9}{%
\subsection{Le Contour tourné}\label{le-contour-tournuxe9}}

\textbf{Image héritée.} Regarder une falaise de biais pour en voir le
relief.

\textbf{Ancrage exact.} Λ(6+it) = 2e\^{}\{-φt\} Re(e\^{}\{6iφ\}
A(6+it)), A = ∫ F(e\^{}\{x+iφ\}) e\^{}\{(6+it)x\} dx, F(w) = Σ τ(n)
e\^{}\{-2πnw\} : la rotation du contour de Mellin, repliée par la
modularité, supprime l'annulation exponentielle ; 2 139 zéros de L(Δ, s)
jusqu'à 1 500 en double précision (accord 10⁻¹³ avec mpmath, complétude
de Turing).

\textbf{Garde héritée.} Il faut la modularité pour replier le contour.

\textbf{Provenance.} 1\_Presentation/20\_Registre\_metaphores\_V6.md:313
; Mise à jour V6.17 - La Rencontre en degré 2 (document 41).
\textbf{Statut :} hérité du registre : statut non explicité dans cette
ligne historique.

\hypertarget{la-loi-de-la-marge}{%
\subsection{La Loi de la marge}\label{la-loi-de-la-marge}}

\textbf{Image héritée.} Une foule qui laisse à chacun une place en
proportion de l'espace disponible.

\textbf{Ancrage exact.} La fréquence d'une configuration est
proportionnelle à la marge de la métrique : à la fenêtre N, marge totale
= densité de Weyl de U(N), marge d'une paire = corrélation de paires du
CUE(N) ; paire conjuguée : marge = sin²θ = (π/2) × densité de Sato-Tate
(démontré) ; zéros de ζ et de Δ (mesuré, déficit à basse hauteur).

\textbf{Garde héritée.} Familles symplectiques (genre ≥ 2) : la densité
de Haar n'est plus le déterminant de la métrique de taille 2g.

\textbf{Provenance.} 1\_Presentation/20\_Registre\_metaphores\_V6.md:314
; Mise à jour V6.17 - La Rencontre en degré 2 (document 41).
\textbf{Statut :} hérité du registre : statut non explicité dans cette
ligne historique.

\hypertarget{les-familles-uxe9tranguxe8res}{%
\subsection{Les Familles
étrangères}\label{les-familles-uxe9tranguxe8res}}

\textbf{Image héritée.} Deux foules qui se traversent sans se voir.

\textbf{Ancrage exact.} Une fonction non primitive superpose des
familles sans répulsion mutuelle : ζ·L(χ5), 15 \% de marges sous 0,1,
toutes mixtes, paire mixte à 8,9·10⁻⁴ ; pendant exact : le changement de
base F\_p → F\_p², P(u²) = P(u)·P(-u), où les collisions existent
(courbes supersingulières) avec RH vraie.

\textbf{Garde héritée.} Sur Q aucune collision n'est connue ; sur F\_p²
elles existent.

\textbf{Provenance.} 1\_Presentation/20\_Registre\_metaphores\_V6.md:315
; Mise à jour V6.17 - La Rencontre en degré 2 (document 41).
\textbf{Statut :} hérité du registre : statut non explicité dans cette
ligne historique.

\hypertarget{la-rencontre-uxe0-n}{%
\subsection{La Rencontre à n}\label{la-rencontre-uxe0-n}}

\textbf{Image héritée.} La Rencontre à trois, vue en relief : une
surface au lieu d'une courbe.

\textbf{Ancrage exact.} Pour n atomes à la fenêtre N, det(G/N) =
det{[}D\_N(θ\_l - θ\_m){]} = (2π/N)\textsuperscript{n} R\_n du CUE(N)
(démontré) ; les triplets des 1 070 zéros de ζ suivent le noyau sinus R3
(χ² réduit 1,34, Poisson 17,3) ; pour ζ·L(χ5), loi de superposition de
deux familles indépendantes (χ² 0,87).

\textbf{Garde héritée.} Mesure à basse hauteur : 0 triplet serré pour
4,4 attendus (effet de taille finie supposé).

\textbf{Provenance.} 1\_Presentation/20\_Registre\_metaphores\_V6.md:323
; Mise à jour V6.19 - La Rencontre à n et la matrice des théories
(documents 43 et 44). \textbf{Statut :} hérité du registre : statut non
explicité dans cette ligne historique.

\hypertarget{le-carillon-de-bohr}{%
\subsection{Le Carillon de Bohr}\label{le-carillon-de-bohr}}

\textbf{Image héritée.} Un carillon dont chaque cloche est un zéro : le
signal sonne sans s'amplifier ni s'éteindre.

\textbf{Ancrage exact.} Ψ(u) =
e\textsuperscript{-u/2}C(e\textsuperscript{u})/u dans l'espace de
Besicovitch B² ; coefficients de Bohr A\_γ(L) ; sous RH,
Σ\textbar A\_γ\textbar² = S (Théorème B).

\textbf{Garde héritée.} Proposé, pas encore calculé (axe E du plan).

\textbf{Provenance.} 1\_Presentation/20\_Registre\_metaphores\_V6.md:324
; Mise à jour V6.19 - La Rencontre à n et la matrice des théories
(documents 43 et 44). \textbf{Statut :} hérité du registre : statut non
explicité dans cette ligne historique.

\hypertarget{le-palier-pruxe9cisuxe9}{%
\subsection{Le Palier (précisé)}\label{le-palier-pruxe9cisuxe9}}

\textbf{Image héritée.} Un palier par endroits, avec des marches cassées
là où le sol se dérobe.

\textbf{Ancrage exact.} Médiane de e\_N² log N/κ = 0,985 × 4/3 de 10⁴ à
5·10⁶, sans montée visible ; compensation b² ≈ 1,16, r\_Q ≈ 0,72, queue
≈ 0,15 ; pics jusqu'à 3,27 × 4/3 là où A\_N s'annule ; sous RH, la
médiane tend vers (S/4κ)·4/3 si la queue s'éteint.

\textbf{Garde héritée.} La limite 4/3 n'est pas atteinte dans les
calculs ; elle reste conditionnelle.

\textbf{Provenance.} 1\_Presentation/20\_Registre\_metaphores\_V6.md:332
; Mise à jour V6.20 - Le Palier recalculé (document 45). \textbf{Statut
:} hérité du registre : statut non explicité dans cette ligne
historique.

\hypertarget{la-surface-virtuelle}{%
\subsection{La Surface virtuelle}\label{la-surface-virtuelle}}

\textbf{Image héritée.} Une surface dont on connaît les croisements sans
en connaître les points.

\textbf{Ancrage exact.} Forme d'intersection sur \{fibres, graphes de
Frobenius\} : dans le Modèle réduit, celle de C × C (indice de Hodge :
425 courbes sur 425 ; surplus exact pour les 220 faux jumeaux) ; pour ζ,
la même par le dictionnaire fibres ↔ pôle, intersections ↔ puissances de
premiers, genre ↔ terme archimédien ; une seule direction positive
depuis les premiers ≤ 439 ; DH : 2 puis 10 en trop. Orthogonalement aux
fibres, c'est l'opposé de la métrique des zéros.

\textbf{Garde héritée.} Le dictionnaire est classique (Weil, Bombieri,
Connes) ; la surface elle-même n'existe pas.

\textbf{Provenance.} 1\_Presentation/20\_Registre\_metaphores\_V6.md:338
; Mise à jour V6.21 - La surface virtuelle (document 46). \textbf{Statut
:} hérité du registre : statut non explicité dans cette ligne
historique.

\hypertarget{la-fonction-auxiliaire}{%
\subsection{La Fonction auxiliaire}\label{la-fonction-auxiliaire}}

\textbf{Image héritée.} Un filet tendu sur tous les points à la fois :
s'il y en a trop, il se déchire.

\textbf{Ancrage exact.} Stepanov-Bombieri : par Riemann-Roch et le
Frobenius, une fonction s'annule fortement en tous les points
rationnels, d'où N\_m ≤ Q + (2g + 1)√Q + 1 ; avec le tordu quadratique,
elle démasque les 220 faux jumeaux (étage médian 4, Castelnuovo 3) ;
Clifford seul n'en démasque aucun.

\textbf{Garde héritée.} Pour ζ, la fonction auxiliaire n'existe pas
encore : il manque le Frobenius.

\textbf{Provenance.} 1\_Presentation/20\_Registre\_metaphores\_V6.md:344
; Mise à jour V6.22 - Le Riemann-Roch virtuel (document 47).
\textbf{Statut :} hérité du registre : statut non explicité dans cette
ligne historique.

\hypertarget{le-frobenius-absolu}{%
\subsection{Le Frobenius absolu}\label{le-frobenius-absolu}}

\textbf{Image héritée.} Une horloge qui tournerait sur les entiers comme
x ↦ x\^{}q tourne sur une courbe.

\textbf{Ancrage exact.} Hypothèse de travail de l'axe J3 : une
application de « Spec Z » dans lui-même dont les points fixes seraient
les nombres premiers ; Connes-Consani proposent l'action d'échelle.

\textbf{Garde héritée.} Rien de tel n'est construit ; l'image n'est pas
un objet.

\textbf{Provenance.} 1\_Presentation/20\_Registre\_metaphores\_V6.md:345
; Mise à jour V6.22 - Le Riemann-Roch virtuel (document 47).
\textbf{Statut :} hérité du registre : statut non explicité dans cette
ligne historique.

\hypertarget{la-duxe9rivuxe9e-nulle}{%
\subsection{La Dérivée nulle}\label{la-duxe9rivuxe9e-nulle}}

\textbf{Image héritée.} Une aiguille qui pointe exactement au même
endroit que les autres, mais qui ne bouge pas quand on la touche.

\textbf{Ancrage exact.} Le Frobenius x ↦ x\^{}p, l'identité et une
permutation prennent les mêmes valeurs ou permutent les points ; seule
la dérivée nulle de x\^{}p fait chuter le rang des conditions de
Stepanov (34 au lieu de 104 pour p = 101) : 13 à 52 fonctions
auxiliaires contre 0.

\textbf{Garde héritée.} Aux tailles calculées, la borne de Stepanov
reste triviale ; seule la construction est vérifiée.

\textbf{Provenance.} 1\_Presentation/20\_Registre\_metaphores\_V6.md:351
; Mise à jour V6.23 - Le Frobenius qui n'est pas un morphisme (document
48). \textbf{Statut :} hérité du registre : statut non explicité dans
cette ligne historique.

\hypertarget{le-frobenius-trop-bon}{%
\subsection{Le Frobenius trop bon}\label{le-frobenius-trop-bon}}

\textbf{Image héritée.} Une horloge si parfaite qu'elle ne perd jamais
une seconde, et ne peut donc rien marquer.

\textbf{Ancrage exact.} En caractéristique un (Connes-Consani), Fr\_λ
est un automorphisme pour tout λ \textgreater{} 0 : il dilate les pentes
au lieu d'annuler les dérivations ; il ne fournit pas l'élément
invisible aux dérivations de Stepanov. La contrainte doit venir de la
positivité sur le carré (route de Weil).

\textbf{Garde héritée.} Lu sur résumés et extraits.

\textbf{Provenance.} 1\_Presentation/20\_Registre\_metaphores\_V6.md:357
; Mise à jour V6.26 - J3b (document 51). \textbf{Statut :} hérité du
registre : statut non explicité dans cette ligne historique.

\hypertarget{index-technique-et-version-interactive}{%
\section{Index technique et version
interactive}\label{index-technique-et-version-interactive}}

Le lecteur interactif \texttt{69\_LANGUE\_DES\_PASSAGES.html} permet
recherche par nom, symbole, image ou garde, filtres de niveau et
consultation d'une fiche à la fois. L'atelier applique les règles typées
et signale les moments perdus.

L'index \texttt{v644\_index\_fonctions.csv} / \texttt{.json}, dans le
dossier de données, donne fichier, ligne, nom qualifié, signature, type
et docstring. Sans docstring, le comportement n'est pas inventé. La
couverture technique est celle de l'AST, non une analyse sémantique
certifiée de tous les programmes.

Le dictionnaire JSON permet d'ajouter des entrées qualifiées en
conservant leur provenance. Les symboles G, P, D, U, R, gamma et kappa
demandent leur contexte : le même caractère écrit ne suffit pas à
définir le même objet.

Les unités fictives accompagnent les calculs : x logarithmique sans
dimension, f d'amplitude A, Q et métriques d'unité A², quotient spectral
sans dimension. Le mot énergie, chaleur ou entropie ne fournit pas une
preuve de signe.

Les preuves complètes et les expériences figurent dans le document 68 ;
les critères globaux, branches, collisions et coefficients des anciennes
familles gardent leurs obligations ouvertes. Cette langue crée des
phrases vérifiables et des ponts de domaines ; elle ne fournit pas
encore la loi arithmétique globale.

\end{document}
